%%%%%%%%%%%%%%%%%%%%%%%%%%%%%%%%%%%%%%%%%%%%%%%%%%%%%%%%%%%%%%%%%%%%%
%% This is a "template" model document for submission to the
%% American Chemical Society (ACS).
%%
%% The guidance here contains information about how you may wish to
%% modify it to match the requirements of various ACS journals. The
%% ACS do *not* typeset accepted articles using LaTeX, so there is
%% no specific class required.
%%
%% This template deliberately does *not* seek to reproduce
%% the layout of the typeset journal: this is explicitly not
%% required by the ACS for LaTeX submissions.
%%
%% Please report any issues with the template at
%% https://github.com/josephwright/acs-template/issues
%%
%% Released under the Creative Commons 0 license
%% https://creativecommons.org/public-domain/cc0/
%% 
%% Copyight (c) 2025 Joseph Wright
%%%%%%%%%%%%%%%%%%%%%%%%%%%%%%%%%%%%%%%%%%%%%%%%%%%%%%%%%%%%%%%%%%%%%
\makeatletter\def\input@path{
{Figures/}
}\makeatother%
\documentclass[journal=jcisd8,manuscript=suppinfo,
keywords=true]{achemso}

% Page numbering for SI: S1, S2, S3, etc. (ACS requirement)
\pagenumbering{arabic}
\renewcommand{\thepage}{S\arabic{page}}

%%%%%%%%%%%%%%%%%%%%%%%%%%%%%%%%%%%%%%%%%%%%%%%%%%%%%%%%%%%%%%%%%%%%%
%% Font setup - delete if you are using LuaLaTeX
%%%%%%%%%%%%%%%%%%%%%%%%%%%%%%%%%%%%%%%%%%%%%%%%%%%%%%%%%%%%%%%%%%%%%
\usepackage[T1]{fontenc}

%%%%%%%%%%%%%%%%%%%%%%%%%%%%%%%%%%%%%%%%%%%%%%%%%%%%%%%%%%%%%%%%%%%%%
%% Adjust the margins and allow for line spacing
%%%%%%%%%%%%%%%%%%%%%%%%%%%%%%%%%%%%%%%%%%%%%%%%%%%%%%%%%%%%%%%%%%%%%
\usepackage{geometry}
\geometry{margin = 1in}
\usepackage{setspace}

%%%%%%%%%%%%%%%%%%%%%%%%%%%%%%%%%%%%%%%%%%%%%%%%%%%%%%%%%%%%%%%%%%%%%
%% Reference support
%%
%% The recommended method for producing the reference section is
%% to use biblatex. If you wish to use a classical BibTeX
%% approach, this is easiest to achieve using the achemso package.
%% In that case, you should remove the biblatex lines.
%%%%%%%%%%%%%%%%%%%%%%%%%%%%%%%%%%%%%%%%%%%%%%%%%%%%%%%%%%%%%%%%%%%%%
% You can adjust the printing of DOI, article title, etc. using
% package options, e.g. "doi = true"; see the biblatex manual for
% details of adjusting the number of authors printed, e.g.
% "maxnames = 15" to print no more than 15 authors.
%\usepackage[style = chem-acs]{biblatex}
%\addbibresource{TADF_Article_References.bib}
% If you are using classical BibTeX, remove the above lines and 
% uncomment:
%\usepackage{achemso}

%%%%%%%%%%%%%%%%%%%%%%%%%%%%%%%%%%%%%%%%%%%%%%%%%%%%%%%%%%%%%%%%%%%%%
%% Graphic inclusion and scheme and chart support
%%%%%%%%%%%%%%%%%%%%%%%%%%%%%%%%%%%%%%%%%%%%%%%%%%%%%%%%%%%%%%%%%%%%%
\usepackage{graphicx}
\usepackage{float}
\newfloat{scheme}{htbp}{los}
\floatname{scheme}{Scheme}
\floatname{chart}{Chart}
\newfloat{graph}{htbp}{loh}

%%%%%%%%%%%%%%%%%%%%%%%%%%%%%%%%%%%%%%%%%%%%%%%%%%%%%%%%%%%%%%%%%%%%%
%% Common support packages
%%%%%%%%%%%%%%%%%%%%%%%%%%%%%%%%%%%%%%%%%%%%%%%%%%%%%%%%%%%%%%%%%%%%%
\usepackage{chemformula,chemfig}
\setchemfig{atom style={scale=0.4}, double bond sep = 4pt}
\usepackage[version=4,arrows=pgf-filled,textfontname=sffamily,
mathfontname=mathsf]{mhchem}
% Typesetting and Layout Utilities
\usepackage{microtype}  % Improves typography (spacing, kerning)
\usepackage[detect-all=true,group-minimum-digits=4]{siunitx}
\sisetup{mode = math, propagate-math-font = true, reset-math-version = false, range-phrase = --}
% siunitx setup
%\usepackage{lscape}     % For landscape pages, if needed
%\usepackage[labelformat=simple]{subcaption}
\usepackage{tabularx, longtable}%
\usepackage{booktabs, multirow,xcolor,cleveref,url}

\usepackage{seqsplit}   % For splitting long sequences in text

%%%%%%%%%%%%%%%%%%%%%%%%%%%%%%%%%%%%%%%%%%%%%%%%%%%%%%%%%%%%%%%%%%%%%
%% Many journals require that sections are unnumbered: this 
%% is activated here
%%%%%%%%%%%%%%%%%%%%%%%%%%%%%%%%%%%%%%%%%%%%%%%%%%%%%%%%%%%%%%%%%%%%%
\setcounter{secnumdepth}{-1}

%%%%%%%%%%%%%%%%%%%%%%%%%%%%%%%%%%%%%%%%%%%%%%%%%%%%%%%%%%%%%%%%%%%%%
%% Place any additional macros here.  Please use \newcommand* where
%% possible, and avoid layout-changing macros (which are not used
%% when typesetting).
%%%%%%%%%%%%%%%%%%%%%%%%%%%%%%%%%%%%%%%%%%%%%%%%%%%%%%%%%%%%%%%%%%%%%

\newcommand{\smiles}[1]{\parbox{5.5cm}{\seqsplit{\ttfamily#1}}}

\newcommand{\deltaest}{$\Delta E_{\text{ST}}$}
\newcommand{\risc}{rISC}
\newcommand{\tict}{TICT}
\DeclareSIUnit\calorie{cal}
\DeclareSIUnit\kcal{\kilo\calorie}
\DeclareSIUnit\kcal{Kcal}
\DeclareSIUnit\atomicunit{a.u.}

%%%%%%%%%%%%%%%%

%%%%%%%%%%%%%%%%%%%%%%%%%%%%%%%%%%%%%%%%%%%%%%%%%%%%%%%%%%%%%%%%%%%%%
%% Author and title data:
%% the authblk package is currently the simplest way to provide this
%%%%%%%%%%%%%%%%%%%%%%%%%%%%%%%%%%%%%%%%%%%%%%%%%%%%%%%%%%%%%%%%%%%%%
%\usepackage{authblk}
\author{Jean-Pierre Tchapet Njafa}
\email{jean-pierre.tchapet@facsciences-uy1.cm}
\affiliation[University of Yaoundé 1]{Department of Physics, Faculty of Science,
University of Yaounde I, P.O.Box
812, Yaounde, Cameroon}

\author{Elvira Vanelle Kameni Tcheuffa}
\affiliation[University of Yaoundé 1]{Department of Physics, Faculty of Science,
University of Yaounde I, P.O.Box
812, Yaounde, Cameroon}

\author{Aissatou Maghame Foumkpou}
\affiliation[University of Yaoundé 1]{Department of Physics, Faculty of Science,
University of Yaounde I, P.O.Box
812, Yaounde, Cameroon}

\author{Serge Guy Nana Engo}
\affiliation[University of Yaoundé 1]{Department of Physics, Faculty of Science,
University of Yaounde I, P.O.Box
812, Yaounde, Cameroon}

\title{xTB-Based High-Throughput Screening of TADF Emitters: 747-Molecule Benchmark}

\begin{document}

%%%%%%%%%%%%%%%%%%%%%%%%%%%%%%%%%%%%%%%%%%%%%%%%%%%%%%%%%%%%%%%%%%%%%
%% Title and Author Information (ACS SI Requirement - Point 3) voir au dessus
%%%%%%%%%%%%%%%%%%%%%%%%%%%%%%%%%%%%%%%%%%%%%%%%%%%%%%%%%%%%%%%%%%%%%

\maketitle

%%%%%%%%%%%%%%%%%%%%%%%%%%%%%%%%%%%%%%%%%%%%%%%%%%%%%%%%%%%%%%%%%%%%%

\section{Stokes Shifts and Estimation of Relaxed Singlet-Triplet Excitation
Energies}\label{sec:SI-stokes}
Direct geometry optimization of the $S_1$ state is not straightforward with the current
xTB~ implementation used. Therefore, an estimation for the relaxed
singlet-triplet gap, $\Delta E_\text{ST,relaxed}$, was made.
The geometric energy relaxation for the $T_1$ state is $\Delta E_\text{relax}(T_1) =
\Delta E_v(S_0\to T_1)_{\text{xTB}} - \Delta E_r(S_0\to T_1)_{\text{xTB}}$.
The Stokes shift for fluorescence ($S_1 \rightarrow S_0$) is approximated by assuming a
similar geometric relaxation energy for $S_1$ as for $T_1$, i.e.,
$S_{shift} \approx \Delta E_\text{relax}(T_1)$. This is a strong approximation.
Then, the relaxed $S_1$ energy is $E_r(S_1) \approx E_v(S_1)_{\text{sTDA/sTD-DFT}} -
S_{shift}$.
And $\Delta E_\text{ST,relaxed} \approx E_r(S_1) - E_r(T_1)_\text{xTB}$.

\newpage
\begin{figure}[!ht]
\centering
\includegraphics[width=\textwidth]{singlet_triplet_with_outliers.pdf}
\caption{Correlation between predicted and reference singlet-triplet energy gaps
(\deltaest)
for TADF molecules. One physically
unrealistic outlier (Spiro-CN, reference \deltaest~ $= \qty{5.50}{\electronvolt}$, shown
as red X) was
excluded from regression analysis.
Statistics show Theil-Sen robust regression with 95\% bootstrap confidence intervals for
slope. Despite moderate scatter, the
correlations indicate that semi-empirical methods capture qualitative trends in
\deltaest.}
\label{fig:st-corrSI}
\end{figure}
\newpage

\section{Dataset}

{\renewcommand{\arraystretch}{2.25}
\begin{longtable}[H]{*{2}{l@{}w{c}{5.5cm}}}
\caption{Molecular architecture of the studied TADF emitters. This table provides the SMILES representation of each molecule.} 
\label{tab:supp-mol-list} \\
\toprule 
Molecule Id & SMILES & Molecule Id & SMILES \\
\midrule
\endfirsthead 
\multicolumn{4}{c}{{\tablename} \thetable{} -- Continued from previous page} \\ 
\toprule 
Molecule Id & SMILES & Molecule Id & SMILES \\
\midrule 
\endhead 
\midrule 
\multicolumn{4}{r}{{Continued on next page}} \\ 
\endfoot 
\bottomrule 
\endlastfoot 
4CzTPN	&	\smiles{N\#Cc1c(-n2c3ccccc3c3ccccc32)c(-n2c3ccccc3c3ccccc32)c(C\#N)c(-n2c3ccccc3c3ccccc32)c1-n1c2ccccc2c2ccccc21}	&	SPBP-DPAC	&	\smiles{O=C1c2ccccc2C2(c3cc(N4c5ccccc5C(c5ccccc5)(c5ccccc5)c5ccccc54)ccc31)c1ccccc1-c1ccccc12}	\\
DMAC-DPS	&	\smiles{CC1(C)c2ccccc2N(c2ccc(S(=O)(=O)c3ccc(N4c5ccccc5C(C)(C)c5ccccc54)cc3)cc2)c2ccccc21}	&	PX-TRZ	&	\smiles{Cc1cc(C)c(B(c2ccc(-c3nc(Cl)nc(-c4ccccc4)n3)cc2)c2c(C)cc(C)cc2C)c(C)c1}	\\
PSPCz	&	\smiles{C1=CC=C(C=C1)S(=O)(=O)C2=CC=CC(=C2)N3C4=CC=CC=C4C5=CC=CC=C53}	&	C2MPI	&	\smiles{Cc1ccc(N(c2ccc(C)cc2)c2ccccc2C2(c3ccccc3N(c3ccc(C)cc3)c3ccc(C)cc3)CCCCC2)cc1}	\\
Px2BP	&	\smiles{C1=CC=C2C(=C1)N(C3=CC=CC=C3O2)C4=CC=C(C=C4)C(=O)C5=CC=C(C=C5)N6C7=CC=CC=C7OC8=CC=CC=C86}	&	p-PCzTPD	&	\smiles{CCCCCCCCN1C(=O)c2c(-c3ccc(-n4c5ccccc5c5ccccc54)cc3)sc(-c3ccc(-n4c5ccccc5c5ccccc54)cc3)c2C1=O}	\\
CzS2	&	\smiles{C1=CC=C2C(=C1)C3=CC=CC=C3N2C4=CC=C(C=C4)S(=O)(=O)C5=CC=C(C=C5)N6C7=CC=CC=C7C8=CC=CC=C86}	&	m-PCzTPD	&	\smiles{CCCCCCCCN1C(=O)c2c(-c3cccc(-n4c5ccccc5c5ccccc54)c3)sc(-c3cccc(-n4c5ccccc5c5ccccc54)c3)c2C1=O}	\\
2TCz-DPS	&	\smiles{CC(C)(C)C1=CC2=C(C=C1)N(C3=C2C=C(C=C3)C(C)(C)C)C4=CC=C(C=C4)S(=O)(=O)C5=CC=C(C=C5)N6C7=C(C=C(C=C7)C(C)(C)C)C8=C6C=CC(=C8)C(C)(C)C}	&	o-PCzTPD	&	\smiles{CCCCCCCCN1C(=O)c2c(-c3ccccc3-n3c4ccccc4c4ccccc43)sc(-c3ccccc3-n3c4ccccc4c4ccccc43)c2C1=O}	\\
TDBA-DI	&	\smiles{CC(C)(C)c1ccc2c(c1)B1c3cc(C(C)(C)C)ccc3Oc3cc(-n4c5ccccc5c5c6c(c7ccccc7n6-c6ccccc6)c6c(c7ccccc7n6-c6ccccc6)c54)cc(c31)O2}	&	CPBzFC	&	\smiles{c1ccc2c(c1)c1ccccc1n2-c1ccc(-c2ccc(-c3ccc4c(c3)Cc3cc(-n5c6ccccc6c6ccccc65)ccc3-4)c3nnsc23)cc1}	\\
OPDPO	&	\smiles{O=CC1=CCC(c2ccccc2N2c3ccccc3Sc3ccccc32)(P(=O)(c2ccccc2)c2ccccc2)C=C1}	&	SiCz2Py2	&	\smiles{c1cncc(-c2ccc([Si](c3ccc(-c4cccnc4)cc3)(c3ccc(-n4c5ccccc5c5ccccc54)cc3)c3ccc(-n4c5ccccc5c5ccccc54)cc3)cc2)c1}	\\
2CzTPEPCz	&	\smiles{C(=C(c1ccc(-n2c3ccccc3c3ccccc32)cc1)c1ccc(-n2c3ccccc3c3ccccc32)cc1)c1ccc(-c2ccc(-n3c4ccccc4c4ccccc43)cc2)cc1}	&	SiCz1Py3	&	\smiles{c1cncc(-c2ccc([Si](c3ccc(-c4cccnc4)cc3)(c3ccc(-c4cccnc4)cc3)c3ccc(-n4c5ccccc5c5ccccc54)cc3)cc2)c1}	\\
TBRb	&	\smiles{CC(C)(C)c1ccc(-c2c3ccc(C(C)(C)C)cc3c(-c3ccccc3)c3c(-c4ccc(C(C)(C)C)cc4)c4ccc(C(C)(C)C)cc4c(-c4ccccc4)c23)cc1}	&	SiCz3Py1	&	\smiles{c1cncc(-c2ccc([Si](c3ccc(-n4c5ccccc5c5ccccc54)cc3)(c3ccc(-n4c5ccccc5c5ccccc54)cc3)c3ccc(-n4c5ccccc5c5ccccc54)cc3)cc2)c1}	\\
TBPe	&	\smiles{CC(C)(C)c1cc2cc(C(C)(C)C)cc3c4cc(C(C)(C)C)cc5cc(C(C)(C)C)cc(c(c1)c23)c54}	&	POBBPE	&	\smiles{CC(C)(C)c1ccc(Oc2ccc(C(C)(C)C)cc2P(=O)(c2ccccc2)c2ccccc2)cc1}	\\
2CzPN	&	\smiles{N\#Cc1cc(-c2cccc3c2[nH]c2ccccc23)c(-c2cccc3c2[nH]c2ccccc23)cc1C\#N}	&	CzppT	&	\smiles{c1ccc(-c2c(-c3ccccc3)c(-c3ccc(-c4cccc5c4[nH]c4ccccc45)cc3)c(-c3ccccn3)c(-c3ccccn3)c2-c2ccc(-c3cccc4c3[nH]c3ccccc34)cc2)cc1}	\\
3CzFTFP	&	\smiles{Fc1nc(F)c(F)c(C2(c3ccc4c(c3)c3ccccc3n4-c3ccccc3)c3ccccc3-c3ccccc32)c1F}	&	SBDBQ-PXZ	&	\smiles{c1ccc(-c2nc3ccc(N4c5ccccc5Oc5ccccc54)cc3nc2-c2ccccc2)cc1}	\\
DDPhCz-DCPP	&	\smiles{N\#Cc1nc2c3ccc(-n4c5ccc(-c6ccccc6)cc5c5cc(-c6ccccc6)ccc54)cc3c3cc(-n4c5ccc(-c6ccccc6)cc5c5cc(-c6ccccc6)ccc54)ccc3c2nc1C\#N}	&	DBQ-3PXZ	&	\smiles{c1ccc2c(c1)Oc1ccccc1N2c1ccc(-c2nc3ccc(N4c5ccccc5Oc5ccccc54)cc3nc2-c2ccc(N3c4ccccc4Oc4ccccc43)cc2)cc1}	\\
ACRXTN	&	\smiles{CC1(C)c2ccccc2N(c2ccc3c(=O)c4ccccc4oc3c2)c2ccccc21}	&	P4CzCN–BCz	&	\smiles{C=Cc1ccc(COc2ccc(C(C)(C)c3ccc(O)cc3)cc2)cc1}	\\
TXO-TPA	&	\smiles{O=C1c2ccccc2S(=O)(=O)c2ccc(-c3ccc(N(c4ccccc4)c4ccccc4)cc3)cc21}	&	NAP-1,5-DPA	&	\smiles{COc1ccc(N(c2ccc(OC)cc2)c2cccc3c(N(c4ccc(OC)cc4)c4ccc(OC)cc4)cccc23)cc1}	\\
PCTXO	&	\smiles{O=C1c2ccccc2S(=O)(=O)c2ccc(-c3ccc(N(c4ccc(Br)cc4)c4ccc(Br)cc4)cc3)cc21}	&	3DPA3CN	&	\smiles{N\#Cc1c(-c2ccc(N(c3ccccc3)c3ccccc3)cc2)c(C\#N)c(-c2ccc(N(c3ccccc3)c3ccccc3)cc2)c(C\#N)c1-c1ccc(N(c2ccccc2)c2ccccc2)cc1}	\\
IP-PPI	&	\smiles{c1ccc(-n2c(-c3ccc(-c4ccc(-c5cn6ccccc6n5)cc4)cc3)nc3c4ccccc4c4ccccc4c32)cc1}	&	TPA-AN-NA	&	\smiles{c1ccc(N(c2ccccc2)c2ccc(-c3c4ccccc4c(-c4ccc5ccccc5c4)c4ccccc34)cc2)cc1}	\\
DMDHNP-DPS	&	\smiles{CC1(C)c2cnccc2Nc2ccncc21.O=S(=O)(c1ccccc1)c1ccccc1}	&	OCzBN	&	\smiles{COc1ccc2c(c1)c1ccccc1n2-c1cc(-n2c3ccccc3c3cc(OC)ccc32)c(C\#N)c(-n2c3ccccc3c3cc(OC)ccc32)c1}	\\
4NTAZ-PPI	&	\smiles{c1ccc(-c2nnc(-c3ccc(-c4ccc(-c5nc6c7ccccc7c7ccccc7c6n5-c5ccccc5)cc4)cc3)n2-c2cccc3ccccc23)cc1}	&	BBCz	&	\smiles{c1ccc2c(c1)[nH]c1cc3[nH]c4ccccc4c3cc12}	\\
TAZ-PPI	&	\smiles{c1ccc(-c2nnc(-c3ccc(-c4ccc(-c5nc6c7ccccc7c7ccccc7c6n5-c5ccccc5)cc4)cc3)n2-c2ccccc2)cc1}	&	IQ-Se	&	\smiles{COc1c(C\#N)c(C)cc2c([Se][Se]c3c(C)nc(N4CCCCC4)c4c(OC)c(C\#N)c(C)cc34)c(C)nc(N3CCCCC3)c12}	\\
PTZ-mPYR	&	\smiles{Cc1cc(N2c3ccccc3Sc3ccccc32)ccc1-c1cc(-c2ccc(N3c4ccccc4Sc4ccccc43)cc2C)ncn1}	&	dTolmp	&	\smiles{Cc1ccccc1P(CP(CP(c1ccccc1C)c1ccccc1C)c1ccccc1)c1ccccc1C}	\\
pCNBCzoCF3	&	\smiles{N\#Cc1ccc(-n2c3ccccc3c3cc(-c4ccc5c(c4)c4ccccc4n5-c4ccc(C\#N)cc4C(F)(F)F)ccc32)c(C(F)(F)F)c1}	&	SiTCNCz	&	\smiles{N\#Cc1cccc(-c2nc(-c3cccc([Si](c4ccccc4)(c4ccccc4)c4ccccc4)c3)nc(-n3c4ccccc4c4ccccc43)n2)c1}	\\
PTZTPE-2	&	\smiles{CC1(C)OB(c2ccc(C(=C(c3ccccc3)c3ccccc3)c3ccccc3)cc2)OC1(C)C}	&	dpmb	&	\smiles{Cc1cc(P(c2ccccc2)c2ccccc2)c(P(c2ccccc2)c2ccccc2)cc1C}	\\
4CzIPN	&	\smiles{N\#Cc1c(-n2c3ccccc3c3ccccc32)c(C\#N)c(-n2c3ccccc3c3ccccc32)c(-n2c3ccccc3c3ccccc32)c1-n1c2ccccc2c2ccccc21}	&	3CzBN	&	\smiles{N\#Cc1c(-n2c3ccccc3c3ccccc32)cc(-n2c3ccccc3c3ccccc32)cc1-n1c2ccccc2c2ccccc21}	\\
DMAC-TRZ	&	\smiles{CC1(C)c2ccccc2N(c2ccc(-c3nc(-c4ccccc4)nc(-c4ccccc4)n3)cc2)c2ccccc21}	&	TCB	&	\smiles{O=C(c1ccccc1)c1ccc(-n2c3ccccc3c3cc(-c4ccc(N(c5ccccc5)c5ccccc5)cc4)ccc32)cc1}	\\
TCzTrz	&	\smiles{c1ccc(-c2nc(-c3ccccc3)nc(-c3cc(-n4c5ccccc5c5ccccc54)c(-n4c5ccccc5c5ccccc54)c(-n4c5ccccc5c5ccccc54)c3)n2)cc1}	&	CPzPyC	&	\smiles{c1ccc2c(c1)c1ccccc1n2-c1cncc(-c2cncc(-n3c4ccccc4c4ccccc43)n2)c1}	\\
$\delta$-2CbPN	&	\smiles{N\#Cc1cc(-n2c3ccccc3c3ncccc32)c(-n2c3ccccc3c3ncccc32)cc1C\#N}	&	CPzPzC	&	\smiles{c1ccc2c(c1)c1ccccc1n2-c1cncc(-c2cncc(-n3c4ccccc4c4ccccc43)n2)n1}	\\
CzDCbTrz	&	\smiles{c1ccc(-c2nc(-c3ccccc3)nc(-c3cc(-n4c5ccccc5c5ncccc54)c(-n4c5ccccc5c5ccccc54)c(-n4c5ccccc5c5ncccc54)c3)n2)cc1}	&	CPzPC	&	\smiles{c1cc(-c2cncc(-n3c4ccccc4c4ccccc43)n2)cc(-n2c3ccccc3c3ccccc32)c1}	\\
$\alpha$-2CbPN	&	\smiles{N\#Cc1cc(-n2c3ccccc3c3cccnc32)c(-n2c3ccccc3c3cccnc32)cc1C\#N}	&	CBIQD	&	\smiles{O=C1c2cccc3cccc(c23)C(=O)N1c1nc2ccc(Cl)cc2s1}	\\
CN-QP	&	\smiles{N\#Cc1ccc2nc(-c3ccc(N4c5ccccc5Oc5ccccc54)cc3)c(-c3ccc(N4c5ccccc5Oc5ccccc54)cc3)nc2c1}	&	PXZ-QL	&	\smiles{c1ccc2c(c1)Oc1ccccc1N2c1ccc2ccccc2n1}	\\
TFM-QP	&	\smiles{FC(F)(F)c1ccc2nc(-c3ccc(N4c5ccccc5Oc5ccccc54)cc3)c(-c3ccc(N4c5ccccc5Oc5ccccc54)cc3)nc2c1}	&	SBA	&	\smiles{c1ccc2c(c1)Nc1ccccc1C21c2ccccc2Nc2ccccc21}	\\
NDN	&	\smiles{c1ccc2c(c1)c1cccnc1n2-c1ccc2sc3ccc(-n4c5ccccc5c5cccnc54)cc3c2c1}	&	TATP-BP	&	\smiles{O=C(c1ccc(-n2c3ccccc3c3c4c(c5ccccc5n4-c4ccccc4)c4c(c5ccccc5n4-c4ccccc4)c32)cc1)c1ccc(-n2c3ccccc3c3c4c(c5ccccc5n4-c4ccccc4)c4c(c5ccccc5n4-c4ccccc4)c32)cc1}	\\
mCP	&	\smiles{c1cc(-n2c3ccccc3c3ccccc32)cc(-n2c3ccccc3c3ccccc32)c1}	&	Pd-1	&	\smiles{CCCCN(CCCC)c1ccc(C=Nc2cc(C\#N)c(C\#N)cc2N=Cc2ccc(N(CCCC)CCCC)cc2O)c(O)c1}	\\
ZDZ	&	\smiles{c1ccc2c(c1)c1ccccc1n2-c1ccc2sc3ccc(-n4c5ccccc5c5ccccc54)cc3c2c1}	&	C3F7	&	\smiles{Cc1cc(C)c(B(c2ccccc2-n2c3ccccc3c3ccccc32)c2c(C)cc(C)cc2C)c(C)c1}	\\
ZDN	&	\smiles{c1ccc2c(c1)c1ccccc1n2-c1ccc2sc3ccc(-n4c5ccccc5c5cccnc54)cc3c2c1}	&	BuCzCF3oB	&	\smiles{Cc1cc(C)c(B(c2ccc(C(F)(F)F)cc2-n2c3ccc(C(C)(C)C)cc3c3cc(C(C)(C)C)ccc32)c2c(C)cc(C)cc2C)c(C)c1}	\\
BTCz3CzTPN	&	\smiles{N\#Cc1c(-n2c3ccccc3c3ccccc32)c(-n2c3ccccc3c3ccccc32)c(C\#N)c(-n2c3ccccc3c3c4sc5ccccc5c4ccc32)c1-n1c2ccccc2c2ccccc21}	&	phenazasiline	&	\smiles{Brc1ccc2nc3c4ccc(Br)cc4c4ccccc4c3nc2c1}	\\
2BTCz2CzTPN	&	\smiles{N\#Cc1c(-n2c3ccccc3c3ccccc32)c(-n2c3ccccc3c3c4sc5ccccc5c4ccc32)c(C\#N)c(-n2c3ccccc3c3ccccc32)c1-n1c2ccccc2c2c3sc4ccccc4c3ccc21}	&	D–A–D	&	\smiles{Brc1ccc2nc3c4ccc(Br)cc4c4ccccc4c3nc2c1}	\\
APDC-DTPA	&	\smiles{N\#Cc1nc2c(nc1C\#N)-c1ccc(-c3ccc(N(c4ccccc4)c4ccccc4)cc3)c3c(-c4ccc(N(c5ccccc5)c5ccccc5)cc4)ccc-2c13}	&	pzpy	&	\smiles{c1ccc(-c2nc3ccccc3n2-c2ccc(-c3c4ccccc4c(-c4ccc(-n5c(-c6ccccc6)nc6ccccc65)cc4)c4ccccc34)cc2)cc1}	\\
BMIM	&	\smiles{N\#Cc1nc2c(nc1C\#N)-c1ccc(-c3ccc(N(c4ccccc4)c4ccccc4)cc3)c3c(-c4ccc(N(c5ccccc5)c5ccccc5)cc4)ccc-2c13}	&	m-PO-ABN	&	\smiles{N\#Cc1ccc(-c2c3ccccc3c(-c3cccc(-c4nc5c6ccccc6c6ccccc6c5o4)c3)c3ccccc23)cc1}	\\
ADO-TPA	&	\smiles{O=C1C(=O)c2ccc(-c3ccc(N(c4ccccc4)c4ccccc4)cc3)c3cccc1c23}	&	PO-ABN	&	\smiles{N\#Cc1ccc(-c2c3ccccc3c(-c3cccc(-c4nc5c6ccccc6c6ccccc6c5o4)c3)c3ccccc23)cc1}	\\
DCzTrz	&	\smiles{c1ccc(-c2nc(-c3ccccc3)nc(-c3cc(-n4c5ccccc5c5ccccc54)cc(-n4c5ccccc5c5ccccc54)c3)n2)cc1}	&	p-PO-ABN	&	\smiles{N\#Cc1ccc(-c2c3ccccc3c(-c3ccc(-c4nc5c6ccccc6c6ccccc6c5o4)cc3)c3ccccc23)cc1}	\\
BDPCC	&	\smiles{c1ccc(-c2ccc3[nH]c4c(-c5ccc6c([nH]c7ccccc76)c5-c5cc(-c6ccccc6)cc6c5[nH]c5ccc(-c7ccccc7)cc56)cc(-c5ccccc5)cc4c3c2)cc1}	&	teracridan	&	\smiles{c1ccc2c(c1)Cc1c(cccc1-c1ccc3c(c1-c1cccc4c1Cc1ccccc1N4)Cc1ccccc1N3)N2}	\\
SpiroAC-TRZ	&	\smiles{c1ccc(-c2nc(-c3ccccc3)nc(-c3ccc(N4c5ccccc5C5(c6ccccc6-c6ccccc65)c5ccccc54)cc3)n2)cc1}	&	PTPC	&	\smiles{CC(C)(C)c1ccc2c(c1)c1cc(C(C)(C)C)ccc1n2-c1ccc(-c2ccc(-c3cc(-c4ccccc4)c(-c4ccc(-c5nc6c7ccccc7c7ccccc7c6n5-c5ccccc5)cc4)cc3-c3ccccc3)cc2)cc1}	\\
DBT-BZ-DMAC	&	\smiles{CC1(C)c2ccccc2Nc2cccc(C(=O)c3ccccc3)c21.c1ccc2c(c1)sc1ccccc12}	&	TPTI	&	\smiles{O=c1[nH]c2ccsc2c2cc3c(=O)[nH]c4ccsc4c3cc12}	\\
DQP	&	\smiles{c1ccc2nc3c(nc2c1)c1nc2ccccc2nc1c1nc2ccccc2nc31}	&	BDBTCz-Trz	&	\smiles{c1ccc(-c2nc(-c3ccccc3)nc(-c3ccc(-n4c5ccc(-c6ccc7sc8ccccc8c7c6)cc5c5cc(-c6ccc7sc8ccccc8c7c6)ccc54)cc3)n2)cc1}	\\
TPA-DCPP	&	\smiles{N\#Cc1nc2c3ccc(-c4ccc(N(c5ccccc5)c5ccccc5)cc4)cc3c3cc(-c4ccc(N(c5ccccc5)c5ccccc5)cc4)ccc3c2nc1C\#N}	&	BDBFCz-Trz	&	\smiles{c1ccc(-c2nc(-c3ccccc3)nc(-c3ccc(-n4c5ccc(-c6ccc7oc8ccccc8c7c6)cc5c5cc(-c6ccc7oc8ccccc8c7c6)ccc54)cc3)n2)cc1}	\\
2PXZ-OXD	&	\smiles{c1ccc2c(c1)Oc1ccccc1N2c1ccc(-c2nnc(-c3ccc(N4c5ccccc5Oc5ccccc54)cc3)o2)cc1}	&	HAP-3MF	&	\smiles{Cc1cc(C2=NC3=NC(c4ccc(F)c(C)c4)=NC4=NC(c5ccc(F)c(C)c5)=NC(=N2)N34)ccc1F}	\\
BmPyPhB	&	\smiles{CC1(C)c2ccccc2N(c2ccc(-c3nc(-c4ccc(N5c6ccccc6C(C)(C)c6ccccc65)cc4)nc(-c4ccc(N5c6ccccc6C(C)(C)c6ccccc65)cc4)n3)cc2)c2ccccc21}	&	FlCz	&	\smiles{CC1(C)c2ccccc2-c2ccc(-c3ncc(-c4ccc(-n5c6ccccc6c6ccccc65)cc4)cn3)cc21}	\\
MPPA	&	\smiles{COc1ccc(N(c2ccc(OC)cc2)c2ccc(-c3ccc4c(c3)C(=O)c3ccc(-c5ccc(N(c6ccc(OC)cc6)c6ccc(OC)cc6)cc5)cc3C4=O)cc2)cc1}	&	SFCz	&	\smiles{c1ccc2c(c1)-c1ccccc1C21c2ccccc2-c2ccc(-c3ncc(-c4ccc(-n5c6ccccc6c6ccccc65)cc4)cn3)cc21}	\\
ACRFLCN	&	\smiles{N\#Cc1ccc2c(c1)C1(c3cc(C\#N)ccc3-2)c2ccccc2N(c2ccccc2)c2ccccc21}	&	o-TPA-AP-TPA	&	\smiles{c1ccc(N(c2ccccc2)c2ccccc2-c2c3ccccc3c(-c3ccc4ccc5c(-c6ccccc6N(c6ccccc6)c6ccccc6)ccc6ccc3c4c65)c3ccccc23)cc1}	\\
PTZ-DCPP	&	\smiles{N\#Cc1nc2c3ccc(N4c5ccccc5Sc5ccccc54)cc3c3cc(N4c5ccccc5Sc5ccccc54)ccc3c2nc1C\#N}	&	p-TPA-AP-TPA	&	\smiles{c1ccc(N(c2ccccc2)c2ccc(-c3c4ccccc4c(-c4ccc5ccc6c(-c7ccc(N(c8ccccc8)c8ccccc8)cc7)ccc7ccc4c5c76)c4ccccc34)cc2)cc1}	\\
BBA-Pyr	&	\smiles{c1ccc(-c2ccc(N(c3ccc(-c4ccccc4)cc3)c3ccc(-c4ccc(-c5cc(-c6ccccc6)nc(-c6ccccc6)n5)cc4)cc3)cc2)cc1}	&	TNPZ	&	\smiles{N\#Cc1ccc(-c2ccc3c(c2)c2cc(-c4ccc(C\#N)cc4)ccc2c2nc4cc(-c5ccc(N(c6ccccc6)c6ccccc6)cc5)c(-c5ccc(N(c6ccccc6)c6ccccc6)cc5)cc4nc32)cc1}	\\
PhICz-Pyr	&	\smiles{c1ccc(-c2ccc3c(c2)c2cc(-c4ccccc4)cc4c5cc(-c6ccc(-c7cc(-c8ccccc8)nc(-c8ccccc8)n7)cc6)ccc5n3c24)cc1}	&	m-TPA-AP-TPA	&	\smiles{c1ccc(N(c2ccccc2)c2cccc(-c3c4ccccc4c(-c4ccc5ccc6c(-c7cccc(N(c8ccccc8)c8ccccc8)c7)ccc7ccc4c5c76)c4ccccc34)c2)cc1}	\\
PhCz-Pyr	&	\smiles{c1ccc(-c2ccc3c(c2)c2cc(-c4ccccc4)ccc2n3-c2ccc(-c3ccc(-c4cc(-c5ccccc5)nc(-c5ccccc5)n4)cc3)cc2)cc1}	&	5,6PXZ-PIDO	&	\smiles{CC1(C)C(=O)c2cc(-c3ccc(N4c5ccccc5Oc5ccccc54)cc3)c(-c3ccc(N4c5ccccc5Oc5ccccc54)cc3)cc2C1=O}	\\
PDA-DP	&	\smiles{CC1(C)c2cc(-c3ccccc3)ccc2N(c2ccc(S(=O)(=O)c3ccc(N4c5ccc(-c6ccccc6)cc5C(C)(C)c5cc(-c6ccccc6)ccc54)cc3)cc2)c2ccc(-c3ccccc3)cc21}	&	TRZ-3SO2	&	\smiles{O=S(=O)(c1ccccc1)c1cccc(-c2nc(-c3cccc(S(=O)(=O)c4ccccc4)c3)nc(-c3cccc(S(=O)(=O)c4ccccc4)c3)n2)c1}	\\
9CzFDBFDPO	&	\smiles{O=P(c1ccccc1)(c1ccccc1)c1cccc2c1oc1c(P(=O)(c3ccccc3)c3ccccc3)cc(-c3ccc(C4(c5ccc(-n6c7ccccc7c7ccccc76)cc5)c5ccccc5-c5ccccc54)cc3)cc12}	&	35tCzBPym	&	\smiles{CC(C)(C)c1ccc2c(c1)c1cc(C(C)(C)C)ccc1n2-c1cc(C(=O)c2cncnc2)cc(-n2c3ccc(C(C)(C)C)cc3c3cc(C(C)(C)C)ccc32)c1}	\\
9PhCzFDBFDPO	&	\smiles{O=P(c1ccccc1)(c1ccccc1)c1cccc2c1oc1c(P(=O)(c3ccccc3)c3ccccc3)cc(-c3ccc(C4(c5ccc(-c6ccc(-n7c8ccccc8c8ccccc87)cc6)cc5)c5ccccc5-c5ccccc54)cc3)cc12}	&	35CzBPym	&	\smiles{O=C(c1cncnc1)c1cc(-n2c3ccccc3c3ccccc32)cc(-n2c3ccccc3c3ccccc32)c1}	\\
9PhCzFDBFSPO	&	\smiles{O=P(c1ccccc1)(c1ccccc1)c1cccc2c1oc1ccc(-c3ccc(C4(c5ccc(-c6ccc(-n7c8ccccc8c8ccccc87)cc6)cc5)c5ccccc5-c5ccccc54)cc3)cc12}	&	25CzBPym	&	\smiles{O=C(c1cncnc1)c1cc(-n2c3ccccc3c3ccccc32)ccc1-n1c2ccccc2c2ccccc21}	\\
9CzFDBFSPO	&	\smiles{O=P(c1ccccc1)(c1ccccc1)c1cccc2c1oc1ccc(-c3ccc(C4(c5ccc(-n6c7ccccc7c7ccccc76)cc5)c5ccccc5-c5ccccc54)cc3)cc12}	&	25tCzBPym	&	\smiles{CC(C)(C)c1ccc2c(c1)c1cc(C(C)(C)C)ccc1n2-c1ccc(-n2c3ccc(C(C)(C)C)cc3c3cc(C(C)(C)C)ccc32)c(C(=O)c2cncnc2)c1}	\\
9ArFDBFxPO	&	\smiles{O=P(c1ccccc1)(c1ccccc1)c1cccc2c1oc1c(P(=O)(c3ccccc3)c3ccccc3)cc(-c3ccc(C4(c5ccc(-n6c7ccccc7c7ccccc76)cc5)c5ccccc5-c5ccccc54)cc3)cc12}	&	BPBiPA	&	\smiles{c1ccc(-c2nc3ccccc3n2-c2ccc(-c3c4ccccc4c(-c4ccc(-n5c(-c6ccccc6)nc6ccccc65)cc4)c4ccccc34)cc2)cc1}	\\
12BTCzTPN	&	\smiles{N\#Cc1c(-n2c3ccccc3c3ccccc32)c(-n2c3ccccc3c3ccc4c5ccccc5sc4c32)c(C\#N)c(-n2c3ccccc3c3ccccc32)c1-n1c2ccccc2c2ccc3c4ccccc4sc3c21}	&	BTZ–DMAC	&	\smiles{CC1(C)c2ccccc2N(c2ccc(N3c4ccccc4C(C)(C)c4ccccc43)c3nsnc23)c2ccccc21}	\\
oDTBPZ-DPXZ	&	\smiles{Cc1ccccc1-c1cc2nc3c4ccc(N5c6ccccc6Oc6ccccc65)cc4c4cc(N5c6ccccc6Oc6ccccc65)ccc4c3nc2cc1-c1ccccc1C}	&	BTZ–DPA	&	\smiles{c1ccc(N(c2ccccc2)c2ccc(N(c3ccccc3)c3ccccc3)c3nsnc23)cc1}	\\
3DMAC-BP-CN	&	\smiles{CC1(C)c2ccccc2N(c2ccc3c(c2)c2cc(N4c5ccccc5C(C)(C)c5ccccc54)ccc2c2nc4cc(N5c6ccccc6C(C)(C)c6ccccc65)c(C\#N)cc4nc32)c2ccccc21}	&	2F	&	\smiles{c1ccc(N(c2ccccc2)c2ccc3c(c2)c2cc(N(c4ccccc4)c4ccccc4)ccc2n3-c2ccc(-c3ccc4nccnc4c3)cc2)cc1}	\\
3DMAC-BP-Br	&	\smiles{CC1(C)c2ccccc2N(c2ccc3c(c2)c2cc(N4c5ccccc5C(C)(C)c5ccccc54)ccc2c2nc4cc(N5c6ccccc6C(C)(C)c6ccccc65)c(Br)cc4nc32)c2ccccc21}	&	PFBP-2a	&	\smiles{CC(C)(C)c1ccc2c(c1)C(C)(C)c1cc(C(C)(C)C)ccc1N2c1c(F)c(F)c(-c2c(F)c(F)c(F)c(F)c2F)c(F)c1F}	\\
TPA-PZCN	&	\smiles{N\#Cc1ccc2c(c1)c1cc(C\#N)ccc1c1nc3cc(-c4ccc(N(c5ccccc5)c5ccccc5)cc4)c(-c4ccc(N(c5ccccc5)c5ccccc5)cc4)cc3nc21}	&	PFBP-2b	&	\smiles{CC(C)(C)c1ccc2c(c1)C(C)(C)c1cc(C(C)(C)C)ccc1N2c1c(F)c(F)c(-c2c(F)c(F)c(N3c4ccc(C(C)(C)C)cc4C(C)(C)c4cc(C(C)(C)C)ccc43)c(F)c2F)c(F)c1F}	\\
1PXZ-BP	&	\smiles{c1ccc2c(c1)Oc1ccccc1N2c1ccc2c(c1)c1ccccc1c1nc3ccccc3nc21}	&	PXZ-PPO	&	\smiles{Cc1cc2nc(N3c4ccccc4Oc4ccccc43)ccc2c(=O)o1}	\\
PTZ-DBTO2	&	\smiles{O=S1(=O)c2ccccc2-c2cc(N3c4ccccc4Sc4ccccc43)ccc21}	&	PXZ-BOO	&	\smiles{Cc1nc2cc(N3c4ccccc4Oc4ccccc43)ccc2c(=O)o1}	\\
BACQ	&	\smiles{CC1(C)c2c(ccc3ccccc23)-n2c3ccccc3c(=O)c3cccc1c32}	&	PPZTPI	&	\smiles{CC(C)(C)c1ccc(-n2c(-c3ccc(N4c5ccccc5N(c5ccccc5)c5ccccc54)cc3)nc3c4ccccc4c4ccccc4c32)cc1}	\\
2Cz-DMAC-TXO	&	\smiles{CC1(C)c2cc(-n3c4ccccc4c4ccccc43)ccc2N(c2ccc3sc4ccccc4c(=O)c3c2)c2ccc(-n3c4ccccc4c4ccccc43)cc21}	&	PPZPPI	&	\smiles{c1ccc(N2c3ccccc3N(c3ccc(-c4nc5c6ccccc6c6ccccc6c5n4-c4ccccc4)cc3)c3ccccc32)cc1}	\\
2Cz-DMAC-TTR	&	\smiles{CC1(C)c2cc(-n3c4ccccc4c4ccccc43)ccc2N(c2ccc3c(c2)S(=O)(=O)c2ccccc2S3(=O)=O)c2ccc(-n3c4ccccc4c4ccccc43)cc21}	&	Pra-2DMAC	&	\smiles{CC1(C)c2ccccc2N(c2ccnc(N3c4ccccc4C(C)(C)c4ccccc43)c2)c2ccccc21}	\\
2Cz-DMAC-BTB	&	\smiles{CC1(C)c2cc(-n3c4ccccc4c4ccccc43)ccc2N(c2ccc3c(c2)-c2ccccc2S3(=O)=O)c2ccc(-n3c4ccccc4c4ccccc43)cc21}	&	BAPCN	&	\smiles{CC1(C)c2ccccc2N(c2ccccc2)c2ccc3cc(-c4ccc(C\#N)cc4)ccc3c21}	\\
DMAC-BPP	&	\smiles{CC1(C)c2ccccc2N(c2cc(-c3cc(-c4ccc(-c5ccccc5)cc4)nc(-c4ccccc4)n3)cc(N3c4ccccc4C(C)(C)c4ccccc43)c2)c2ccccc21}	&	BACN	&	\smiles{CC1(C)c2ccccc2N(c2ccccc2)c2ccc3cc(C\#N)ccc3c21}	\\
Me-MOC	&	\smiles{COc1ccc2c(c1)c1ccccc1n2-c1ccc(-c2c(C\#N)c(C)nc(C)c2C\#N)cc1}	&	BDQC-2	&	\smiles{c1ccc(-c2ccc(N(c3ccc(-c4ccccc4)cc3)c3ccc4c5ccccc5n(-c5ccc6nc(-c7ccccc7)c(-c7ccccc7)nc6c5)c4c3)cc2)cc1}	\\
Me-DOC	&	\smiles{COc1ccc2c(c1)c1cc(OC)ccc1n2-c1ccc(-c2c(C\#N)c(C)nc(C)c2C\#N)cc1}	&	NIDPA-1	&	\smiles{c1ccc(N(c2ccccc2)c2ccc3c(c2)cc2c4cccc(N(c5ccccc5)c5ccccc5)c4n4c5ccccc5c3c24)cc1}	\\
3DCPO	&	\smiles{O=P(c1ccccc1)(c1ccc2c(c1)c1ccccc1n2-c1ccccc1)c1ccc2c(c1)c1ccccc1n2-c1ccccc1}	&	NIDPA-2	&	\smiles{c1ccc(N(c2ccccc2)c2ccc3c(c2)c2cc(N(c4ccccc4)c4ccccc4)cc4c5c6ccccc6ccc5n3c24)cc1}	\\
TDBA-TPDICz	&	\smiles{CC(C)(C)c1ccc2c(c1)B1c3cc(C(C)(C)C)ccc3Oc3cc(-n4c5ccccc5c5c(-c6ccccc6)c6c(c(-c7ccccc7)c54)c4ccccc4n6-c4ccccc4)cc(c31)O2}	&	SABIQ	&	\smiles{c1ccc(N2c3ccccc3C3(c4ccccc4-c4nc5cc6ccccc6cc5nc43)c3ccccc32)cc1}	\\
TRZ-pIC	&	\smiles{CC(C)(C)c1ccc2c(c1)B1c3cc(C(C)(C)C)ccc3Oc3cc(-n4c5ccccc5c5c(-c6ccccc6)c6c(c(-c7ccccc7)c54)c4ccccc4n6-c4ccccc4)cc(c31)O2}	&	Perylene	&	\smiles{c1cc2cccc3c4cccc5cccc(c(c1)c23)c54}	\\
SBF-PXZ	&	\smiles{c1ccc2c(c1)Oc1ccccc1N2c1ccc2c(c1)C1(c3ccccc3-c3ccccc31)c1cc(N3c4ccccc4Oc4ccccc43)ccc1-2}	&	SAIP	&	\smiles{c1ccc(N2c3ccccc3C3(c4ccccc4-c4nccnc43)c3ccccc32)cc1}	\\
PtOEP	&	\smiles{CCC1=C(CC)c2cc3[nH]c(cc4nc(cc5[nH]c(cc1n2)c(CC)c5CC)C(CC)=C4CC)c(CC)c3CC.[Pt+2]}	&	FAC	&	\smiles{c1ccc2c(c1)Nc1ccccc1C21c2ccccc2-c2ccccc21}	\\
PxCYP	&	\smiles{N\#Cc1ccc(-c2ccc(N3c4ccccc4Oc4ccccc43)cc2)nc1-c1ccc(N2c3ccccc3Oc3ccccc32)cc1}	&	phenylperylene	&	\smiles{c1ccc(-c2ccc3cccc4c5cccc6cccc(c2c34)c65)cc1}	\\
AcCYP	&	\smiles{CC1(C)c2ccccc2N(c2ccc(-c3ccc(C\#N)c(-c4ccc(N5c6ccccc6C(C)(C)c6ccccc65)cc4)n3)cc2)c2ccccc21}	&	ImPy-2	&	\smiles{COc1ccc2cc(-c3nc(-c4ccccc4)c4ccccn34)ccc2c1}	\\
BCN	&	\smiles{CC1(C)c2ccccc2N2c3ccc4cc(C\#N)ccc4c3C(C)(C)c3cccc1c32}	&	ImPy-1	&	\smiles{c1ccc(-c2nc(-c3ccc4ccccc4c3)n3ccccc23)cc1}	\\
DPPSPF	&	\smiles{O=S(=O)(c1ccccc1)c1ccc(-c2ccc3c(c2)C(c2ccccc2)(c2ccccc2)c2ccccc2-3)cc1}	&	ImPy-3	&	\smiles{c1ccc(-c2nc(-c3ccc4ccc5cccc6ccc3c4c56)n3ccccc23)cc1}	\\
Me2ACBO	&	\smiles{CC(C)(c1ccc2c(c1)B1c3ccc(N4c5ccccc5C(C)(C)c5ccccc54)cc3Oc3cccc(c31)O2)c1ccc2c(c1)B1c3ccc(N4c5ccccc5C(C)(C)c5ccccc54)cc3Oc3cccc(c31)O2}	&	BTaz-Ar-PXZ	&	\smiles{c1ccc2c(c1)Oc1ccccc1N2c1ccc(-c2nc3ccccc3s2)cc1}	\\
F2AcBO	&	\smiles{CC1(C)c2ccccc2N(c2ccc3c(c2)Oc2cccc4c2B3c2cc(C3(c5ccc6c(c5)B5c7ccc(N8c9ccccc9C(C)(C)c9ccccc98)cc7Oc7cccc(c75)O6)c5ccccc5-c5ccccc53)ccc2O4)c2ccccc21}	&	FSF4A	&	\smiles{CC1(C)c2ccccc2-c2ccc(N(c3ccccc3-c3ccccc3)c3cccc4c3-c3ccccc3C43c4ccccc4-c4ccccc43)cc21}	\\
PSPSF	&	\smiles{O=S(=O)(c1ccccc1)c1ccc(-c2ccc3c(c2)C2(c4ccccc4-c4ccccc42)c2ccccc2-3)cc1}	&	TPA-L-BN	&	\smiles{N\#Cc1ccc(CCCc2ccc(N(c3ccccc3)c3ccccc3)cc2)cc1}	\\
TXO-P-Si	&	\smiles{O=C1c2ccccc2S(=O)(=O)c2ccc(-c3ccc([Si](c4ccccc4)(c4ccccc4)c4ccccc4)cc3)cc21}	&	IpCm-PhBzAc	&	\smiles{CC(C)c1ccc2occ(-c3ccc(N4c5ccccc5C(C)(C)c5cc6c(cc54)oc4ccccc46)cc3)c(=O)c2c1}	\\
CzB-FMPIM	&	\smiles{FC(F)(F)c1ccc(-n2c(-c3ccc(-c4ccc(-n5c6ccccc6c6ccccc65)cc4)cc3)nc(-c3ccccc3)c2-c2ccccc2)cc1}	&	DphAn-Br	&	\smiles{O=C1c2ccccc2C(c2ccccc2)(c2ccc(Br)cc2)c2ccccc21}	\\
CzB-FMPPI	&	\smiles{FC(F)(F)c1ccc(-n2c(-c3ccc(-c4ccc(-n5c6ccccc6c6ccccc65)cc4)cc3)nc3c4ccccc4c4ccccc4c32)cc1}	&	5-DBFBI	&	\smiles{c1ccc2c(c1)nc1n(-c3ccc(-c4ccc5oc6ccccc6c5c4)cn3)c3ccccc3n21}	\\
BACF	&	\smiles{CC1(C)c2ccccc2N(c2ccccc2)c2ccc3cc(F)ccc3c21}	&	4-DBFBI	&	\smiles{c1ccc2c(c1)nc1n(-c3cc(-c4ccc5oc6ccccc6c5c4)ccn3)c3ccccc3n21}	\\
BACH	&	\smiles{CC1(C)c2ccccc2N(c2ccccc2)c2ccc3ccccc3c21}	&	DMACMNPTO	&	\smiles{Cc1cc(N2c3ccccc3S(=O)(=O)c3ccccc32)c(C)cc1N1c2ccccc2C(C)(C)c2ccccc21}	\\
PzTDBA	&	\smiles{CC(C)(C)c1ccc2c(c1)B1c3cc(C(C)(C)C)ccc3Oc3cc(N4c5ccccc5N(c5cc6c7c(c5)Oc5ccc(C(C)(C)C)cc5B7c5cc(C(C)(C)C)ccc5O6)c5ccccc54)cc(c31)O2}	&	1CzSO	&	\smiles{CC(C)(C)c1ccc2c(c1)c1cc(C(C)(C)C)ccc1n2-c1ccc2c(c1)-c1cc(-n3c4ccc(C(C)(C)C)cc4c4cc(C(C)(C)C)ccc43)ccc1S2(=O)=O}	\\
3TCPM	&	\smiles{O=C(c1ccccc1)c1ccc(-n2c3ccccc3c3cc(-c4ccc(N(c5ccccc5)c5ccccc5)cc4)ccc32)cc1}	&	DpAn-BzAc	&	\smiles{CC1(C)c2ccccc2N(c2ccc(C3(c4ccccc4)c4ccccc4C(=O)c4ccccc43)cc2)c2c1ccc1c2oc2ccccc21}	\\
TPA-PPO	&	\smiles{c1ccc(-c2nc(-c3ccc(-c4ccc(N(c5ccccc5)c5ccccc5)cc4)cn3)oc2-c2ccccc2)cc1}	&	BzITz	&	\smiles{N\#Cc1cc(-c2nc(-c3ccccc3)nc(-c3ccccc3)n2)ccc1-n1c2ccccc2n2c3ccccc3nc12}	\\
CzPh-PPO	&	\smiles{c1ccc(-c2nc(-c3ccc(-c4ccc(-n5c6ccccc6c6ccccc65)cc4)cn3)oc2-c2ccccc2)cc1}	&	TPXZ–as-TAZ	&	\smiles{c1ccc2c(c1)Oc1ccccc1N2c1ccc(-c2nnc(-c3ccc(N4c5ccccc5Oc5ccccc54)cc3)c(-c3ccc(N4c5ccccc5Oc5ccccc54)cc3)n2)cc1}	\\
m-CF3PIBI	&	\smiles{FC(F)(F)c1cccc(-n2c(-c3ccc(-c4ccc(-c5ccc(-c6nc7ccccc7n6-c6ccccc6)cc5)cc4)cc3)nc3c4ccccc4c4ccccc4c32)c1}	&	TPA-CN-N4-2PY	&	\smiles{CC(C)(C)c1ccc(N(c2ccc(-c3cc(N)c(N)cc3C\#N)cc2)c2ccc(C(C)(C)C)cc2)cc1}	\\
DPTPEA	&	\smiles{c1ccc(C(=C(c2ccccc2)c2ccc(N(c3ccccc3)c3ccccc3)cc2)c2ccccc2)cc1}	&	TPA-CN-N4	&	\smiles{c1ccc2nc3c4cccnc4c4ncccc4c3nc2c1}	\\
PICNAnCz	&	\smiles{N\#Cc1ccc(-n2c(-c3ccc(-c4c5ccccc5c(Br)c5ccccc45)cc3)nc3c4ccccc4c4ccccc4c32)cc1}	&	6,7-DCNQx-DICz	&	\smiles{N\#Cc1cc2nc(-c3ccccc3)c(-c3ccc(-n4c5ccccc5c5c6c(c7ccccc7n6-c6ccccc6)c6c(c7ccccc7n6-c6ccccc6)c54)cc3)nc2cc1C\#N}	\\
DPA-DPQ	&	\smiles{c1ccc(-c2cc3ncc(N(c4ccccc4)c4ccccc4)nc3cc2-c2ccccc2)cc1}	&	PBICT	&	\smiles{c1ccc(-c2nc(-n3c4ccccc4c4ccc5c6ccccc6n(-c6ccccc6)c5c43)nc(-n3c4ccccc4c4ccc5c6ccccc6n(-c6ccccc6)c5c43)n2)cc1}	\\
PhQLPXZ	&	\smiles{c1ccc(-c2cc(-c3ccccc3)c3cc(N4c5ccccc5Oc5ccccc54)ccc3n2)cc1}	&	DCBNPE	&	\smiles{O=[N+]([O-])c1ccc(/C(=C(/c2ccc([N+](=O)[O-])cc2)n2c3ccccc3c3ccccc32)n2c3ccccc3c3ccccc32)cc1}	\\
2DPA-DPQ	&	\smiles{c1ccc(-c2cc3nc(N(c4ccccc4)c4ccccc4)c(N(c4ccccc4)c4ccccc4)nc3cc2-c2ccccc2)cc1}	&	TPA-BT-Q	&	\smiles{Cc1ccc(N(c2ccc(C)cc2)c2ccc(C3(c4ccc(N(c5ccc(C)cc5)c5ccc(C)cc5)cc4)CCCCC3)cc2)cc1}	\\
FDQPXZ	&	\smiles{Fc1cc2nc(-c3ccccc3)c(-c3ccccc3)nc2cc1N1c2ccccc2Oc2ccccc21}	&	D5	&	\smiles{c1cc(-c2ccc(-n3c4ccccc4c4ccccc43)cc2)cc(-c2ccc3c4ccccc4c4ccccc4c3c2)c1}	\\
DTRZSFX	&	\smiles{c1ccc(-c2nc(-c3ccccc3)nc(-c3ccc4c(c3)C3(c5ccccc5Oc5ccccc53)c3cc(-c5nc(-c6ccccc6)nc(-c6ccccc6)n5)ccc3-4)n2)cc1}	&	D6	&	\smiles{c1cc(-c2ccc3c4ccccc4c4ccccc4c3c2)cc(-n2c3ccccc3c3ccccc32)c1}	\\
E-o-DHMPA	&	\smiles{OCc1ccccc1-c1c2ccccc2c(-c2ccccc2CO)c2ccccc12}	&	PxPYM	&	\smiles{c1ccc2c(c1)Oc1ccccc1N2c1ccc(-c2cnc(-c3ccc(N4c5ccccc5Oc5ccccc54)cc3)nc2)cc1}	\\
Z-o-DHMPA	&	\smiles{OCc1ccccc1-c1c2ccccc2c(-c2ccccc2CO)c2ccccc12}	&	D1–D6	&	\smiles{c1cc(-c2ccc3c4ccccc4c4ccccc4c3c2)cc(-n2c3ccccc3c3ccccc32)c1}	\\
p-DHMPA	&	\smiles{OCc1ccc(-c2c3ccccc3c(-c3ccc(CO)cc3)c3ccccc23)cc1}	&	DTPCZTZ	&	\smiles{CC(C)(C)c1ccc2c(c1)c1cc(C(C)(C)C)ccc1n2-c1ccc(-c2ccc(-c3nnc(-c4ccccc4)n3-c3ccccc3)cc2)cc1}	\\
STAC-BP-STAC	&	\smiles{O=C(c1ccc(N2c3ccccc3C3(c4ccccc4Sc4ccccc43)c3ccccc32)cc1)c1ccc(N2c3ccccc3C3(c4ccccc4Sc4ccccc43)c3ccccc32)cc1}	&	DHPZ-2BI	&	\smiles{c1ccc(-n2c(-c3ccc(N4c5ccccc5N(c5ccc(-c6nc7ccccc7n6-c6ccccc6)cc5)c5ccccc54)cc3)nc3ccccc32)cc1}	\\
SXAC-BP-SXAC	&	\smiles{O=C(c1ccc(N2c3ccccc3C3(c4ccccc4Oc4ccccc43)c3ccccc32)cc1)c1ccc(N2c3ccccc3C3(c4ccccc4Oc4ccccc43)c3ccccc32)cc1}	&	DHPZ-2BN	&	\smiles{N\#Cc1ccc(N2c3ccccc3N(c3ccc(C\#N)cc3)c3ccccc32)cc1}	\\
m-DHMPA	&	\smiles{OCc1cccc(-c2c3ccccc3c(-c3cccc(CO)c3)c3ccccc23)c1}	&	SAF-2NP	&	\smiles{c1ccc2c(c1)-c1ccccc1C21c2ccccc2N(c2ccc3nc4c5cccnc5c5ncccc5c4nc3c2)c2ccccc21}	\\
SFAC-BP-SFAC	&	\smiles{O=C(c1ccc(N2c3ccccc3C3(c4ccccc4-c4ccccc43)c3ccccc32)cc1)c1ccc(N2c3ccccc3C3(c4ccccc4-c4ccccc43)c3ccccc32)cc1}	&	BTPO	&	\smiles{O=S1(=O)c2ccccc2N2c3ccccc3S(=O)(=O)c3cccc1c32}	\\
PBTPA	&	\smiles{CCN1c2ccccc2Sc2cc(-c3ccc(-c4ccc(N(c5ccccc5)c5ccccc5)cc4)c4nsnc34)ccc21}	&	PyDCN–DMAC	&	\smiles{CC1(C)c2ccccc2N(c2ccc(-c3cc(-c4ccc(C\#N)cc4)nc(-c4ccc(C\#N)cc4)c3)cc2)c2ccccc21}	\\
HDP	&	\smiles{O=C1c2ccccc2C(=O)c2nc3c(cc(O)c4ccccc43)nc21}	&	PyDCN	&	\smiles{N\#Cc1ccc(-c2cc(-c3ccccc3)cc(-c3ccc(C\#N)cc3)n2)cc1}	\\
PY-TPA	&	\smiles{CC(C)(C)c1ccc(N(c2ccc(-c3cc4nc5c6ccc(C\#N)cc6c6cc(C\#N)ccc6c5nc4cc3-c3cccnc3)cc2)c2ccc(C(C)(C)C)cc2)cc1}	&	PXZ-OXD	&	\smiles{c1ccc(-c2nnc(-c3ccc(N4c5ccccc5Oc5ccccc54)cc3)o2)cc1}	\\
CN-TPA	&	\smiles{CC(C)(C)c1ccc(N(c2ccc(-c3cc4nc5c6ccc(C\#N)cc6c6cc(C\#N)ccc6c5nc4cc3C\#N)cc2)c2ccc(C(C)(C)C)cc2)cc1}	&	TCz-3PA-TCz	&	\smiles{CC(C)(C)c1ccc(-c2ccc3c(c2)c2cc(-c4cnc(-c5ccc6c(c5)c5cc(-c7ccc(C(C)(C)C)cc7)ccc5n6-c5ccc(C(C)(C)C)cc5)cn4)ccc2n3-c2ccc(C(C)(C)C)cc2)cc1}	\\
DMIC-TRZ	&	\smiles{CC1(C)CC(c2cccc(-c3nc(-c4ccccc4)nc(-c4ccccc4)n3)c2)C=c2[nH]c3cc4c(cc3c21)-c1ccccc1C=4}	&	NAI–BiFA	&	\smiles{Cc1c(-c2ccccc2)cc(N2C(=O)c3cccc4c(N(c5ccc(-c6ccccc6)cc5)c5ccc6c(c5)C(C)(C)c5ccccc5-6)ccc(c34)C2=O)cc1-c1ccccc1}	\\
NyDPAc	&	\smiles{CC1(C)c2ccccc2Nc2cccc(-c3ccc(-c4ccc5nc(-c6ccc(-c7cccc8c7C(C)(C)c7ccccc7N8)cc6)ccc5n4)cc3)c21}	&	NAI–PhBiFA	&	\smiles{Cc1c(-c2ccccc2)cc(N2C(=O)c3cccc4c(-c5ccc(N(c6ccc(-c7ccccc7)cc6)c6ccc7c(c6)C(C)(C)c6ccccc6-7)cc5)ccc(c34)C2=O)cc1-c1ccccc1}	\\
oCBP	&	\smiles{c1ccc(-n2c3ccccc3c3ccccc32)c(-c2ccccc2-n2c3ccccc3c3ccccc32)c1}	&	pDTCz-2DPyS	&	\smiles{CC(C)(C)c1ccc2c(c1)c1cc(C(C)(C)C)ccc1n2-c1ccc(S(=O)(=O)c2ccc(-n3c4ccc(C(C)(C)C)cc4c4cc(C(C)(C)C)ccc43)cn2)nc1}	\\
PABP	&	\smiles{Cc1cc(-c2c3ccccc3c(-c3ccccc3)c3ccccc23)ccc1-c1ccc(-c2c3ccccc3c(-c3ccccc3)c3ccccc23)cc1C}	&	2SPBAC-BP	&	\smiles{O=C(c1ccc(N2c3ccccc3C3(c4ccccc4-c4ccccc43)c3cc4ccccc4cc32)cc1)c1ccc(N2c3ccccc3C3(c4ccccc4-c4ccccc43)c3cc4ccccc4cc32)cc1}	\\
DHID-DPS	&	\smiles{CCn1c2ccccc2c2c1c1ccccc1n2-c1ccc(S(=O)(=O)c2ccc(-n3c4ccccc4c4c3c3ccccc3n4CC)cc2)cc1}	&	SPBAC-PH	&	\smiles{c1ccc(N2c3ccccc3C3(c4ccccc4-c4ccccc43)c3cc4ccccc4cc32)cc1}	\\
4FBTPI	&	\smiles{CC(C)(C)c1ccc(-n2c(-c3ccc(-c4c(F)c(F)c(-c5ccc(-c6nc7c8ccccc8c8ccccc8c7n6-c6ccc(C(C)(C)C)cc6)cc5)c(F)c4F)cc3)nc3c4ccccc4c4ccccc4c32)cc1}	&	DFPBI	&	\smiles{CC1(C)c2cc(-c3nc(-c4ccccc4)c(-c4ccccc4)n3-c3ccccc3)ccc2-c2ccc(N(c3ccccc3)c3ccccc3)cc21}	\\
BCzBMes	&	\smiles{Cc1cc(C)c(B(c2c(C)cc(C)cc2C)n2c3ccccc3c3cc(-c4ccc5c(c4)c4ccccc4n5-c4ccccc4)ccc32)c(C)c1}	&	TDPA–TRZ	&	\smiles{c1ccc(N(c2ccccc2)c2ccc(-c3nc(-c4ccc(N(c5ccccc5)c5ccccc5)cc4)nc(-c4ccc(N(c5ccccc5)c5ccccc5)cc4)n3)cc2)cc1}	\\
CCDMB	&	\smiles{Cc1cc(C)c(B(c2ccc(-n3c4ccccc4c4cc(-n5c6ccccc6c6ccccc65)ccc43)cc2)c2c(C)cc(C)cc2C)c(C)c1}	&	DMAC-11-DPPZ	&	\smiles{CC1(C)c2ccccc2N(c2ccc3nc4c5ccccc5c5ccccc5c4nc3c2)c2ccccc21}	\\
PCDMB	&	\smiles{Cc1cc(C)c(B(c2ccc(-n3c4ccccc4c4cc(N5c6ccccc6Oc6ccccc65)ccc43)cc2)c2c(C)cc(C)cc2C)c(C)c1}	&	PXZ-11-DPPZ	&	\smiles{c1ccc2c(c1)Oc1ccccc1N2c1ccc2nc3c4ccccc4c4ccccc4c3nc2c1}	\\
BTrztCz	&	\smiles{CC(C)(C)c1ccc2c(c1)c1cc(C(C)(C)C)ccc1n2-c1ccccc1-c1nc(-c2ccccc2)nc(-c2ccc(-c3nc(-c4ccccc4)nc(-c4ccccc4-n4c5ccc(C(C)(C)C)cc5c5cc(C(C)(C)C)ccc54)n3)cc2)n1}	&	PXZ-10-DPPZ	&	\smiles{c1ccc2c(c1)Oc1ccccc1N2c1cccc2nc3c4ccccc4c4ccccc4c3nc12}	\\
DHID-DBS	&	\smiles{CCCCn1c2ccccc2c2c1c1ccccc1n2-c1ccc2c(c1)-c1cc(-n3c4ccccc4c4c3c3ccccc3n4CCCC)ccc1S2(=O)=O}	&	CCC	&	\smiles{c1ccc2c(c1)Cc1c-2ccc2c1-c1ccccc1C2}	\\
CZ-DPS-DMAC	&	\smiles{CC1(C)c2ccccc2N(c2ccc(S(=O)(=O)c3cc(-n4c5ccccc5c5ccccc54)cc(-n4c5ccccc5c5ccccc54)c3)cc2)c2ccccc21}	&	CzmPPC	&	\smiles{c1ccc(-c2cccc(-c3nc(-c4cccc(-c5ccccc5)c4)nc(-c4cccc(-c5ccccc5)c4)n3)c2)cc1}	\\
XT-T	&	\smiles{Cc1ccc(N(c2ccc(C)cc2)c2ccccc2C2(c3ccccc3N(c3ccc(C)cc3)c3ccc(C)cc3)CCCCC2)cc1}	&	C2FLA-1	&	\smiles{CCCC1(CCC)c2ccccc2-c2ccc(C\#Cc3ccc4c(c3)C(CCC)(CCC)c3cc(C\#N)ccc3-4)cc21}	\\
3DPTIA	&	\smiles{CC1(C)c2ccccc2N(c2nc(-c3ccccc3)nc(-c3ccccc3)n2)c2cc3c4ccccc4n(-c4ccccc4)c3cc21}	&	MDP3FL	&	\smiles{Cc1cccc(N(c2ccccc2)c2ccc3c(c2)C(C)(C)c2cc(-c4ccc5c(c4)C(C)(C)c4cc(-c6ccc7c(c6)C(C)(C)c6cc(N(c8ccccc8)c8cccc(C)c8)ccc6-7)ccc4-5)ccc2-3)c1}	\\
PyIAnTPh	&	\smiles{c1ccc(-c2cc(-c3ccccc3)cc(-c3c4ccccc4c(-c4ccc(-c5nc6c7cccc8ccc9cccc(c9c87)c6n5-c5ccccc5)cc4)c4ccccc34)c2)cc1}	&	C3FLA-2	&	\smiles{CCCCC1(CCCC)c2cc(C\#Cc3ccc4c(c3)C(CCC)(CCC)c3cc(C\#N)ccc3-4)ccc2-c2ccc(C\#Cc3ccc4c(c3)C(CCC)(CCC)c3cc(C\#N)ccc3-4)cc21}	\\
2CzP-INN	&	\smiles{N\#Cc1c(-c2ccc(-n3c4ccccc4c4ccccc43)cc2)cncc1-c1ccc(-n2c3ccccc3c3ccccc32)cc1}	&	9CzFDPEPO	&	\smiles{O=P(c1ccccc1)(c1ccccc1)c1ccccc1Oc1ccc(-c2ccc(C3(c4ccc(-n5c6ccccc6c6ccccc65)cc4)c4ccccc4-c4ccccc43)cc2)cc1P(=O)(c1ccccc1)c1ccccc1}	\\
VBN0	&	\smiles{C=Cc1ccc(COc2ccc3c(c2)c2ccccc2n3-c2cccc(-n3c4ccccc4c4cc(OCc5ccc(C=C)cc5)ccc43)c2C\#N)cc1}	&	9CzFDPESPO	&	\smiles{O=P(c1ccccc1)(c1ccccc1)c1ccccc1Oc1ccc(-c2ccc(C3(c4ccc(-n5c6ccccc6c6ccccc65)cc4)c4ccccc4-c4ccccc43)cc2)cc1}	\\
PIAnTPh	&	\smiles{c1ccc(-c2cc(-c3ccccc3)cc(-c3c4ccccc4c(-c4ccc(-c5nc6c7ccccc7c7ccccc7c6n5-c5ccccc5)cc4)c4ccccc34)c2)cc1}	&	CTM	&	\smiles{O=C(c1ccc(-n2c3ccccc3c3ccccc32)cc1)c1ccc2sc3ccccc3c2c1}	\\
BCZ-DPS-AD	&	\smiles{CC1(C)c2ccccc2N(c2cc(N3c4ccccc4C(C)(C)c4ccccc43)cc(S(=O)(=O)c3ccc4c(c3)c3ccccc3n4-c3ccccc3)c2)c2ccccc21}	&	tCTM	&	\smiles{CC(C)(C)c1ccc2c(c1)c1cc(C(C)(C)C)ccc1n2-c1ccc(C(=O)c2ccc3sc4ccccc4c3c2)cc1}	\\
2BTPA-INN	&	\smiles{N\#Cc1c(-c2ccc(N(c3ccc(-c4ccccc4)cc3)c3ccc(-c4ccccc4)cc3)cc2)cncc1-c1ccc(N(c2ccc(-c3ccccc3)cc2)c2ccc(-c3ccccc3)cc2)cc1}	&	PSBC	&	\smiles{CC(C)(C)c1ccc(-n2c(-c3ccc(-c4ccc(S(=O)(=O)c5ccc(-c6ccc(-n7c8ccccc8c8ccccc87)cc6)cc5)cc4)cc3)nc3c4ccccc4c4ccccc4c32)cc1}	\\
SpiroAc	&	\smiles{c1ccc2c(c1)Nc1ccccc1C21c2ccccc2-c2ccccc21}	&	26DAcBPy	&	\smiles{CC1(C)c2ccccc2N(c2ccc(C(=O)c3cccc(C(=O)c4ccc(N5c6ccccc6C(C)(C)c6ccccc65)cc4)n3)cc2)c2ccccc21}	\\
BuPCzBCO	&	\smiles{Cc1ccc2c3cc(Br)c(C)cc3n(-c3ccc(C(C)(C)C)cc3)c2c1}	&	25DAcBPy	&	\smiles{CC1(C)c2ccccc2N(c2ccc(C(=O)c3ccc(C(=O)c4ccc(N5c6ccccc6C(C)(C)c6ccccc65)cc4)nc3)cc2)c2ccccc21}	\\
AQ-S-CZ	&	\smiles{O=C1c2ccccc2C(=O)c2cc(Sc3ccc(-n4c5ccccc5c5ccccc54)cc3)ccc21}	&	26DPXZBPy	&	\smiles{O=C(c1ccc(N2c3ccccc3Oc3ccccc32)cc1)c1cccc(C(=O)c2ccc(N3c4ccccc4Oc4ccccc43)cc2)n1}	\\
DBP-3TPA	&	\smiles{c1ccc(N(c2ccccc2)c2ccc(-c3ccc4nc5c6ccc(-c7ccc(N(c8ccccc8)c8ccccc8)cc7)cc6c6cc(-c7ccc(N(c8ccccc8)c8ccccc8)cc7)ccc6c5nc4c3)cc2)cc1}	&	BTITrz	&	\smiles{N\#Cc1cc(-c2nc(-c3ccccc3)nc(-c3ccccc3)n2)ccc1-n1c2ccccc2c2c3ccccc3sc21}	\\
DCzIPN	&	\smiles{N\#Cc1cc(C\#N)c(-n2c3ccccc3c3ccccc32)cc1-n1c2ccccc2c2ccccc21}	&	mTPA–PPI	&	\smiles{c1ccc(N(c2ccccc2)c2ccc(-c3cccc(-c4nc5c6ccccc6c6ccccc6c5n4-c4ccccc4)c3)cc2)cc1}	\\
DMeCzIPN	&	\smiles{Cc1cccc2c3ccccc3n(-c3cc(-n4c5ccccc5c5cccc(C)c54)c(C\#N)cc3C\#N)c12}	&	PPI	&	\smiles{c1ccc(-c2nc3c4ccccc4c4ccccc4c3n2-c2ccccc2)cc1}	\\
SAFDPA	&	\smiles{c1ccc(N(c2ccccc2)c2ccc3c(c2)-c2ccccc2C32c3ccccc3N(c3ccccc3)c3ccccc32)cc1}	&	CzTPA-m-Trz	&	\smiles{c1ccc(-c2ccc(N(c3ccc(-c4ccc5c(c4)c4ccccc4n5-c4ccccc4)cc3)c3cccc(-c4nc(-c5ccccc5)nc(-c5ccccc5)n4)c3)cc2)cc1}	\\
t-DABNA	&	\smiles{CC(C)(C)c1ccc(N2c3ccc(C(C)(C)C)cc3B3c4cc(C(C)(C)C)ccc4N(c4ccc(C(C)(C)C)cc4)c4cccc2c43)cc1}	&	TPA–2PPITPA	&	\smiles{c1ccc(N(c2ccc(-c3ccc(-c4nc5c6ccccc6c6ccccc6c5n4-c4ccc(N(c5ccccc5)c5ccccc5)cc4)cc3)cc2)c2ccc(-c3ccc(-c4nc5c6ccccc6c6ccccc6c5n4-c4ccc(N(c5ccccc5)c5ccccc5)cc4)cc3)cc2)cc1}	\\
3-AcBPtCN	&	\smiles{CC1(C)c2ccccc2N(c2ccc3c(c2)c2cc(C\#N)ccc2c2nc4cc(C\#N)c(C\#N)cc4nc32)c2ccccc21}	&	TPA–2SPPITPA	&	\smiles{C(=C/c1nc2c3ccccc3c3ccccc3c2n1-c1ccc(N(c2ccccc2)c2ccccc2)cc1)$\backslash$c1ccc(-c2ccc(N(c3ccccc3)c3ccc(-c4ccc(/C=C/c5nc6c7ccccc7c7ccccc7c6n5-c5ccc(N(c6ccccc6)c6ccccc6)cc5)cc4)cc3)cc2)cc1}	\\
3-AcBPdCN	&	\smiles{CC1(C)c2ccccc2N(c2ccc3c(c2)c2ccccc2c2nc4cc(C\#N)c(C\#N)cc4nc32)c2ccccc21}	&	2Cz2tCzBn	&	\smiles{CC(C)(C)c1ccc2c(c1)c1cc(C(C)(C)C)ccc1n2-c1c(-n2c3ccccc3c3ccccc32)cc(-n2c3ccccc3c3ccccc32)c(-n2c3ccc(C(C)(C)C)cc3c3cc(C(C)(C)C)ccc32)c1C\#N}	\\
11-AcBPdCN	&	\smiles{CC1(C)c2ccccc2N(c2ccc3nc4c5ccc(C\#N)cc5c5cc(C\#N)ccc5c4nc3c2)c2ccccc21}	&	2tCz2CzBn	&	\smiles{CC(C)(C)c1ccc2c(c1)c1cc(C(C)(C)C)ccc1n2-c1cc(-n2c3ccc(C(C)(C)C)cc3c3cc(C(C)(C)C)ccc32)c(-n2c3ccccc3c3ccccc32)c(C\#N)c1-n1c2ccccc2c2ccccc21}	\\
PID-3	&	\smiles{c1ccc(-n2c(-c3ccc(-c4cccs4)s3)nc3c4ccccc4c4ccccc4c32)cc1}	&	PmPmP	&	\smiles{c1cncc(-c2nc(-c3ccc(N4c5ccccc5Oc5ccccc54)cc3)cc(-c3cncnc3)n2)c1}	\\
3-ACBPmCN	&	\smiles{CC1(C)c2ccccc2N(c2ccc3c(c2)c2cc(C\#N)ccc2c2nc4ccccc4nc32)c2ccccc21}	&	PLQYs	&	\smiles{CC(C)(C)c1ccc2c(c1)c1cc(C(C)(C)C)ccc1n2-c1cc2ccc3cc(-n4c5ccc(C(C)(C)C)cc5c5cc(C(C)(C)C)ccc54)cc4ccc(c1)c2c34}	\\
3-DTPA-BBI	&	\smiles{O=c1c2ccc(-c3ccc(N(c4ccccc4)c4ccccc4)cc3)c3c(-c4ccc(N(c5ccccc5)c5ccccc5)cc4)ccc(c32)c2nc3ccccc3n12}	&	PXZ-ND	&	\smiles{c1ccc2c(c1)Oc1ccccc1N2c1ccc2ccc(N3c4ccccc4Oc4ccccc43)nc2n1}	\\
IABTCN	&	\smiles{CC1(C)c2ccccc2-c2cc3c(cc21)N(c1ccccc1)c1ccc(-c2ccc(C\#N)c4nsnc24)cc1C3(c1ccccc1)c1ccccc1}	&	CPPyC	&	\smiles{c1cc(-c2cncc(-n3c4ccccc4c4ccccc43)c2)cc(-n2c3ccccc3c3ccccc32)c1}	\\
BTZPP	&	\smiles{CC(C)(C)c1ccc(-n2c(-c3ccc(-c4cccc5nsnc45)cc3)cc3c2cc(-c2ccc(-c4cccc5nsnc45)cc2)n3-c2ccc(C(C)(C)C)cc2)cc1}	&	DMAC-ND	&	\smiles{CC1(C)c2ccccc2N(c2ccc3ccc(N4c5ccccc5C(C)(C)c5ccccc54)nc3n2)c2ccccc21}	\\
DPimdPPA	&	\smiles{N\#C/C(=C$\backslash$c1ccc(-c2nc(-c3ccccc3)c(-c3ccccc3)[nH]2)cc1)c1ccccc1}	&	PrDPhAc	&	\smiles{c1ccc(C2(c3ccccc3)c3ccccc3N(c3cncc(N4c5ccccc5C(c5ccccc5)(c5ccccc5)c5ccccc54)n3)c3ccccc32)cc1}	\\
BCzPh	&	\smiles{c1ccc(-n2c3ccccc3c3cc(-c4ccc5c(c4)c4ccccc4n5-c4ccccc4)ccc32)cc1}	&	2Mi2SB	&	\smiles{Cc1nc2ccccc2n1-c1cc(S(=O)(=O)c2ccc(C(C)(C)C)cc2)c(-n2c(C)nc3ccccc32)cc1S(=O)(=O)c1ccc(C(C)(C)C)cc1}	\\
TBPM2CN	&	\smiles{N\#Cc1ccc(-c2ccc3c(c2)c2cc(-c4ccc(C\#N)cc4)ccc2c2c3nc(-c3ccc(-c4ccc(N(c5ccccc5)c5ccccc5)cc4)cc3)n2-c2ccccc2)cc1}	&	BFDMAc-NAI	&	\smiles{CC(C)(C)c1ccc(N2C(=O)c3cccc4c(N5c6ccccc6C(C)(C)c6ccc7c(oc8ccccc87)c65)ccc(c34)C2=O)cc1}	\\
TPM2CN	&	\smiles{N\#Cc1ccc(-c2ccc3c(c2)c2cc(-c4ccc(C\#N)cc4)ccc2c2c3nc(-c3ccc(N(c4ccccc4)c4ccccc4)cc3)n2-c2ccccc2)cc1}	&	HPIPXZ	&	\smiles{c1ccc(N(c2ccc(-c3ccc(N(c4ccccc4)c4cccc5ccccc45)cc3)cc2)c2cccc3ccccc23)cc1}	\\
PT-BZ-DMAC	&	\smiles{CC1(C)c2ccccc2N(c2ccc(C(=O)c3ccc(Sc4ccccc4)cc3)cc2)c2ccccc21}	&	3BPy-mDTC	&	\smiles{CC(C)(C)c1ccc2c(c1)c1cc(C(C)(C)C)ccc1n2-c1cc(C(=O)c2cccnc2)cc(-n2c3ccc(C(C)(C)C)cc3c3cc(C(C)(C)C)ccc32)c1}	\\
PS-BZ-DMAC	&	\smiles{CC1(C)c2ccccc2N(c2ccc(C(=O)c3ccc(S(=O)(=O)c4ccccc4)cc3)cc2)c2ccccc21}	&	2BPy-mDTC	&	\smiles{CC(C)(C)c1ccc2c(c1)c1cc(C(C)(C)C)ccc1n2-c1cc(C(=O)c2ccccn2)cc(-n2c3ccc(C(C)(C)C)cc3c3cc(C(C)(C)C)ccc32)c1}	\\
2PCzBN-tPh	&	\smiles{CC(C)(C)c1ccc(-c2cc(-n3c4ccc(-c5ccccc5)cc4c4cc(-c5ccccc5)ccc43)c(C\#N)c(-n3c4ccc(-c5ccccc5)cc4c4cc(-c5ccccc5)ccc43)c2)cc1}	&	BP-mDTC	&	\smiles{CC(C)(C)c1ccc2c(c1)c1cc(C(C)(C)C)ccc1n2-c1cc(C(=O)c2ccccc2)cc(-n2c3ccc(C(C)(C)C)cc3c3cc(C(C)(C)C)ccc32)c1}	\\
2PCzBN	&	\smiles{N\#Cc1c(-n2c3ccc(-c4ccccc4)cc3c3cc(-c4ccccc4)ccc32)cccc1-n1c2ccc(-c3ccccc3)cc2c2cc(-c3ccccc3)ccc21}	&	Cz2ICz	&	\smiles{c1ccc2c(c1)c1ccccc1n2-c1ccc2c(c1)c1cccc3c4cc(-n5c6ccccc6c6ccccc65)ccc4n2c13}	\\
TPA-2	&	\smiles{O=C1c2ccccc2C(=O)c2c1ccc1nc(-c3ccc(N(c4ccccc4)c4ccccc4)cc3)c(-c3ccc(N(c4ccccc4)c4ccccc4)cc3)nc21}	&	Ac-BPCN	&	\smiles{CC1(C)c2ccccc2N(c2ccc3nc4c5cc(C\#N)ccc5c5ccc(C\#N)cc5c4nc3c2)c2ccccc21}	\\
36DTPAFM	&	\smiles{N\#CC(C\#N)=C1c2ccc(-c3ccc(N(c4ccccc4)c4ccccc4)cc3)cc2-c2cc(-c3ccc(N(c4ccccc4)c4ccccc4)cc3)ccc21}	&	5CzICz	&	\smiles{c1ccc2c(c1)c1ccccc1n2-c1ccc2c(c1)c1cccc3c4ccccc4n2c31}	\\
AQTC	&	\smiles{N\#Cc1cc2nc3c(nc2c(C\#N)c1C\#N)-c1cccc2cccc-3c12}	&	DCPDAPM	&	\smiles{CC1(C)c2ccccc2N(c2ccc(C(=O)c3cc(-n4c5ccccc5c5ccccc54)cc(-n4c5ccccc5c5ccccc54)c3)cc2)c2ccccc21}	\\
3TPAFM	&	\smiles{N\#CC(C\#N)=C1c2ccccc2-c2cc(-c3ccc(N(c4ccccc4)c4ccccc4)cc3)ccc21}	&	SPFS-PXZ	&	\smiles{O=S1(=O)c2ccc(N3c4ccccc4Oc4ccccc43)cc2C2(c3ccccc3-c3ccccc32)c2cc(N3c4ccccc4Oc4ccccc43)ccc21}	\\
BBPA	&	\smiles{CCCCOc1ccccc1-c1c2ccccc2c(-c2ccccc2OCCCC)c2ccccc12}	&	DPS-PXZ	&	\smiles{O=S(=O)(c1ccc(N2c3ccccc3Oc3ccccc32)cc1)c1ccc(N2c3ccccc3Oc3ccccc32)cc1}	\\
DMPA	&	\smiles{COc1ccc(-c2c3ccccc3c(-c3ccc(OC)cc3OC)c3ccccc23)c(OC)c1}	&	TX-CzPh	&	\smiles{O=c1c2ccccc2sc2ccc(-c3ccc4c(c3)c3ccccc3n4-c3ccccc3)cc12}	\\
PXZ-NAI	&	\smiles{O=C1c2cccc3c(N4c5ccccc5Oc5ccccc54)ccc(c23)C(=O)N1c1ccccc1}	&	bpipH	&	\smiles{Brc1cccc(-c2nc3c4cccnc4c4ncccc4c3[nH]2)n1}	\\
PTZ-NAI	&	\smiles{O=C1c2cccc3c(N4c5ccccc5Sc5ccccc54)ccc(c23)C(=O)N1c1ccccc1}	&	Tri-PXZ-PCN	&	\smiles{N\#Cc1c(-c2ccc(N3c4ccccc4Oc4ccccc43)cc2)nc(-c2ccc(N3c4ccccc4Oc4ccccc43)cc2)c(C\#N)c1-c1ccc(N2c3ccccc3Oc3ccccc32)cc1}	\\
AI-4Cz	&	\smiles{O=C1c2c(c(-n3c4ccccc4c4ccccc43)c(-n3c4ccccc4c4ccccc43)c(-n3c4ccccc4c4ccccc43)c2-n2c3ccccc3c3ccccc32)C(=O)N1c1ccccc1}	&	ACID–BPPZ	&	\smiles{CC1(C)c2ccccc2N(c2ccc3c(c2)c2ccccc2c2nc4c5cccnc5c5ncccc5c4nc32)c2cc3c4ccccc4n(-c4ccccc4)c3cc21}	\\
DAcIPN	&	\smiles{CC1(C)c2ccccc2N(c2cc(N3c4ccccc4C(C)(C)c4ccccc43)c(C\#N)cc2C\#N)c2ccccc21}	&	C-BPyIA	&	\smiles{c1ccc(-n2c(-c3ccc(-c4c5ccccc5c(-c5ccc(-c6nc7c8cccc9ccc\%10cccc(c\%10c98)c7n6-c6ccccc6)cc5)c5ccccc45)cc3)nc3c4cccc5ccc6cccc(c6c54)c32)cc1}	\\
TCz1	&	\smiles{CCCCC(CC)Cn1c2ccc(-n3c4ccccc4c4ccccc43)cc2c2cc(-n3c4ccccc4c4ccccc43)ccc21}	&	N-BPyIA	&	\smiles{c1ccc(-c2nc3c4cccc5ccc6cccc(c6c54)c3n2-c2ccc(-c3c4ccccc4c(-c4ccc(-n5c(-c6ccccc6)nc6c7cccc8ccc9cccc(c9c87)c65)cc4)c4ccccc34)cc2)cc1}	\\
MBF	&	\smiles{F[B-]1(F)OC(c2ccccc2)=Cc2cc(-c3ccc(N4c5ccccc5Oc5ccccc54)cc3)nc[n+]21}	&	2-TA-AN	&	\smiles{c1ccc2c(c1)Sc1ccc(-c3c4ccccc4cc4ccccc34)cc1S2}	\\
MPA	&	\smiles{O=C(Cc1cc(-c2ccc(N3c4ccccc4Oc4ccccc43)cc2)ncn1)c1ccccc1}	&	ANOP-MeFAC	&	\smiles{Cc1ccc2c(c1)C1(c3ccccc3-c3ccccc31)c1cc(C)ccc1N2c1ccc2c3c(cccc13)-c1nc3nonc3nc1-2}	\\
3,3-NHC	&	\smiles{O=c1oc2ccccc2c(O)c1C(c1ccc2ccccc2c1)c1c(O)c2ccccc2oc1=O}	&	ANOP-DMAC	&	\smiles{CC1(C)c2ccccc2N(c2ccc3c4c(cccc24)-c2nc4nonc4nc2-3)c2ccccc21}	\\
3CP-DPS-DMAC	&	\smiles{CC1(C)c2ccccc2N(c2ccc(S(=O)(=O)c3ccc4c(c3)c3ccccc3n4-c3ccccc3)cc2)c2ccccc21}	&	4PyCNICz	&	\smiles{N\#Cc1c(-c2ccccc2)nc(-c2ccccc2)nc1-c1ccc(-n2c3ccccc3c3ccc4c(c5ccccc5n4-c4ccccc4)c32)cc1}	\\
3CP-DPS-PXZ	&	\smiles{O=S(=O)(c1ccc(N2c3ccccc3Oc3ccccc32)cc1)c1ccc2c(c1)c1ccccc1n2-c1ccccc1}	&	2PyCNBCz	&	\smiles{N\#Cc1c(-c2ccccc2)nc(-c2ccc(-n3c4ccccc4c4cc(-c5ccc6c(c5)c5ccccc5n6-c5ccccc5)ccc43)cc2)nc1-c1ccccc1}	\\
3CzAnN	&	\smiles{c1ccc(-n2c3ccccc3c3cc(-c4c5ccccc5c(-c5ccc6ccccc6c5)c5ccccc45)ccc32)cc1}	&	CzA-PPO	&	\smiles{Cc1ccc(N(c2ccc(C)cc2)c2c3ccccc3c(-c3ccc(P(=O)(c4ccccc4)c4ccc(-n5c6ccccc6c6ccccc65)cc4)cc3)c3ccccc23)cc1}	\\
DPQ-DMAC	&	\smiles{CC1(C)c2ccccc2N(c2ccc3nc(-c4ccccn4)c(-c4ccccn4)nc3c2)c2ccccc21}	&	2PyCNICz	&	\smiles{N\#Cc1c(-c2ccccc2)nc(-c2ccc(-n3c4ccccc4c4ccc5c(c6ccccc6n5-c5ccccc5)c43)cc2)nc1-c1ccccc1}	\\
DPAA	&	\smiles{c1ccc(-c2c3ccccc3c(-c3ccccc3)c3ccccc23)cc1}	&	TAA-PPO	&	\smiles{Cc1ccc(N(c2ccc(C)cc2)c2c3ccccc3c(-c3ccc(P(=O)(c4ccccc4)c4ccccc4)cc3)c3ccccc23)cc1}	\\
t3Cz-SO	&	\smiles{CC(C)(C)c1ccc2c(c1)c1cc(C(C)(C)C)ccc1n2-c1ccc2c(c1)c1cc(-n3c4ccc(C(C)(C)C)cc4c4cc(C(C)(C)C)ccc43)ccc1n2-c1ccc(S(=O)(=O)c2ccccc2)cc1}	&	pSFIAc2	&	\smiles{Cc1ccc2c(c1)c1c3c4c(c5c1n2-c1ccc(C(C)(C)C)cc1C51c2ccccc2-c2ccccc21)c1cc(C)ccc1n4-c1ccc(C(C)(C)C)cc1C31c2ccccc2-c2ccccc21}	\\
Cz-SO	&	\smiles{O=S(=O)(c1ccccc1)c1ccc(-n2c3ccccc3c3ccccc32)cc1}	&	TBN-TPA	&	\smiles{CC(C)(C)c1ccc(N2c3ccc(C(C)(C)C)cc3B3c4cc(C(C)(C)C)ccc4N(c4ccc(C(C)(C)C)cc4)c4cc(-n5c6ccc(C(C)(C)C)cc6c6cc(C(C)(C)C)ccc65)cc2c43)cc1}	\\
DPQ-PXZ	&	\smiles{c1ccc(-c2nc3ccc(N4c5ccccc5Oc5ccccc54)cc3nc2-c2ccccn2)nc1}	&	CzBPDCb	&	\smiles{c1cc(-c2cccc(-n3c4ccccc4c4ncccc43)c2)cc(-n2c3ccccc3c3ccccc32)c1}	\\
pCzAnN	&	\smiles{c1ccc2cc(-c3c4ccccc4c(-c4ccc(-n5c6ccccc6c6ccccc65)cc4)c4ccccc34)ccc2c1}	&	{[2,1-b]IF}	&	\smiles{C1=CC2(C=C3C1=CC1=C3C3(C4=c5ccccc5=CC4=C1)c1ccccc1Sc1ccccc13)c1ccccc1Sc1ccccc12}	\\
3-BHH	&	\smiles{O=c1oc2ccccc2c(O)c1C(c1ccc(Br)s1)c1c(O)c2ccccc2oc1=O}	&	cyclohexanone	&	\smiles{BrCCCCCCC1(CCCCCCBr)c2cc(Br)ccc2-c2ccc(Br)cc21}	\\
DPP-PXZ	&	\smiles{c1ccc(-c2nc3cc(N4c5ccccc5Oc5ccccc54)cnc3nc2-c2ccccn2)nc1}	&	thio-ether	&	\smiles{C1=CC2(C=C3C1=CC1=C3C3(C4=c5ccccc5=CC4=C1)c1ccccc1Sc1ccccc13)c1ccccc1Sc1ccccc12}	\\
APPT-PXZ	&	\smiles{c1ccc2c(c1)Oc1ccccc1N2c1ccc2c3c(cccc13)-c1nc3c4cccnc4c4ncccc4c3nc1-2}	&	ICz-2	&	\smiles{c1ccc(-c2cccnc2-c2cc(-c3c[nH]nn3)cc(-c3ccccn3)n2)cc1}	\\
HITPE	&	\smiles{Oc1ccc(-c2ccc(C(=C(c3ccccc3)c3ccccc3)c3ccccc3)cc2)cc1-c1nc(-c2ccccc2)c(-c2ccccc2)n1-c1ccccc1}	&	DDCzTrz	&	\smiles{c1ccc(-c2nc(-c3cc(-n4c5ccccc5c5ccccc54)cc(-n4c5ccccc5c5ccccc54)c3)nc(-c3cc(-n4c5ccccc5c5ccccc54)cc(-n4c5ccccc5c5ccccc54)c3)n2)cc1}	\\
DV-3CzCN	&	\smiles{C=Cc1ccc(COc2cccc3c2c2ccccc2n3-c2cc(-n3c4ccccc4c4c(OCc5ccc(C=C)cc5)cccc43)c(C\#N)c(-n3c4ccccc4c4c(OCc5ccc(C=C)cc5)cccc43)c2)cc1}	&	B-O-dpa	&	\smiles{CC(C)(C)c1cc2c3c(c1)N(c1ccccc1)c1ccccc1B3c1ccccc1O2}	\\
SF-DPSO	&	\smiles{O=S(=O)(c1ccc(-c2ccc3c(c2)C2(c4ccccc4-c4ccccc42)c2ccccc2-3)cc1)c1ccc(-c2ccc3c(c2)C2(c4ccccc4-c4ccccc42)c2ccccc2-3)cc1}	&	B-dpa-SpiroAC	&	\smiles{CC(C)(C)c1cc2c3c(c1)N1c4ccccc4C4(c5ccccc5-c5ccccc54)c4cccc(c41)B3c1ccccc1N2c1ccccc1}	\\
BDQ-tBuCz	&	\smiles{CC(C)(C)c1ccc2c(c1)c1cc(C(C)(C)C)ccc1n2-c1ccc(-n2c3ccc(C(C)(C)C)cc3c3cc(C(C)(C)C)ccc32)c2nc(-c3ccccc3)c(-c3ccccc3)nc12}	&	DCz-OTP	&	\smiles{c1cc(-c2ccccc2-c2cccc(-n3c4ccccc4c4ccccc43)c2)cc(-n2c3ccccc3c3ccccc32)c1}	\\
BDQDMAC	&	\smiles{CC1(C)c2ccccc2N(c2ccc(N3c4ccccc4C(C)(C)c4ccccc43)c3nc(-c4ccccc4)c(-c4ccccc4)nc23)c2ccccc21}	&	B-dpa-Cz	&	\smiles{CC(C)(C)c1cc2c3c(c1)-n1c4ccccc4c4cccc(c41)B3c1ccccc1N2c1ccccc1}	\\
BDQPCz	&	\smiles{c1ccc(-c2nc3c(-c4ccc(-n5c6ccccc6c6ccccc65)cc4)ccc(-c4ccc(-n5c6ccccc6c6ccccc65)cc4)c3nc2-c2ccccc2)cc1}	&	c-D2T	&	\smiles{CC(C)(C)c1ccc2c(c1)c1cc(C(C)(C)C)ccc1n2-c1ccc(N2c3ccccc3C3(c4ccccc4-c4cc5c(cc43)C3(c4ccccc4-5)c4ccccc4N(c4ccc(-c5nc(-c6ccccc6)nc(C6CCCCC6)n5)cc4)c4ccccc43)c3ccccc32)cc1}	\\
4-SCZ	&	\smiles{c1ccc2c(c1)-c1ccccc1C21c2ccccc2-c2c1ccc1[nH]c3ccccc3c21}	&	DPXZ-DPPM	&	\smiles{O=C(c1ccccc1)c1ccc2nc3c4ccc(N5c6ccccc6Oc6ccccc65)cc4c4cc(N5c6ccccc6Oc6ccccc65)ccc4c3nc2c1}	\\
Py-TPA	&	\smiles{CC(C)(C)c1ccc(N(c2ccc(-c3cc4nc5c6ccc(-c7ccncc7)nc6c6nc(-c7ccncc7)ccc6c5nc4cc3C\#N)cc2)c2ccc(C(C)(C)C)cc2)cc1}	&	CBPMCN	&	\smiles{N\#Cc1ccc(-n2c(-c3ccc(-n4c5ccccc5c5ccccc54)cc3)nc3c4ccccc4c4ccccc4c32)cc1}	\\
DPACPhTPI	&	\smiles{CC(C)(C)c1ccc(-n2c(-c3ccc(-c4ccc(N5c6ccccc6C(c6ccccc6)(c6ccccc6)c6ccccc65)cc4)cc3)nc3c4ccccc4c4ccccc4c32)cc1}	&	CPBPMCN	&	\smiles{N\#Cc1ccc(-n2c(-c3ccc(-c4ccc(-n5c6ccccc6c6ccccc65)cc4)cc3)nc3c4ccccc4c4ccccc4c32)cc1}	\\
oPTBC	&	\smiles{N\#Cc1ccc(-c2cc(C\#N)cc(-c3ccc(C\#N)cc3)c2N2c3ccccc3Oc3ccccc32)cc1}	&	TPB-AC	&	\smiles{N\#Cc1ccc(-c2cc(-c3ccccc3)c(-c3ccc(N(c4ccccc4)c4ccccc4)cc3)cc2-c2ccccc2)cc1}	\\
mPTBC	&	\smiles{N\#Cc1ccc(-c2cc(N3c4ccccc4Oc4ccccc43)cc(-c3ccc(C\#N)cc3)c2C\#N)cc1}	&	Me-DMAC	&	\smiles{Cc1nc(C)c(C\#N)c(-c2ccc(N3c4ccccc4C(C)(C)c4ccccc43)cc2)c1C\#N}	\\
P3C	&	\smiles{CCCCC(CC)Cn1c2cc(C\#N)ccc2c2ccc(C\#Cc3ccc4ccc5cccc6ccc3c4c56)cc21}	&	iTPBI-CN	&	\smiles{N\#Cc1c(-n2c(-c3ccccc3)nc3ccccc32)cc(-n2c(-c3ccccc3)nc3ccccc32)cc1-n1c(-c2ccccc2)nc2ccccc21}	\\
DBCP	&	\smiles{c1ccc2c(c1)Oc1ccccc1N2c1cc2c3ccccc3c3ccc4ccc5c6ccccc6c1c1c5c4c3c21}	&	TCP	&	\smiles{c1ccc2c(c1)c1ccccc1n2-c1cc(-n2c3ccccc3c3ccccc32)cc(-n2c3ccccc3c3ccccc32)c1}	\\
NAQ3	&	\smiles{O=C1c2ccccc2C(=O)c2cc(N(c3ccc4c(c3)C(=O)c3ccccc3C4=O)c3ccc4c(c3)C(=O)c3ccccc3C4=O)ccc21}	&	CDBP-BP-PXZ	&	\smiles{Cc1cc(-n2c3ccccc3c3ccccc32)ccc1-c1ccc(-n2c3ccccc3c3cc(C(=O)c4ccc(N5c6ccccc6Oc6ccccc65)cc4)ccc32)cc1C}	\\
mCP-BP-DMAC	&	\smiles{O=C(c1ccc(N2c3ccccc3Oc3ccccc32)cc1)c1ccc2c(c1)c1ccccc1n2-c1cccc(-n2c3ccccc3c3ccccc32)c1}	&	CDBP-BP-DMAC	&	\smiles{Cc1cc(-n2c3ccccc3c3ccccc32)ccc1-c1ccc(-n2c3ccccc3c3cc(C(=O)c4ccc(N5c6ccccc6C(C)(C)c6ccccc65)cc4)ccc32)cc1C}	\\
BFCZPZ3	&	\smiles{c1ccc2c(c1)oc1ccc3c4ccccc4n(-c4cnc(-n5c6ccccc6c6ccc7oc8ccccc8c7c65)cn4)c3c12}	&	3CzPhpPM	&	\smiles{c1ccc(-c2ccc3c(c2)c2ccccc2n3-c2ccc(-c3cc(-c4ccccc4)nc(-c4ccccc4)n3)cc2)cc1}	\\
BTCZPZ1	&	\smiles{c1ccc2c(c1)sc1c2ccc2c3ccccc3n(-c3cnc(-n4c5ccccc5c5ccc6c7ccccc7sc6c54)cn3)c21}	&	SFI23pPM	&	\smiles{c1ccc(-c2cc(-c3ccc(-n4c5ccccc5c5cc6c(cc54)C4(c5ccccc5-c5ccccc54)c4ccccc4-6)cc3)nc(-c3ccccc3)n2)cc1}	\\
BFCZPZ2	&	\smiles{c1ccc2c(c1)oc1c2ccc2c1c1ccccc1n2-c1cnc(-n2c3ccccc3c3c4oc5ccccc5c4ccc32)cn1}	&	3PCzPFP	&	\smiles{c1ccc(-n2c3ccccc3c3cc(-c4cnc5oc6ncccc6c5c4)ccc32)cc1}	\\
CZ9CZPZ	&	\smiles{c1ccc2c(c1)c1ccccc1n2-c1ccc2c(c1)c1ccccc1n2-c1cnc(-n2c3ccccc3c3cc(-n4c5ccccc5c5ccccc54)ccc32)cn1}	&	SFI34pPM	&	\smiles{c1ccc(-c2cc(-c3ccc(-n4c5ccccc5c5c6c(ccc54)-c4ccccc4C64c5ccccc5-c5ccccc54)cc3)nc(-c3ccccc3)n2)cc1}	\\
BFCZPZ1	&	\smiles{c1ccc2c(c1)oc1c2ccc2c3ccccc3n(-c3cnc(-n4c5ccccc5c5ccc6c7ccccc7oc6c54)cn3)c21}	&	SBA-2DPS	&	\smiles{O=S(=O)(c1ccccc1)c1ccc(N2c3ccccc3C3(c4ccccc42)c2ccccc2N(c2ccc(S(=O)(=O)c4ccccc4)cc2)c2ccccc23)cc1}	\\
TPBPPI-PY	&	\smiles{c1ccc(-c2cc(-c3ccc(-c4nc5c6ccccc6c6ccccc6c5n4-c4ccc(-c5cccnc5)cc4)cc3)c(-c3ccccc3)c(-c3ccccc3)c2-c2ccccc2)cc1}	&	2DMAC-BP-F	&	\smiles{CC1(C)c2ccccc2N(c2ccc3c(c2)c2cc(N4c5ccccc5C(C)(C)c5ccccc54)ccc2c2nc4c(F)cccc4nc32)c2ccccc21}	\\
TPBPPI-PBI	&	\smiles{c1ccc(-c2cc(-c3ccc(-c4nc5c6ccccc6c6ccccc6c5n4-c4ccc(-c5ccc(-c6nc7ccccc7n6-c6ccccc6)cc5)cc4)cc3)c(-c3ccccc3)c(-c3ccccc3)c2-c2ccccc2)cc1}	&	DCA	&	\smiles{N\#Cc1cc(C\#N)cc(-c2c3ccccc3c(-c3cc(C\#N)cc(C\#N)c3)c3ccccc23)c1}	\\
PTZ-2PTO	&	\smiles{O=S1(=O)c2ccccc2N(c2ccccc2)c2cc(N3c4ccccc4Sc4ccccc43)ccc21}	&	DPFPhe	&	\smiles{CCCCCCN1C2=CCC(=CC3c4cc(-c5ccccc5)ccc4-c4ccc(-c5ccccc5)cc43)C=C2Oc2ccccc21}	\\
4F-m-$\nu$-DABNA	&	\smiles{Cc1cccc(N(c2cccc(C)c2)c2cc3c4c(c2)N(c2c(F)cccc2F)c2cc5c(cc2B4c2ccc(C)cc2N3c2cccc(C)c2)B2c3ccc(C)cc3N(c3cccc(C)c3)c3cc(N(c4cccc(C)c4)c4cccc(C)c4)cc(c32)N5c2c(F)cccc2F)c1}	&	DFPFPhe	&	\smiles{CCCCCCN1C2=CCC(=CC3c4cc(-c5ccc(F)cc5)ccc4-c4ccc(-c5ccc(F)cc5)cc43)C=C2Oc2ccccc21}	\\
m-$\nu$-DABNA	&	\smiles{Cc1cccc(N(c2cccc(C)c2)c2cc3c4c(c2)N(c2ccccc2)c2cc5c(cc2B4c2ccc(C)cc2N3c2cccc(C)c2)B2c3ccc(C)cc3N(c3cccc(C)c3)c3cc(N(c4cccc(C)c4)c4cccc(C)c4)cc(c32)N5c2ccccc2)c1}	&	3ai	&	\smiles{CCCCc1cc2c(s1)-c1sc(CCCC)cc1P2(=O)c1ccc(-n2c3ccccc3c3ccccc32)cc1}	\\
TPA-PyI	&	\smiles{Brc1cccc(-n2c(-c3ccc(-c4ccc(N(c5ccccc5)c5ccccc5)cc4)cc3)nc(-c3ccccc3)c2-c2ccccc2)c1}	&	DNFPh	&	\smiles{CCCCCCN1C2=CCC(=CC3c4cc(-c5cccc6ccccc56)ccc4-c4ccc(-c5cccc6ccccc56)cc43)C=C2Sc2ccccc21}	\\
4F-$\nu$-DABNA	&	\smiles{Fc1cccc(F)c1N1c2cc3c(cc2B2c4ccccc4N(c4ccccc4)c4cc(N(c5ccccc5)c5ccccc5)cc1c42)B1c2ccccc2N(c2ccccc2)c2cc(N(c4ccccc4)c4ccccc4)cc(c21)N3c1c(F)cccc1F}	&	pCzAnBzt	&	\smiles{N\#Cc1ccc(-c2c3ccccc3c(-c3ccc(-n4c5ccccc5c5ccccc54)cc3)c3ccccc23)cc1}	\\
CBPTZ	&	\smiles{c1ccc(-c2nc(-c3ccccc3)nc(-c3cccc(-c4cc(-n5c6ccccc6c6ccccc65)ccc4-c4ccc(-n5c6ccccc6c6ccccc65)cc4)c3)n2)cc1}	&	DPE-DPXZ	&	\smiles{C(\#Cc1ccc(N2c3ccccc3Oc3ccccc32)cc1)c1ccc(N2c3ccccc3Oc3ccccc32)cc1}	\\
CBPBM	&	\smiles{c1ccc(-n2c(-c3ccc(-c4cc(-n5c6ccccc6c6ccccc65)ccc4-c4ccc(-n5c6ccccc6c6ccccc65)cc4)cc3)nc3ccccc32)cc1}	&	3CzAnBzt	&	\smiles{N\#Cc1ccc(-c2c3ccccc3c(-c3ccc4c(c3)c3ccccc3n4-c3ccccc3)c3ccccc23)cc1}	\\
CBPPY	&	\smiles{c1cncc(-c2cc(-n3c4ccccc4c4ccccc43)ccc2-c2ccc(-n3c4ccccc4c4ccccc43)cc2)c1}	&	TABD-COOH	&	\smiles{O=C(O)c1ccc(/C(=C$\backslash$C=C($\backslash$c2ccccc2)c2ccc(C(=O)O)cc2)c2ccccc2)cc1}	\\
N-B	&	\smiles{Cc1cc(C)c(B(c2ccccc2)c2c(C)cc(C)cc2C)c(C)c1}	&	ADN	&	\smiles{c1ccc2cc(-c3c4ccccc4c(-c4ccc5ccccc5c4)c4ccccc34)ccc2c1}	\\
DPACpBr	&	\smiles{Cc1cc(N2c3ccccc3C(c3ccccc3)(c3ccccc3)c3ccccc32)cc(C)c1Br}	&	Et3SiH	&	\smiles{FC1=C(F)C(c2c(F)c(F)c(F)c(F)c2F)(c2c(F)c(F)c(F)c(F)c2F)C(c2c(F)c(F)c(F)c(F)c2F)(c2c(F)c(F)c(F)c(F)c2F)B1c1c(F)c(F)c(F)c(F)c1F}	\\
DPCN	&	\smiles{CC1(C)c2ccccc2N(c2ccc(-c3ccc(-c4nc5c6ccccc6c6ccccc6c5n4-c4ccc(C\#N)cc4)cc3)cc2)c2ccccc21}	&	DPE-DDMAc	&	\smiles{CC1(C)c2ccccc2N(c2ccc(C\#Cc3ccc(N4c5ccccc5C(C)(C)c5ccccc54)cc3)cc2)c2ccccc21}	\\
BZC-PXZ	&	\smiles{O=C(CC(=O)c1ccccc1O)c1ccc(N2c3ccccc3Oc3ccccc32)cc1}	&	mCz-TAn-CN	&	\smiles{CC(C)(C)c1ccc2c(-c3cccc(-n4c5ccccc5c5ccccc54)c3)c3cc(C(C)(C)C)ccc3c(-c3ccc(C\#N)cc3)c2c1}	\\
CzBN	&	\smiles{N\#Cc1c(-n2c3ccccc3c3ccccc32)cccc1-n1c2ccccc2c2ccccc21}	&	indenoacridine	&	\smiles{c1ccc2[nH]c3ccc4c5ccccc5cc4c3cc2c1}	\\
2CzBN	&	\smiles{CC(C)(C)c1ccc2c(c1)c1cc(C(C)(C)C)ccc1n2-c1cccc(-n2c3ccc(C(C)(C)C)cc3c3cc(C(C)(C)C)ccc32)c1C\#N}	&	MPAc	&	\smiles{Cc1ccc2c(c1)C(c1ccccc1)(c1ccccc1)c1cc(C)ccc1N2}	\\
TPA-DPPI	&	\smiles{c1ccc(-c2nc3c4ccccc4c4ccccc4c3n2-c2ccc(-n3c(-c4ccc(-c5ccc(N(c6ccccc6)c6ccccc6)cc5)cc4)nc4c5ccccc5c5ccccc5c43)cc2)cc1}	&	TBAc	&	\smiles{CC(C)(C)C1=CC=CC(C(C)(C)C)(C2(C3(C(C)(C)C)C=CC=C(C(C)(C)C)C3)c3ccccc3N(c3ccccc3)c3ccccc32)C1}	\\
ANQDC-PSTA	&	\smiles{N\#Cc1cc2nc3c(nc2cc1C\#N)-c1ccc(N2c4ccc(-c5ccccc5)cc4C4(c5ccccc5-c5ccccc54)c4ccc5c(sc6ccccc65)c42)c2cccc-3c12}	&	poCz-SO	&	\smiles{O=P(c1ccccc1)(c1ccccc1)c1ccc2[nH]c3ccc(P(=O)(c4ccccc4)c4ccccc4)cc3c2c1}	\\
pDPAPQ-PXZ	&	\smiles{c1ccc2c(c1)Oc1ccccc1N2c1ccc2c3c(cccc13)-c1nc3cc(-c4ccncc4)c(-c4ccncc4)cc3nc1-2}	&	H4TCPE	&	\smiles{O=C(O)c1ccc(CC(c2ccc(C(=O)O)cc2)(c2ccc(C(=O)O)cc2)c2ccc(C(=O)O)cc2)cc1}	\\
m-CZ-DPS-DMAC	&	\smiles{CC1(C)c2ccccc2N(c2cccc(S(=O)(=O)c3cc(-n4c5ccccc5c5ccccc54)cc(-n4c5ccccc5c5ccccc54)c3)c2)c2ccccc21}	&	CbBPCb	&	\smiles{c1cc(-c2cccc(-n3c4ccccc4c4cccnc43)c2)cc(-n2c3ccccc3c3cccnc32)c1}	\\
pDPBPZ-PXZ	&	\smiles{c1ccc2c(c1)Oc1ccccc1N2c1ccc2c(c1)c1ccccc1c1nc3cc(-c4ccncc4)c(-c4ccncc4)cc3nc21}	&	TMOE	&	\smiles{COc1ccc(C(=C(c2ccc(OC)cc2)c2ccc(OC)cc2)c2ccc(OC)cc2)cc1}	\\
DDPhCz-DPPZ	&	\smiles{c1ccc(-c2ccc3c(c2)c2cc(-c4ccccc4)ccc2n3-c2cc3nc4c5cccnc5c5ncccc5c4nc3cc2-n2c3ccc(-c4ccccc4)cc3c3cc(-c4ccccc4)ccc32)cc1}	&	FAPbCl3	&	\smiles{c1ccc(-c2ccc(-c3ccccc3)c(-c3ccccc3)c2-c2ccccc2)cc1}	\\
DCz-DPPZ	&	\smiles{c1cnc2c(c1)c1nc3cc(-n4c5ccccc5c5ccccc54)c(-n4c5ccccc5c5ccccc54)cc3nc1c1cccnc12}	&	DBT-SADF	&	\smiles{O=S1(=O)c2ccccc2-c2cc(N3c4ccccc4C4(c5ccccc5-c5ccccc54)c4ccccc43)ccc21}	\\
Oxd-bCZ	&	\smiles{c1ccc2c(c1)c1ccccc1n2-c1ccc(-c2ccc(-c3nnc(-c4ccc(-c5ccc(-n6c7ccccc7c7ccccc76)cc5)cc4)o3)cc2)cc1}	&	mPXTMP	&	\smiles{Cc1c(-c2ccc3c(c2)Oc2ccccc2S3)cccc1-c1ccc2c(c1)Oc1ccccc1S2}	\\
Fene	&	\smiles{c1ccc2c(c1)Oc1ccccc1N2c1ccc2cccc(N3c4ccccc4Oc4ccccc43)c2n1}	&	APDC-tPh	&	\smiles{N\#Cc1nc2c(nc1C\#N)-c1ccc(-c3cc(-c4ccccc4)cc(-c4ccccc4)c3)c3cccc-2c13}	\\
Fens	&	\smiles{c1ccc2c(c1)Sc1ccccc1N2c1ccc2cccc(N3c4ccccc4Sc4ccccc43)c2n1}	&	POZDBPHZ	&	\smiles{c1ccc2c(c1)Oc1ccccc1N2c1ccc2c(ccc3nc4ccc5cc(N6c7ccccc7Oc7ccccc76)ccc5c4nc32)c1}	\\
Yad	&	\smiles{CC1(C)c2ccccc2N(c2ccc3cccc(N4c5ccccc5C(C)(C)c5ccccc54)c3n2)c2ccccc21}	&	CPTBF	&	\smiles{CC(C)(C)c1ccc2c(c1)C(c1ccc(-n3c4ccccc4c4ccccc43)cc1)(c1ccc(-n3c4ccccc4c4ccccc43)cc1)c1cc(C(C)(C)C)ccc1-2}	\\
DTPCZPYTZ	&	\smiles{CC(C)(C)c1ccc2c(c1)c1cc(C(C)(C)C)ccc1n2-c1ccc(-c2ccc3ccc4c(-c5ccc(-c6nnc(-c7ccccc7)n6-c6ccccc6)cc5)ccc5ccc2c3c54)cc1}	&	trimesitylboron	&	\smiles{Cc1cc(C)c(B(c2c(C)cc(C)cc2C)c2c(C)cc(C)cc2C)c(C)c1}	\\
TPADSO2	&	\smiles{O=S1(=O)c2ccccc2S(=O)(=O)c2cc(-c3ccc(N(c4ccccc4)c4ccccc4)cc3)ccc21}	&	CMP	&	\smiles{c1ccc(C(c2ccccc2)(c2ccccc2)c2ccccc2)cc1}	\\
B4PyPPM	&	\smiles{c1ccc(-c2nc(-c3cc(-c4ccncc4)cc(-c4ccncc4)c3)cc(-c3cc(-c4ccncc4)cc(-c4ccncc4)c3)n2)cc1}	&	CH3OH–CHCl3	&	\smiles{O=P(c1ccccc1)(c1ccccc1)c1cccc(-c2ccc3ccc4cccc5ccc2c3c45)c1}	\\
mAP	&	\smiles{CC1(C)c2ccccc2N(c2cccc(N3c4ccccc4C(C)(C)c4ccccc43)c2)c2ccccc21}	&	s-triazine	&	\smiles{CC(C)(C)c1ccc2c(c1)c1cc(C(C)(C)C)ccc1n2-c1ccc([Si](c2ccccc2)(c2ccccc2)c2ccc(B(O)O)cc2)cc1}	\\
DBBPZ-DPXZ	&	\smiles{c1ccc(-c2cc3nc4c5ccc(N6c7ccccc7Oc7ccccc76)cc5c5cc(N6c7ccccc7Oc7ccccc76)ccc5c4nc3cc2-c2ccccc2)cc1}	&	{[BMIM][PF6]}	&	\smiles{CCCCCCC1(CCCCCC)c2ccccc2-c2ccc(Br)cc21}	\\
DPIBPZ-DPXZ	&	\smiles{c1ccc2c(c1)Oc1ccccc1N2c1ccc2c(c1)c1cc(N3c4ccccc4Oc4ccccc43)ccc1c1nc3cc(-c4cncnc4)c(-c4cncnc4)cc3nc21}	&	1SBFN	&	\smiles{c1ccc2c(c1)-c1ccccc1C21c2ccccc2-c2ccc(-c3cccc4ccccc34)cc21}	\\
m-DTPACO	&	\smiles{O=C(c1ccc(N(c2ccccc2)c2ccccc2)cc1)c1cccc(C(=O)c2ccc(N(c3ccccc3)c3ccccc3)cc2)c1}	&	2SBFN	&	\smiles{c1ccc2c(c1)-c1ccccc1C21c2ccccc2-c2ccc(-c3ccc4ccccc4c3)cc21}	\\
mDPBPZ-DPXZ	&	\smiles{c1cncc(-c2cc3nc4c5ccc(N6c7ccccc7Oc7ccccc76)cc5c5cc(N6c7ccccc7Oc7ccccc76)ccc5c4nc3cc2-c2cccnc2)c1}	&	F2PA	&	\smiles{CC1(C)c2ccccc2-c2ccc(N(c3ccccc3)c3ccc4c(c3)C(C)(C)c3ccccc3-4)cc21}	\\
p-DTPACO	&	\smiles{O=C(c1ccc(C(=O)c2ccc(N(c3ccccc3)c3ccccc3)cc2)cc1)c1ccc(N(c2ccccc2)c2ccccc2)cc1}	&	TSBFB	&	\smiles{c1ccc2c(c1)-c1ccccc1C21c2ccccc2-c2ccc(-c3cc(-c4ccc5c(c4)C4(c6ccccc6-c6ccccc64)c4ccccc4-5)cc(-c4ccc5c(c4)C4(c6ccccc6-c6ccccc64)c4ccccc4-5)c3)cc21}	\\
2TDBA-SBA	&	\smiles{CC(C)(C)c1ccc2c(c1)B1c3cc(C(C)(C)C)ccc3Oc3cc(N4c5ccccc5C5(c6ccccc64)c4ccccc4N(c4cc6c7c(c4)Oc4ccc(C(C)(C)C)cc4B7c4cc(C(C)(C)C)ccc4O6)c4ccccc45)cc(c31)O2}	&	DPPP	&	\smiles{c1ccc(-c2cc(-c3ccccc3)cc(-c3ccccc3-c3ccc4ccc5c(-c6ccccc6-c6cc(-c7ccccc7)cc(-c7ccccc7)c6)ccc6ccc3c4c65)c2)cc1}	\\
TPA-APm	&	\smiles{c1ccc(N(c2ccccc2)c2ccc(-c3cnc4nc5c(nc4n3)-c3cccc4cccc-5c34)cc2)cc1}	&	TPA-NZP	&	\smiles{c1ccc(-c2c3ccccc3c(-c3ccc(N(c4ccccc4)c4ccccc4)cc3)c3nsnc23)cc1}	\\
TPA-AP	&	\smiles{c1ccc(N(c2ccccc2)c2ccc(-c3ccc4nc5c(nc4c3)-c3cccc4cccc-5c34)cc2)cc1}	&	mTPA-PPI	&	\smiles{c1ccc(N(c2ccccc2)c2ccc(-c3cccc(-c4nc5c6ccccc6c6ccccc6c5n4-c4ccccc4)c3)cc2)cc1}	\\
TDBA-SBA	&	\smiles{CC(C)(C)c1ccc(N2c3ccccc3C3(c4ccccc42)c2ccccc2N(c2cc4c5c(c2)Oc2ccc(C(C)(C)C)cc2B5c2cc(C(C)(C)C)ccc2O4)c2ccccc23)cc1}	&	BSBFB	&	\smiles{c1cc(-c2ccc3c(c2)C2(c4ccccc4-c4ccccc42)c2ccccc2-3)cc(-c2ccc3c(c2)C2(c4ccccc4-c4ccccc42)c2ccccc2-3)c1}	\\
TPA-APy	&	\smiles{c1ccc(N(c2ccccc2)c2ccc(-c3cnc4nc5c(nc4c3)-c3cccc4cccc-5c34)cc2)cc1}	&	BN-DBC-Ph2	&	\smiles{C1=CB2C(=C3C(=C4C=C(c5ccccc5)C=CN4c4ccc(-c5ccccc5)cc43)c3ccccc32)C=C1}	\\
TPD4PA	&	\smiles{c1ccc(N(c2ccccc2)c2cc3c4c(c2)N(c2ccccc2)c2cc5c6c(c2B4c2ccccc2N3c2ccccc2)Oc2ccccc2B6c2ccccc2O5)cc1}	&	TCbTrz	&	\smiles{c1ccc(-c2nc(-c3ccccc3)nc(-c3cc(-n4c5ccccc5c5ncccc54)c(-n4c5ccccc5c5ncccc54)c(-n4c5ccccc5c5ncccc54)c3)n2)cc1}	\\
o-Ac2BBP	&	\smiles{CC1(C)c2ccccc2N(c2ccc(C(=O)c3ccccc3C(=O)c3ccc(N4c5ccccc5C(C)(C)c5ccccc54)cc3)cc2)c2ccccc21}	&	Xm-mCP	&	\smiles{O=c1c2ccccc2oc2ccc(-c3cc(-n4c5ccccc5c5ccccc54)cc(-n4c5ccccc5c5ccccc54)c3)cc12}	\\
TPA-PyF2	&	\smiles{Cc1ccc(N(c2ccc(C3(c4cccnc4)c4cc(C(C)(C)C)ccc4-c4ccc(C(C)(C)C)cc43)cc2)c2ccc(C3(c4cccnc4)c4cc(C(C)(C)C)ccc4-c4ccc(C(C)(C)C)cc43)cc2)cc1}	&	Xp-mCP	&	\smiles{CC1(C)OB(c2cc(-n3c4ccccc4c4ccccc43)cc(-n3c4ccccc4c4ccccc43)c2)OC1(C)C}	\\
o-Px2BBP	&	\smiles{O=C(c1ccc(N2c3ccccc3Oc3ccccc32)cc1)c1ccccc1C(=O)c1ccc(N2c3ccccc3Oc3ccccc32)cc1}	&	ethylacetate	&	\smiles{O=S(=O)(c1ccc(N2c3ccccc3C=Cc3ccccc32)cc1)c1ccc(N2c3ccccc3C=Cc3ccccc32)cc1}	\\
D-BP-PXZ	&	\smiles{O=C(c1ccc(N2c3ccccc3Oc3ccccc32)cc1)c1ccccc1-c1ccccc1C(=O)c1ccc(N2c3ccccc3Oc3ccccc32)cc1}	&	3DMACIPN	&	\smiles{CC1(C)c2ccccc2N(c2c(F)c(N3c4ccccc4C(C)(C)c4ccccc43)c(C\#N)c(N3c4ccccc4C(C)(C)c4ccccc43)c2C\#N)c2ccccc21}	\\
BTD-PXZ	&	\smiles{c1ccc2c(c1)Oc1ccccc1N2c1ccc2nsnc2c1}	&	ICZ	&	\smiles{CCCCCc1c2[nH]c3ccccc3c2cc2[nH]c3ccccc3c12}	\\
tBuCzDBA	&	\smiles{Cc1cc(-n2c3ccc(C(C)(C)C)cc3c3cc(C(C)(C)C)ccc32)cc(C)c1B1c2ccccc2B(c2c(C)cc(-n3c4ccc(C(C)(C)C)cc4c4cc(C(C)(C)C)ccc43)cc2C)c2ccccc21}	&	BDAVBi	&	\smiles{C(=Cc1ccc(N(c2ccccc2)c2ccccc2)cc1)c1ccc(-c2ccc(C=Cc3ccc(N(c4ccccc4)c4ccccc4)cc3)cc2)cc1}	\\
5-Cz-Sac	&	\smiles{CC(C)(C)N1C(=O)c2cc(-n3c4ccccc4c4ccccc43)ccc2S1(=O)=O}	&	DPVBi	&	\smiles{C(=C(c1ccccc1)c1ccccc1)c1ccc(-c2ccc(C=C(c3ccccc3)c3ccccc3)cc2)cc1}	\\
6-Cz-Sac	&	\smiles{CC(C)(C)N1C(=O)c2ccc(-n3c4ccccc4c4ccccc43)cc2S1(=O)=O}	&	CBTD2	&	\smiles{CC(C)(C)c1ccc2c(c1)c1cc(C(C)(C)C)ccc1n2-c1ccc(-c2ccc(-c3ccc(-n4c5ccc(C(C)(C)C)cc5c5cc(C(C)(C)C)ccc54)cc3)c3nsnc23)cc1}	\\
DiCz-Sac	&	\smiles{CC(C)(C)N1C(=O)c2cc(-n3c4ccccc4c4ccccc43)c(-n3c4ccccc4c4ccccc43)cc2S1(=O)=O}	&	MeOCzPCN	&	\smiles{COc1ccc2c(c1)c1cc(OC)ccc1n2-c1ncnc(-n2c3ccc(OC)cc3c3cc(OC)ccc32)c1C\#N}	\\
2CM	&	\smiles{CCCCCCCCn1c2ccccc2c2cc(Nc3nc(Nc4ccc5c(c4)c4ccccc4n5CCCCCCCC)nc(OC)n3)ccc21}	&	Na-DMAC	&	\smiles{CC1(C)c2ccccc2N(c2ccc(-c3c(C\#N)c(-c4cccc5ccccc45)nc(-c4cccc5ccccc45)c3C\#N)cc2)c2ccccc21}	\\
CzoB	&	\smiles{Cc1cc(C)c(B(c2ccccc2-n2c3ccccc3c3ccccc32)c2c(C)cc(C)cc2C)c(C)c1}	&	3Py-DMAC	&	\smiles{CC1(C)c2ccccc2N(c2ccc(-c3c(C\#N)c(-c4cccnc4)nc(-c4cccnc4)c3C\#N)cc2)c2ccccc21}	\\
CzBroB	&	\smiles{Cc1cc(C)c(B(c2ccc(Br)cc2-n2c3ccccc3c3ccccc32)c2c(C)cc(C)cc2C)c(C)c1}	&	6tBPA	&	\smiles{CC(C)(C)c1ccc(N(c2ccc(C(C)(C)C)cc2)c2c3ccc(C(C)(C)C)cc3c(N(c3ccc(C(C)(C)C)cc3)c3ccc(C(C)(C)C)cc3)c3ccc(C(C)(C)C)cc23)cc1}	\\
PxCN	&	\smiles{N\#Cc1c(-c2ccc(N3c4ccccc4Oc4ccccc43)cc2)cncc1-c1ccc(N2c3ccccc3Oc3ccccc32)cc1}	&	C-O	&	\smiles{COc1ccc2c(c1)c1ccccc1n2-c1cc2c(cc1-n1c3ccccc3c3ccccc31)Cc1ccccc1O2}	\\
AcCN	&	\smiles{CC1(C)c2ccccc2N(c2ccc(-c3cncc(-c4ccc(N5c6ccccc6C(C)(C)c6ccccc65)cc4)c3C\#N)cc2)c2ccccc21}	&	IDFL-2DPA	&	\smiles{CC1(C)c2ccccc2-c2c1cc1c3cc(N(c4ccccc4)c4ccccc4)ccc3n3c4ccc(N(c5ccccc5)c5ccccc5)cc4c2c13}	\\
CzCNoB	&	\smiles{Cc1cc(C)c(B(c2ccc(C\#N)cc2-n2c3ccccc3c3ccccc32)c2c(C)cc(C)cc2C)c(C)c1}	&	ACRF	&	\smiles{N\#Cc1ccc2c(c1)C1(c3cc(C\#N)ccc3-2)c2ccccc2N(c2ccccc2)c2ccccc21}	\\
CNPP-TPA	&	\smiles{N\#Cc1nc2c3ccc(-c4ccc(N(c5ccccc5)c5ccccc5)cc4)nc3c3nc(-c4ccc(N(c5ccccc5)c5ccccc5)cc4)ccc3c2nc1C\#N}	&	TPSi-F	&	\smiles{c1ccc(C2(c3ccc([Si](c4ccccc4)(c4ccccc4)c4ccccc4)cc3)c3ccccc3-c3ccccc32)cc1}	\\
SPAC-FO	&	\smiles{O=C1c2ccc(N3c4ccccc4C4(c5ccccc5-c5ccccc54)c4ccccc43)cc2-c2cc(N3c4ccccc4C4(c5ccccc5-c5ccccc54)c4ccccc43)ccc21}	&	PAB	&	\smiles{c1ccc(N(c2ccccc2)c2cc3c4c(c2)N(c2ccccc2)c2ccccc2B4c2ccccc2N3c2ccccc2)cc1}	\\
DMAC-FO	&	\smiles{CC1(C)c2ccccc2N(c2ccc3c(c2)-c2cc(N4c5ccccc5C(C)(C)c5ccccc54)ccc2C3=O)c2ccccc21}	&	2tPAB	&	\smiles{CC(C)(C)c1ccc(N2c3ccc(C(C)(C)C)cc3B3c4cc(C(C)(C)C)ccc4N(c4ccc(C(C)(C)C)cc4)c4cc(N(c5ccccc5)c5ccccc5)cc2c43)cc1}	\\
6AcBIQ	&	\smiles{CC1(C)c2ccccc2N(c2ccc3c4c(cccc24)C(=O)N(c2ccccc2)C3=O)c2ccccc21}	&	pDTPA-DTPZ	&	\smiles{c1ccc(N(c2ccccc2)c2ccc(-c3ccc(-c4ccc(N(c5ccccc5)c5ccccc5)cc4)c4nc5c(nc34)c3sccc3c3ccsc35)cc2)cc1}	\\
PtB	&	\smiles{CC(C)(C)c1ccc(-n2c(-c3ccc4ccc5cccc6ccc3c4c56)nc(-c3ccccc3)c2-c2ccccc2)cc1}	&	oDTPA-DTPZ	&	\smiles{c1ccc(N(c2ccccc2)c2ccc(-c3cc4nc5c(nc4cc3-c3ccc(N(c4ccccc4)c4ccccc4)cc3)c3sccc3c3ccsc35)cc2)cc1}	\\
BCPBT	&	\smiles{c1ccc2c(c1)c1ccccc1n2-c1cc(-c2ccc(-c3cc(-n4c5ccccc5c5ccccc54)cc(-n4c5ccccc5c5ccccc54)c3)c3nsnc23)cc(-n2c3ccccc3c3ccccc32)c1}	&	oTPAPO-DTPZ	&	\smiles{O=P(c1ccccc1)(c1ccccc1)c1cc2nc3c(nc2cc1-c1ccc(N(c2ccccc2)c2ccccc2)cc1)c1sccc1c1ccsc13}	\\
InDM	&	\smiles{CN(C)c1ccc(-n2c(-c3ccccc3)nc(-c3ccccc3)c2-c2ccccc2)cc1}	&	BPZPP	&	\smiles{c1ccc(-c2cc3nc4c5ccc(N6c7ccccc7Oc7ccccc76)nc5c5nc(N6c7ccccc7Oc7ccccc76)ccc5c4nc3cc2-c2ccccc2)cc1}	\\
PnB	&	\smiles{CCCCc1ccc(-n2c(-c3ccc4ccc5cccc6ccc3c4c56)nc(-c3ccccc3)c2-c2ccccc2)cc1}	&	BIPOCz	&	\smiles{c1ccc(-c2nc(-c3ccc(OCCCCc4cccc5c4[nH]c4ccccc45)cc3)n(-c3ccccc3)c2-c2ccccc2)cc1}	\\
BANE	&	\smiles{c1ccc(-c2ccc(-c3c4ccccc4c(-c4c5ccccc5c(-c5ccc(-c6ccccc6)cc5)c5ccccc45)c4ccccc34)cc2)cc1}	&	tDPAC-BN	&	\smiles{CC(C)(C)c1ccc(N2c3ccc(C(C)(C)C)cc3B3c4cc(C(C)(C)C)ccc4N(c4ccc(C(C)(C)C)cc4)c4cc(N5c6ccccc6C(c6ccccc6)(c6ccccc6)c6ccccc65)cc2c43)cc1}	\\
CBM	&	\smiles{CCCCn1c2ccc(C=C(C\#N)C\#N)cc2c2cc(C=C(C\#N)C\#N)ccc21}	&	tDMAC-BN	&	\smiles{CC(C)(C)c1ccc(N(c2ccc(C(C)(C)C)cc2)c2cc3c4c(c2)N2c5ccccc5C(C)(C)c5cccc(c52)B4c2cc(C(C)(C)C)ccc2N3c2ccc(C(C)(C)C)cc2)cc1}	\\
DCzBNPh-1	&	\smiles{N\#Cc1ccc(-c2cccc(-c3ccc(C\#N)c(-n4c5ccccc5c5ccccc54)c3)c2)cc1-n1c2ccccc2c2ccccc21}	&	2ICz-OXD	&	\smiles{c1ccc(-n2c3ccccc3c3c2ccc2c4ccccc4n(-c4ccc(-c5nnc(-c6ccc(-n7c8ccccc8c8ccc9c(c\%10ccccc\%10n9-c9ccccc9)c87)cc6)o5)cc4)c23)cc1}	\\
MCZ-P2-DTM	&	\smiles{COc1ccc(-c2ccc3c(c2)c2cc(-c4ccc(OC)cc4)ccc2n3-c2ccc(C(=O)c3ccc4sc5ccccc5c4c3)cn2)cc1}	&	SIR-B	&	\smiles{Cc1ccc2c(c1)C1(c3ccc(-c4ccccc4)cc3C(=O)c3ccc(-c4ccccc4)cc31)c1cc(C)ccc1N2c1ccc(C(C)(C)C)cc1}	\\
DBTO-PTZ	&	\smiles{O=S1(=O)c2ccc(N3c4ccccc4Oc4ccccc43)cc2-c2cc(N3c4ccccc4Oc4ccccc43)ccc21}	&	SIR-C	&	\smiles{Cc1ccc2c(c1)C1(c3ccc(-c4ccccc4-c4ccccc4)cc3C(=O)c3ccc(-c4ccccc4-c4ccccc4)cc31)c1cc(C)ccc1N2c1ccc(C(C)(C)C)cc1}	\\
$\alpha$-DMAC-DBP	&	\smiles{CC1(C)c2ccccc2N(c2cccc3nc4c5ccccc5c5ccccc5c4nc23)c2ccccc21}	&	SIR-A	&	\smiles{Cc1ccc2c(c1)C1(c3ccccc3C(=O)c3ccccc31)c1cc(C)ccc1N2c1ccc(C(C)(C)C)cc1}	\\
$\alpha$-2DMAC-DBP	&	\smiles{CC1(C)c2ccccc2N(c2ccc(N3c4ccccc4C(C)(C)c4ccccc43)c3nc4c5ccccc5c5ccccc5c4nc23)c2ccccc21}	&	TRZ3DPA	&	\smiles{c1ccc(N(c2ccccc2)c2ccc(-c3nc(-c4ccc(N(c5ccccc5)c5ccccc5)cc4)nc(-c4ccc(N(c5ccccc5)c5ccccc5)cc4)n3)cc2)cc1}	\\
$\beta$-DI-DBP	&	\smiles{c1ccc(-n2c3ccccc3c3c2c2c4ccccc4n(-c4ccc5nc6c7ccccc7c7ccccc7c6nc5c4)c2c2c4ccccc4n(-c4ccccc4)c32)cc1}	&	BO3N	&	\smiles{Nc1ccc2c(c1)Oc1cc(N)cc3c1B2c1ccc(Nc2cccc4cc5ccc6cc7cc8ccccc8cc7cc6c5cc24)cc1O3}	\\
DIC-TRZ	&	\smiles{c1ccc(-c2nc(-c3ccccc3)nc(-n3c4ccccc4c4ccc5c6ccccc6n(-c6ccccc6)c5c43)n2)cc1}	&	BPPyA	&	\smiles{Cc1ccc(C)c(N(c2ccc3ccc4c(N(c5cc(C)ccc5C)c5cccc6c5oc5ccccc56)ccc5ccc2c3c54)c2cccc3c2oc2ccccc23)c1}	\\
PO-TPA	&	\smiles{CC(C)(C)c1ccc(N(c2ccc(-c3cc4c(cc3C\#N)=NC3c5ccc(P(=O)(c6ccccc6)c6ccccc6)nc5-c5nc(P(=O)(c6ccccc6)c6ccccc6)ccc5C3N=4)cc2)c2ccc(C(C)(C)C)cc2)cc1}	&	SiCN2Cz	&	\smiles{N\#Cc1cc(-c2nc(-n3c4ccccc4c4ccccc43)nc(-n3c4ccccc4c4ccccc43)n2)cc([Si](c2ccccc2)(c2ccccc2)c2ccccc2)c1}	\\
DMAC-Cz-TTR	&	\smiles{CC1(C)c2ccccc2N(c2ccc3c(c2)c2cc(N4c5ccccc5C(C)(C)c5ccccc54)ccc2n3-c2ccc3c(c2)S(=O)(=O)c2ccccc2S3(=O)=O)c2ccccc21}	&	PhPC	&	\smiles{c1ccc(-c2c3ccccc3c(-c3ccc(-n4c5ccccc5c5ccccc54)cc3)c3ccccc23)cc1}	\\
DCzPO	&	\smiles{Cc1ccc2c(c1)c1cc(C)ccc1n2-c1cccc(-n2c3ccc(C)cc3c3cc(C)ccc32)c1P(=O)(c1ccccc1)c1ccccc1}	&	1MOBI	&	\smiles{COc1cc(/C=C/c2ccc(-n3cnc4ccccc43)cc2)c(OC)cc1/C=C/c1ccc(-n2cnc3ccccc32)cc1}	\\
IQ-oTPA	&	\smiles{O=C1c2ccccc2-c2nc3cc(-c4ccc(N(c5ccccc5)c5ccccc5)cc4)c(-c4ccc(N(c5ccccc5)c5ccccc5)cc4)cc3nc21}	&	3BrBI	&	\smiles{Brc1cc(/C=C/c2ccc(-n3cnc4ccccc43)cc2)c(Br)cc1/C=C/c1ccc(-n2cnc3ccccc32)cc1}	\\
2C	&	\smiles{N\#CC(C\#N)=Cc1ccc(-c2ccc3ccc4cccc5ccc2c3c45)cc1F}	&	DMBFTX	&	\smiles{CC1(C)c2ccccc2-c2cccc(-c3ccc4sc5ccc(-c6cccc7c6C(C)(C)c6ccccc6-7)cc5c(=O)c4c3)c21}	\\
C-H	&	\smiles{N\#CC(C\#N)=Cc1c(F)cc(-c2ccc3ccc4cccc5ccc2c3c45)cc1F}	&	SBF-AnCN	&	\smiles{N\#Cc1ccc(-c2c3ccccc3c(-c3ccc4c(c3)C3(c5ccccc5-c5ccccc53)c3ccccc3-4)c3ccccc23)cc1}	\\
PSPP	&	\smiles{O=S(=O)(c1cccnc1)c1cccc(N2c3ccccc3Oc3ccccc32)c1}	&	DMF-AnCN	&	\smiles{CC1(C)c2ccccc2-c2ccc(-c3c4ccccc4c(-c4ccc(C\#N)cc4)c4ccccc34)cc21}	\\
2DN	&	\smiles{CC(C)c1cccc(C(C)C)c1N1C(=O)c2cc(N3c4ccccc4C(C)(C)c4ccccc43)cc3cc(N4c5ccccc5C(C)(C)c5ccccc54)cc(c23)C1=O}	&	DPF-AnCN	&	\smiles{N\#Cc1ccc(-c2c3ccccc3c(-c3ccc4c(c3)C(c3ccccc3)(c3ccccc3)c3ccccc3-4)c3ccccc23)cc1}	\\
PSPBP	&	\smiles{O=S(=O)(c1cccnc1)c1cc(N2c3ccccc3Oc3ccccc32)cc(N2c3ccccc3Oc3ccccc32)c1}	&	DPA	&	\smiles{c1ccc(-c2c3ccccc3c(-c3ccccc3)c3ccccc23)cc1}	\\
2PN	&	\smiles{CC(C)c1cccc(C(C)C)c1N1C(=O)c2cc(N3c4ccccc4Oc4ccccc43)cc3cc(N4c5ccccc5Oc5ccccc54)cc(c23)C1=O}	&	TPA–DAP–TPA	&	\smiles{c1ccc(N(c2ccccc2)c2ccc(-c3c4ccccc4c(-c4ccc5ccc6c(-c7c8ccccc8c(-c8ccc(N(c9ccccc9)c9ccccc9)cc8)c8ccccc78)ccc7ccc4c5c76)c4ccccc34)cc2)cc1}	\\
pCbz-PYR	&	\smiles{c1ccc(-c2ccc3c(c2)c2cc(-c4ccccc4)ccc2n3-c2cc(-n3c4ccc(-c5ccccc5)cc4c4cc(-c5ccccc5)ccc43)ncn2)cc1}	&	AcDbp	&	\smiles{CC1(C)c2ccccc2N(c2ccc3c4ccc(N5c6ccccc6C(C)(C)c6ccccc65)cc4c4nc5ccccc5nc4c3c2)c2ccccc21}	\\
PPI-An-NPIM	&	\smiles{c1ccc(-c2nc(-c3ccccc3)n(-c3ccc(-c4c5ccccc5c(-c5ccc(-c6nc7c8ccccc8c8ccccc8c7n6-c6ccccc6)cc5)c5ccccc45)cc3)c2-c2ccccc2)cc1}	&	4-Ac	&	\smiles{CC(C)(C)c1ccc2c(c1)C(C)(C)c1cc(C(C)(C)C)ccc1N2c1ccc(C(F)(F)F)cc1}	\\
2PS-2FPh	&	\smiles{Fc1cc(N2c3ccccc3Sc3ccccc32)c(F)cc1N1c2ccccc2Sc2ccccc21}	&	czpzpy	&	\smiles{c1cc(-c2cccc3c2[nH]c2ccccc23)nc(-n2cccn2)c1}	\\
9CzTPE	&	\smiles{C(=Cc1ccc(-n2c3ccccc3c3ccccc32)cc1)c1ccc(C(=C(c2ccccc2)c2ccccc2)c2ccccc2)cc1}	&	NyDPO	&	\smiles{c1ccc2c(c1)Oc1ccccc1N2c1ccc(-c2ccc3nc(-c4ccc(N5c6ccccc6Oc6ccccc65)cc4)ccc3n2)cc1}	\\
2CzTPE	&	\smiles{CCn1c2ccccc2c2ccc(-c3ccc(C=Cc4ccc(C(=C(c5ccccc5)c5ccccc5)c5ccccc5)cc4)cc3)cc21}	&	1MPA	&	\smiles{Cc1nc(N2c3ccccc3C(c3ccccc3)(c3ccccc3)c3ccccc32)cc(N2c3ccccc3C(c3ccccc3)(c3ccccc3)c3ccccc32)n1}	\\
FTAT-FBO	&	\smiles{O=P(c1ccccc1)(c1ccccc1)c1ccc([Si](c2ccccc2)(c2ccccc2)c2ccccc2)cc1}	&	LEt	&	\smiles{CCc1ccccc1P(c1ccccc1CC)c1ccccc1P(c1ccccc1CC)c1ccccc1CC}	\\
DmPy-PPI	&	\smiles{c1ccc(-n2c(-c3ccc(-c4cc(-c5cccc(-c6cccnc6)c5)cc(-c5cccc(-c6cccnc6)c5)c4)cc3)nc3c4ccccc4c4ccccc4c32)cc1}	&	DC-TC	&	\smiles{CC(C)(C)c1ccc2c(c1)c1cc(C(C)(C)C)ccc1n2-c1ccc(C(=O)C(=O)c2ccc(-n3c4ccc(C(C)(C)C)cc4c4cc(C(C)(C)C)ccc43)cc2)cc1}	\\
mCP-BmPy	&	\smiles{CC(C)(C)c1ccc2c(c1)C(c1cc(-c3cccnc3)cc(-c3cccnc3)c1)(c1ccc3c(c1)c1ccccc1n3-c1cccc(-n3c4ccccc4c4ccccc43)c1)c1cc(C(C)(C)C)ccc1-2}	&	DC-ACR	&	\smiles{CC1(C)c2ccccc2N(c2ccc(C(=O)C(=O)c3ccc(N4c5ccccc5C(C)(C)c5ccccc54)cc3)cc2)c2ccccc21}	\\
mCP-Py	&	\smiles{CC(C)(C)c1ccc2c(c1)C(c1cccnc1)(c1ccc3c(c1)c1ccccc1n3-c1cccc(-n3c4ccccc4c4ccccc43)c1)c1cc(C(C)(C)C)ccc1-2}	&	1G-F1	&	\smiles{c1ccc(P(c2ccccc2)c2cscc2P(c2ccccc2)c2ccccc2)cc1}	\\
mCP-Ph	&	\smiles{CC(C)(C)c1ccc2c(c1)C(c1ccccc1)(c1ccc3c(c1)c1ccccc1n3-c1cccc(-n3c4ccccc4c4ccccc43)c1)c1cc(C(C)(C)C)ccc1-2}	&	P–Cu–C	&	\smiles{c1ccc(P(c2ccccc2)c2ccccc2-c2ccccc2P(c2ccccc2)c2ccccc2)cc1}	\\
DMTAC	&	\smiles{Cc1ccc(-c2c(-c3ccccc3)c(-c3ccccc3)c3c(N)cc4c5ccccc5c(N)cc4c3c2-c2ccc(C)c(C)c2)cc1C}	&	BPCP-2CPC	&	\smiles{N\#Cc1cc(-n2c3ccccc3c3cc(-c4ccc(C5(c6ccc(-n7c8ccccc8c8ccccc87)cc6)CCCCC5)cc4)ccc32)c(-n2c3ccccc3c3cc(-c4ccc(C5(c6ccc(-n7c8ccccc8c8ccccc87)cc6)CCCCC5)cc4)ccc32)cc1C\#N}	\\
IPTAC	&	\smiles{Cc1ccc(-c2c(-c3ccc(C(C)C)cc3)c(-c3ccc(C(C)C)cc3)c3c(c(N)cc4c5ccccc5c(N)cc43)c2-c2ccc(C)cc2)cc1}	&	2H-Qz	&	\smiles{c1ccc(-c2cc(-c3ccccn3)cc(-c3nc(-c4ccc(N5c6ccccc6Oc6ccccc65)cc4)c4ccccc4n3)c2)nc1}	\\
TBTAC	&	\smiles{Cc1ccccc1-c1c(-c2ccc(C(C)(C)C)cc2)c(-c2ccc(C(C)(C)C)cc2)c2c(c(N)cc3c4ccccc4c(N)cc32)c1-c1ccccc1C}	&	PPINCN-Cz	&	\smiles{N\#Cc1ccc(-n2c(-c3ccc4ccc5cccc6ccc3c4c56)nc3c4ccc(-c5ccc(-n6c7ccccc7c7ccccc76)cc5)cc4c4cc(-c5ccc(-n6c7ccccc7c7ccccc76)cc5)ccc4c32)c2ccccc12}	\\
p2Cz2Trz	&	\smiles{c1ccc(-c2nc(-c3ccccc3)nc(-c3cc(-n4c5ccccc5c5ccccc54)c(-c4nc(-c5ccccc5)nc(-c5ccccc5)n4)cc3-n3c4ccccc4c4ccccc43)n2)cc1}	&	PPINCN	&	\smiles{N\#Cc1ccc(-n2c(-c3ccc4ccc5cccc6ccc3c4c56)nc3c4ccccc4c4ccccc4c32)c2ccccc12}	\\
m2Cz2Trz	&	\smiles{c1ccc(-c2nc(-c3ccccc3)nc(-c3cc(-c4nc(-c5ccccc5)nc(-c5ccccc5)n4)c(-n4c5ccccc5c5ccccc54)cc3-n3c4ccccc4c4ccccc43)n2)cc1}	&	BDAPM	&	\smiles{O=C(c1ccc(N2c3ccccc3[Si](c3ccccc3)(c3ccccc3)c3ccccc32)cc1)c1ccc(N2c3ccccc3[Si](c3ccccc3)(c3ccccc3)c3ccccc32)cc1}	\\
2,4-CzPyCl3	&	\smiles{Clc1nc(-n2c3ccccc3c3ccccc32)c(Cl)c(-n2c3ccccc3c3ccccc32)c1Cl}	&	12AcCz-PM	&	\smiles{CC1(C)c2ccccc2N(c2ccc(-c3cc(-c4ccccc4)nc(-c4ccccc4)n3)cc2)c2c1ccc1c2c2ccccc2n1-c1ccccc1}	\\
DMAC-PCN	&	\smiles{CC1(C)c2ccccc2N(c2ccc(-c3c(C\#N)c(-c4ccccc4)nc(-c4ccccc4)c3C\#N)cc2)c2ccccc21}	&	23AcCz-PM	&	\smiles{CC1(C)c2ccccc2N(c2ccc(-c3cc(-c4ccccc4)nc(-c4ccccc4)n3)cc2)c2cc3c(cc21)c1ccccc1n3-c1ccccc1}	\\
2PXZP	&	\smiles{Cn1cnc2c(-c3ccc(N4c5ccccc5Oc5ccccc54)cc3)nc(-c3ccc(N4c5ccccc5Oc5ccccc54)cc3)nc21}	&	CDDPI	&	\smiles{c1ccc2c(-n3c4ccccc4c4ccccc43)ccc(-c3nc4c5ccccc5c5ccccc5c4n3-c3ccc4c(c3)OCCO4)c2c1}	\\
2,4,6-CzPyCl2	&	\smiles{Clc1c(-n2c3ccccc3c3ccccc32)nc(-n2c3ccccc3c3ccccc32)c(Cl)c1-n1c2ccccc2c2ccccc21}	&	DDBP	&	\smiles{c1ccc(-c2nc(-c3ccc(-c4ccc(N(c5ccccc5)c5ccccc5)cc4)c4ccccc34)n(-c3ccc4c(c3)OCCO4)c2-c2ccccc2)cc1}	\\
1PXZP	&	\smiles{Cn1cnc2c(-c3ccc(N4c5ccccc5Oc5ccccc54)cc3)ncnc21}	&	CCDD	&	\smiles{O=C(c1ccccc1)c1ccc(-n2c3ccc(N(c4ccccc4)c4ccccc4)cc3c3cc(N(c4ccccc4)c4ccccc4)ccc32)cc1}	\\
3,6-T2C	&	\smiles{N\#Cc1ccc(-c2nc3c4ccc(-c5ccc(N(c6ccccc6)c6ccccc6)cc5)cc4c4cc(-c5ccc(N(c6ccccc6)c6ccccc6)cc5)ccc4c3n2-c2ccc(C\#N)cc2)cc1}	&	DPNC	&	\smiles{COc1ccc(N(c2ccc(OC)cc2)c2ccc3c(c2)c2cc(N(c4ccc(OC)cc4)c4ccc(OC)cc4)ccc2n3-c2cccc3ccccc23)cc1}	\\
3,6-TC	&	\smiles{N\#Cc1ccc(-c2nc3c4ccc(-c5ccc(N(c6ccccc6)c6ccccc6)cc5)cc4c4cc(-c5ccc(N(c6ccccc6)c6ccccc6)cc5)ccc4c3n2-c2ccccc2)cc1}	&	LiPr	&	\smiles{CC(C)c1ccccc1P(c1ccccc1C(C)C)c1ccccc1P(c1ccccc1C(C)C)c1ccccc1C(C)C}	\\
Ph-CF3	&	\smiles{c1ccc2c(c1)c1ccccc1n2-c1ccc(Nc2ccc(-n3c4ccccc4c4ccccc43)cc2)cc1}	&	Spiro-CN	&	\smiles{Cc1ccc(N(c2ccc(C)cc2)c2ccc3c(c2)C2(c4cc(C\#N)ccc4-c4ccc(C\#N)cc42)c2cc(N(c4ccc(C)cc4)c4ccc(C)cc4)ccc2-3)cc1}	\\
arylene	&	\smiles{c1ccc2c(c1)c1ccccc1n2-c1ccc(Nc2ccc(-n3c4ccccc4c4ccccc43)cc2)cc1}	&	TMDBP	&	\smiles{CC1(C)c2ccccc2Nc2c1ccc1c2C(C)(C)c2ccccc2N1}	\\
TS-1	&	\smiles{O=S(=O)(c1ccccc1)c1ccc(-n2c3ccccc3c3ccccc32)cc1}	&	ACRSA	&	\smiles{O=C1c2ccccc2C2(c3ccccc31)c1ccccc1N(c1ccccc1)c1ccccc12}	\\
3,6-MeCzOXD	&	\smiles{Cc1ccc2c(c1)c1cc(C)ccc1n2-c1ccccc1-c1nnc(-c2ccccc2-n2c3ccc(C)cc3c3cc(C)ccc32)o1}	&	CAPI	&	\smiles{c1ccc2c(-n3c(-c4ccc(-c5c6ccccc6c(-n6c7ccccc7c7ccccc76)c6ccccc56)cc4)nc4c5ccccc5c5ccccc5c43)cccc2c1}	\\
p-PO15NCzDPA	&	\smiles{O=P(c1ccccc1)(c1ccccc1)c1ccc(-c2c3ccccc3c(-c3ccc(-n4c5cccnc5c5cccnc54)cc3)c3ccccc23)cc1}	&	PhCzSPOTz	&	\smiles{O=P(c1ccccc1)(c1ccccc1)c1nc(-c2ccccc2)nc(-n2c3ccccc3c3ccccc32)n1}	\\
2,7-MeCzOXD	&	\smiles{Cc1ccc2c3ccc(C)cc3n(-c3ccccc3-c3nnc(-c4ccccc4-n4c5cc(C)ccc5c5ccc(C)cc54)o3)c2c1}	&	CCAPI	&	\smiles{N\#Cc1ccc(-n2c(-c3ccc(-c4c5ccccc5c(-n5c6ccccc6c6ccccc65)c5ccccc45)cc3)nc3c4ccccc4c4ccccc4c32)c2ccccc12}	\\
BPCN-Cz2Ph	&	\smiles{N\#Cc1ccccc1-c1ccc(-n2c3ccc(-c4ccccc4)cc3c3cc(-c4ccccc4)ccc32)cc1}	&	2CzCbPy	&	\smiles{c1cc(-n2c3ccccc3c3ccccc32)nc(-n2c3ccc(-n4c5ccccc5c5ccccc54)cc3c3ncccc32)c1}	\\
o-CzOXD	&	\smiles{c1ccc(-n2c3ccccc3c3ccccc32)c(-c2nnc(-c3ccccc3-n3c4ccccc4c4ccccc43)o2)c1}	&	CzCbPy	&	\smiles{c1ccc(-n2c3ccc(-n4c5ccccc5c5ccccc54)cc3c3ncccc32)nc1}	\\
BPCN-2Cz	&	\smiles{N\#Cc1ccccc1-c1ccc(-n2c3ccccc3c3cc(-c4ccc5c(c4)c4ccccc4n5-c4ccccc4)ccc32)cc1}	&	DCE	&	\smiles{CC(C)CCC[C@@H](C)[C@H]1CC[C@H]2[C@@H]3CC=C4C[C@@H](O)CC[C@]4(C)[C@H]3CC[C@]12C}	\\
Cbz-MI	&	\smiles{CCCCCCN1C(=O)C(c2ccc(-c3ccc(-n4c5ccccc5c5ccccc54)cc3)cc2)=C(c2ccc(-c3ccc(-n4c5ccccc5c5ccccc54)cc3)cc2)C1=O}	&	PICFOCz	&	\smiles{FC(F)(F)c1cccc(-n2c(-c3ccc(OCCCCn4c5ccccc5c5ccccc54)cc3)nc3c4ccccc4c4ccccc4c32)c1}	\\
DBFDPO	&	\smiles{O=P(c1ccccc1)(c1ccccc1)c1cccc2c1oc1c(P(=O)(c3ccccc3)c3ccccc3)cccc12}	&	dpmt	&	\smiles{Cc1sc(C)c(P(c2ccccc2)c2ccccc2)c1P(c1ccccc1)c1ccccc1}	\\
XCBP	&	\smiles{Cc1cc(-n2c3ccccc3c3ccccc32)c(C)cc1-c1ccc(C(=O)c2ccccc2)cc1}	&	BICFOCz	&	\smiles{FC(F)(F)c1cccc(-n2c(-c3ccc(OCCCCn4c5ccccc5c5ccccc54)cc3)nc(-c3ccccc3)c2-c2ccccc2)c1}	\\
Cbz-PYR	&	\smiles{c1ccc2c(c1)c1ccccc1n2-c1cc(-n2c3ccccc3c3ccccc32)ncn1}	&	19a	&	\smiles{c1ccc2cc(-c3ccc4cc5cc(-c6ccc7ccccc7c6)ccc5cc4c3)ccc2c1}	\\
tCbz-mPYRs	&	\smiles{CSc1nc(-n2c3ccc(C(C)(C)C)cc3c3c(C(C)(C)C)cccc32)c(C)c(-n2c3ccc(C(C)(C)C)cc3c3c(C(C)(C)C)cccc32)n1}	&	BTPBFCz	&	\smiles{c1ccc(-c2nc(-c3ccccc3-n3c4ccccc4c4c5oc6ccccc6c5ccc43)c3sc4ccccc4c3n2)cc1}	\\
PCBP	&	\smiles{O=C(c1ccccc1)c1ccc(-c2ccc(-n3c4ccccc4c4ccccc43)cc2)cc1}	&	Ph-OBNA	&	\smiles{c1ccc(-n2c3ccccc3c3c4c5c(cc32)Oc2ccccc2B5c2ccccc2O4)cc1}	\\
MeCZ-TRZ	&	\smiles{Cc1ccc2c(c1)c1ccccc1n2-c1ccc(-c2nc(-c3ccc(-n4c5ccccc5c5cc(C)ccc54)c(C)c3)nc(-c3ccc(-n4c5ccccc5c5cc(C)ccc54)c(C)c3)n2)cc1C}	&	HPVAnt	&	\smiles{CCCCCCc1ccc(C=Cc2ccc3cc4ccccc4cc3c2)cc1}	\\
3ca	&	\smiles{COc1ccc(-c2cc(-c3ccc(C(F)(F)F)cc3)c(-c3nc4ccccc4[nH]3)s2)cc1}	&	BTPBFCz-D3	&	\smiles{c1ccc(-c2nc(-c3ccc(-n4c5ccccc5c5ccc6c(c7ccccc7n6-c6ccccc6)c54)cc3-n3c4ccccc4c4c5oc6ccccc6c5ccc43)c3sc4ccccc4c3n2)cc1}	\\
2MeCZ-TRZ	&	\smiles{Cc1ccc2c(c1)c1cc(C)ccc1n2-c1ccc(-c2nc(-c3ccc(-n4c5ccc(C)cc5c5cc(C)ccc54)c(C)c3)nc(-c3ccc(-n4c5ccc(C)cc5c5cc(C)ccc54)c(C)c3)n2)cc1C}	&	ICBNTrz4	&	\smiles{CC1(C)c2ccccc2-c2cc3c4ccccc4n(-c4ccc(-c5nc(-c6ccccc6)nc(-c6ccccc6)n5)cc4C\#N)c3cc21}	\\
CZ-TRZ	&	\smiles{Cc1cc(-c2nc(-c3ccc(-n4c5ccccc5c5ccccc54)c(C)c3)nc(-c3ccc(-n4c5ccccc5c5ccccc54)c(C)c3)n2)ccc1-n1c2ccccc2c2ccccc21}	&	1PCTRz	&	\smiles{c1ccc(-c2nc(-c3ccccc3)nc(-c3ccc(-n4c5ccccc5c5cc(-n6c7ccccc7c7ccccc76)ccc54)c(-c4ccccc4)c3)n2)cc1}	\\
DMAC-AQ	&	\smiles{CC1(C)c2ccccc2N(c2ccc3c(c2)C(=O)c2ccccc2C3=O)c2ccccc21}	&	DMACN-B	&	\smiles{Cc1cc(C)c(N(c2ccc3c(c2)C(c2ccccc2)(c2ccccc2)c2cccc4c2N3c2ccccc2C4(C)C)c2c(C)cc(C)cc2C)c(C)c1}	\\
DTPTCzDP-CN	&	\smiles{CC(C)(C)c1nc(-c2ccc(-n3c4ccc(N(c5ccccc5)c5ccccc5)cc4c4cc(N(c5ccccc5)c5ccccc5)ccc43)c(C\#N)c2)nc(C(C)(C)C)n1}	&	CzSO-1CO	&	\smiles{CC1(OCOC(=O)c2ccccc2)CC=Cc2c1c1ccccc1n2-c1ccc(S(=O)(=O)c2ccccc2)cc1}	\\
mP2MPC	&	\smiles{Cc1cc(-c2cccc(-n3c4ccccc4c4ccccc43)c2)c(C)cc1-c1cccc(-c2nc3c4ccccc4c4ccccc4c3n2-c2ccccc2)c1}	&	PPO21	&	\smiles{O=P(c1ccccc1)(c1ccccc1)c1ccc(-n2c3ccccc3c3cc(P(=O)(c4ccccc4)c4ccccc4)ccc32)cc1}	\\
XPhos	&	\smiles{CC(C)C1=CC(C(C)C)=C(c2ccccc2)C(C(C)C)(P(C2CCCCC2)C2CCCCC2)C1}	&	OXD-7	&	\smiles{CC(C)(C)c1ccc(-c2nnc(-c3cccc(-c4nnc(-c5ccc(C(C)(C)C)cc5)o4)c3)o2)cc1}	\\
TPXZ-as-TAZ	&	\smiles{c1ccc2c(c1)Oc1ccccc1N2c1ccc(-c2nnc(-c3ccc(N4c5ccccc5Oc5ccccc54)cc3)c(-c3ccc(N4c5ccccc5Oc5ccccc54)cc3)n2)cc1}	&	TCZ1	&	\smiles{CCCCC(CC)Cn1c2ccc(-n3c4ccccc4c4ccccc43)cc2c2cc(-n3c4ccccc4c4ccccc43)ccc21}	\\
2Trz6CNDBT	&	\smiles{N\#Cc1cccc2c1sc1ccc(-c3nc(-c4ccccc4)nc(-c4ccccc4)n3)cc12}	&	APPT	&	\smiles{c1cc2c3c(cccc3c1)-c1nc3c4cccnc4c4ncccc4c3nc1-2}	\\
TPAPyPI	&	\smiles{c1ccc(N(c2ccccc2)c2ccc(-c3ccc(-c4nc5c6ccccc6c6ccccc6c5n4-c4cccc(-c5ccc6ccc7cccc8ccc5c6c78)c4)cc3)cc2)cc1}	&	ANQDC	&	\smiles{N\#Cc1cc2nc3c(nc2cc1C\#N)-c1cccc2cccc-3c12}	\\
TPACzPI	&	\smiles{c1ccc(N(c2ccccc2)c2ccc(-c3ccc(-c4nc5c6ccccc6c6ccccc6c5n4-c4cccc(-c5ccc6c(c5)c5ccccc5n6-c5ccccc5)c4)cc3)cc2)cc1}	&	ECDPTT	&	\smiles{CCn1c2ccccc2c2cc(C(=C(c3ccccc3)c3ccccc3)c3ccc(-c4ccc5c(c4)S(=O)(=O)c4ccccc4S5(=O)=O)cc3)ccc21}	\\
NA-TNA	&	\smiles{CC(C)(C)c1ccc(N2C(=O)c3cccc4c(/C=C/c5ccc6cc(N(c7ccccc7)c7ccccc7)ccc6c5)ccc(c34)C2=O)cc1}	&	1CM	&	\smiles{CCCCCCCCn1c2ccccc2c2cc(Nc3nc(OC)nc(OC)n3)ccc21}	\\
CzPrSBI	&	\smiles{O=C1c2ccccc2S(=O)(=O)N1CCCn1c2ccccc2c2ccccc21}	&	NPE-AcDPS	&	\smiles{CC(c1ccccc1)N1c2cc(N3c4ccccc4C(C)(C)c4ccccc43)ccc2S(=O)(=O)c2ccc(N3c4ccccc4C(C)(C)c4ccccc43)cc21}	\\
m-CzCNDMT	&	\smiles{Cc1cc(-c2cccc(-n3c4ccccc4c4cc(C\#N)ccc43)c2)c(C)cc1-c1cccc(-n2c3ccccc3c3cc(C\#N)ccc32)c1}	&	DMAC-Cz	&	\smiles{CC1(C)c2ccccc2N(c2ccc3[nH]c4ccc(N5c6ccccc6C(C)(C)c6ccccc65)cc4c3c2)c2ccccc21}	\\
DTXO-TPA2	&	\smiles{Cc1cc(N(c2ccccc2)c2cccc3ccccc23)ccc1-c1ccc(N(c2ccccc2)c2cccc3ccccc23)cc1C}	&	VBC-PTZ	&	\smiles{CCN(CC)c1ccc2c(c1)oc(=O)c1c(=O)oc3cc(N4c5ccccc5Sc5ccccc54)ccc3c12}	\\
24PQ-Cz	&	\smiles{c1ccc2c(-c3ccc(-n4c5ccccc5c5ccccc54)cc3)nc(-c3ccc(-n4c5ccccc5c5ccccc54)cc3)nc2c1}	&	t-BuPCz	&	\smiles{CC(C)(C)c1ccc(-c2ccc3[nH]c4ccc(-c5ccc(C(C)(C)C)cc5)cc4c3c2)cc1}	\\
SFPC	&	\smiles{O=S(=O)(c1ccc(N2c3ccccc3Sc3ccccc32)cc1)c1ccc(-n2c3ccccc3c3ccccc32)cc1}	&	4BPy-mDTC	&	\smiles{CC(C)(C)c1ccc2c(c1)c1cc(C(C)(C)C)ccc1n2-c1cc(C(=O)c2ccncc2)cc(-n2c3ccc(C(C)(C)C)cc3c3cc(C(C)(C)C)ccc32)c1}	\\
2CzMC	&	\smiles{COc1cc(=O)oc2cc3c(cc12)c1ccccc1n3-c1ccccc1}	&	DC–TC	&	\smiles{Cc1ccc(N(c2ccc(C)cc2)c2ccc(C3(c4ccc(N(c5ccc(C)cc5)c5ccc(C)cc5)cc4)CCCCC3)cc2)cc1}	\\
PPI-PIM-2	&	\smiles{FC(F)(F)c1ccc(-n2c(-c3ccc(-c4ccc(-n5c(-c6ccccc6)nc(-c6ccccc6)c5-c5ccccc5)cc4)cc3)nc3c4ccccc4c4ccccc4c32)cc1}	&	DFDBQPXZ	&	\smiles{Fc1cc2nc(-c3ccc(N4c5ccccc5Oc5ccccc54)cc3)c(-c3ccc(N4c5ccccc5Oc5ccccc54)cc3)nc2cc1F}	\\
4CzMC	&	\smiles{COc1cc(=O)oc2c1ccc1c2c2ccccc2n1-c1ccccc1}	&	DPE	&	\smiles{C(\#Cc1ccc(N2c3ccccc3Oc3ccccc32)cc1)c1ccc(N2c3ccccc3Oc3ccccc32)cc1}	\\
PPI-PIM-1	&	\smiles{FC(F)(F)c1ccc(-n2c(-c3ccc(-c4ccc(-c5nc(-c6ccccc6)c(-c6ccccc6)n5-c5ccccc5)cc4)cc3)nc3c4ccccc4c4ccccc4c32)cc1}	&	CBZ	&	\smiles{CCCCCCOc1ccc2c3ccc(OCCCCCC)cc3n(-c3ccnc(-c4ccc(F)c(C\#N)c4F)c3)c2c1}	\\
CzTPA	&	\smiles{c1ccc(N(c2ccccc2)c2ccccc2-n2c3ccccc3c3ccccc32)cc1}	&	tBuCz2	&	\smiles{CC(C)(C)c1ccc2c(c1)c1cc(C(C)(C)C)ccc1n2-c1ccc2[nH]c3ccc(-n4c5ccc(C(C)(C)C)cc5c5cc(C(C)(C)C)ccc54)cc3c2c1}	\\
SFX-bPy	&	\smiles{c1ccc2c(c1)Oc1ccc(-c3ccc4ccccc4n3)cc1C21c2ccccc2-c2ccccc21}	&	26DCzPPy	&	\smiles{c1cc(-c2cccc(-c3cccc(-n4c5ccccc5c5ccccc54)c3)n2)cc(-n2c3ccccc3c3ccccc32)c1}	\\
t–BuCz2BP	&	\smiles{CC(C)(C)c1ccc2c(c1)c1cc(C(C)(C)C)ccc1n2-c1ccc(C(=O)c2ccc(-n3c4ccc(C(C)(C)C)cc4c4cc(C(C)(C)C)ccc43)cc2)cc1}	&	mCBP	&	\smiles{c1cc(-c2cccc(-n3c4ccccc4c4ccccc43)c2)cc(-n2c3ccccc3c3ccccc32)c1}	\\
SAF-2CzB	&	\smiles{CC(C)(C)c1ccc2c(c1)c1cc(C(C)(C)C)ccc1n2-c1ccc(-c2ccc3c(c2)C2(c4cc(-c5ccc(-n6c7ccc(C(C)(C)C)cc7c7cc(C(C)(C)C)ccc76)cc5)ccc4-3)c3ccccc3N(c3ccccc3)c3ccccc32)cc1}	&	benzoxazole	&	\smiles{Cc1sc2ccccc2c1C1=C(c2c(C)sc3ccccc23)C(F)(F)C(F)(F)C1(F)F}	\\
F-2CzB	&	\smiles{CC(C)(C)c1ccc2c(c1)c1cc(C(C)(C)C)ccc1n2-c1ccc(-c2ccc3c(c2)C(C)(C)c2cc(-c4ccc(-n5c6ccc(C(C)(C)C)cc6c6cc(C(C)(C)C)ccc65)cc4)ccc2-3)cc1}	&	dbG2TAZ	&	\smiles{c1ccc2c(c1)ccc1[nH]c3ccc4ccccc4c3c12}	\\
SAF-CzB	&	\smiles{CC(C)(C)c1ccc2c(c1)c1cc(C(C)(C)C)ccc1n2-c1ccc(-c2ccc3c(c2)C2(c4ccccc4-3)c3ccccc3N(c3ccccc3)c3ccccc32)cc1}	&	V–TPA	&	\smiles{C=Cc1ccc(COc2cccc3c2c2ccccc2n3-c2ccc(N(c3ccc(-n4c5ccccc5c5ccccc54)cc3)c3ccc(-n4c5ccccc5c5ccccc54)cc3)cc2)cc1}	\\
BPy-TP2	&	\smiles{c1ccc(-c2ccc(-c3ccc4c5ccc(-c6ccc(-c7ccccn7)nc6)cc5c5ccccc5c4c3)cn2)nc1}	&	CCDC	&	\smiles{O=C(c1ccccc1)c1ccc(-n2c3ccc(-n4c5ccccc5c5ccccc54)cc3c3cc(-n4c5ccccc5c5ccccc54)ccc32)cc1}	\\
Cz–AQ	&	\smiles{O=C1c2ccccc2C(=O)c2cc(-n3c4ccccc4c4ccccc43)ccc21}	&	PXZPDO	&	\smiles{O=C(C=C(O)c1ccc(N2c3ccccc3Oc3ccccc32)cc1)c1ccc(N2c3ccccc3Oc3ccccc32)cc1}	\\
-ph-pCzAnBzt	&	\smiles{Cc1cc(-c2c3ccccc3c(-c3ccc(C\#N)cc3)c3ccccc23)c(C)cc1-c1ccc(-n2c3ccccc3c3ccccc32)cc1}	&	DMAC-BP	&	\smiles{CC1(C)c2ccccc2N(c2ccc(C(=O)c3ccc(N4c5ccccc5C(C)(C)c5ccccc54)cc3)cc2)c2ccccc21}	\\
PyI	&	\smiles{c1cc2ccc3cccc4c5[nH]cnc5c(c1)c2c34}	&	SBTF	&	\smiles{c1ccc2c(c1)-c1ccccc1C21c2c(ccc3cc(-c4ccc5ccccc5c4)ccc23)-c2ccc(-c3cc4ccccc4c4ccccc34)c3cccc1c23}	\\
DVCz-2CzCN	&	\smiles{C=Cc1ccc(COc2cccc3c2c2ccccc2n3CCCCCCOc2cccc3c2c2ccccc2n3-c2cccc(-n3c4ccccc4c4c(OCCCCCCn5c6ccccc6c6c(OCc7ccc(C=C)cc7)cccc65)cccc43)c2C\#N)cc1}	&	3-DPM	&	\smiles{N\#CC(C\#N)=C(c1cccc(N(c2ccccc2)c2ccccc2)c1)c1cccc(N(c2ccccc2)c2ccccc2)c1}	\\
ATDBT	&	\smiles{c1ccc(-c2c3ccccc3c(-c3ccc(-n4c5ccccc5c5c6sc7ccccc7c6ccc54)cc3)c3ccccc23)cc1}	&	4-DPM	&	\smiles{N\#CC(C\#N)=C(c1ccc(N(c2ccccc2)c2ccccc2)cc1)c1ccc(N(c2ccccc2)c2ccccc2)cc1}	\\
DPhBCz	&	\smiles{c1ccc(-n2c3ccccc3c3cc(-c4ccc5c(c4)c4ccccc4n5-c4ccccc4)ccc32)cc1}	&	2CzPy	&	\smiles{c1ccc(-c2ccc(-n3c4ccccc4c4ccccc43)cc2-n2c3ccccc3c3ccccc32)nc1}	\\
AcHPM	&	\smiles{CC1(C)c2ccccc2N(c2ccc(-c3cc(-c4ccc(N5c6ccccc6C(C)(C)c6ccccc65)cc4)ncn3)cc2)c2ccccc21}	&	2PBCzPy	&	\smiles{c1ccc(-n2c3ccccc3c3cc(-c4ccc5c(c4)c4ccccc4n5-c4ccc(-c5ccccn5)c(-n5c6ccccc6c6cc(-c7ccc8c(c7)c7ccccc7n8-c7ccccc7)ccc65)c4)ccc32)cc1}	\\
PhCO-DMAC	&	\smiles{CC1(C)c2ccccc2N(c2ccc(C(=O)c3ccccc3)cc2)c2ccccc21}	&	2BFCzPy	&	\smiles{c1ccc(-c2ccc(-n3c4ccccc4c4c5oc6ccccc6c5ccc43)cc2-n2c3ccccc3c3c4oc5ccccc5c4ccc32)nc1}	\\
Diphenylsulfane	&	\smiles{c1ccc(-c2cccc(-c3nc(-c4cccc(-c5ccccc5)c4)nc(-c4cccc(-c5ccccc5)c4)n3)c2)cc1}	&	CNTPA-CZ	&	\smiles{N\#Cc1ccc(N(c2ccccc2)c2cc(-n3c4ccccc4c4ccccc43)cc(-n3c4ccccc4c4ccccc43)c2)cc1}	\\
2BrCPT	&	\smiles{Brc1ccc2c(c1)c1cc(Br)ccc1n2-c1ccc(-c2nc(-c3ccccc3)nc(-c3ccccc3)n2)cc1}	&	CNTPA-PXZ	&	\smiles{N\#Cc1ccc(N(c2ccccc2)c2cc(N3c4ccccc4Oc4ccccc43)cc(N3c4ccccc4Oc4ccccc43)c2)cc1}	\\
DCbz-Ph-Bz	&	\smiles{c1ccc(-n2c(-c3ccc(-c4cc(-n5c6ccccc6c6ccccc65)cc(-n5c6ccccc6c6ccccc65)c4)cc3)nc3ccccc32)cc1}	&	BMZ-TZ	&	\smiles{c1cc(-c2nc3ccccc3[nH]2)cc(-c2nnnc(-c3cccc(-c4nc5ccccc5[nH]4)c3)c2-c2cccc(-c3nc4ccccc4[nH]3)c2)c1}	\\
DDPA-Ph-Bz	&	\smiles{c1ccc(N(c2ccccc2)c2cc(-c3ccc(-c4nc5ccccc5n4-c4ccccc4)cc3)cc(N(c3ccccc3)c3ccccc3)c2)cc1}	&	TPAPm	&	\smiles{N\#Cc1c(-c2ccc(N(c3ccccc3)c3ccccc3)cc2)ncnc1-c1ccc(N(c2ccccc2)c2ccccc2)cc1}	\\
IDAC-MCO	&	\smiles{Cc1cc2cc(-n3c4ccccc4c4cc5c(cc43)C(C)(C)c3ccccc3N5c3ccccc3)ccc2c(=O)o1}	&	CzPm	&	\smiles{N\#Cc1c(-c2ccc(-n3c4ccccc4c4ccccc43)cc2)ncnc1-c1ccc(-n2c3ccccc3c3ccccc32)cc1}	\\
diarylbiindole	&	\smiles{c1ccc(-c2c(-c3ccccc3)c3ccc4cccc5cc(-c6ccccc6)c(c2-c2ccccc2)c3c45)cc1}	&	TCPZ	&	\smiles{C1=C(c2cccc(-n3c4ccccc4c4ccccc43)c2)NN(c2cccc(-n3c4ccccc4c4ccccc43)c2)N=C1c1cccc(-n2c3ccccc3c3ccccc32)c1}	\\
CNDSB	&	\smiles{N\#CC(=Cc1ccccc1)c1ccc(C(C\#N)=Cc2ccccc2)cc1}	&	TrisCz-Trz	&	\smiles{c1ccc(-n2c3ccccc3c3ccccc32)c(-c2nc(-c3ccccc3-n3c4ccccc4c4ccccc43)nc(-c3ccccc3-n3c4ccccc4c4ccccc43)n2)c1}	\\
TMAB	&	\smiles{Cc1cc(C)c(CC\#Cc2ccc(N(C)C)cc2)c(B(c2c(C)c(C)cc(C)c2CC\#Cc2ccc(N(C)C)cc2)c2c(C)c(C)cc(C)c2CC\#Cc2ccc(N(C)C)cc2)c1C}	&	PTZ–TAZ	&	\smiles{c1ccc2c(c1)Oc1ccccc1N2c1ccc(-c2nnc(-c3ccc(N4c5ccccc5Oc5ccccc54)cc3)o2)cc1}	\\
MOF	&	\smiles{c1cncc(CC(c2nc3ccccc3[nH]2)N(C(Cc2cccnc2)c2nc3ccccc3[nH]2)C(Cc2cccnc2)c2nc3ccccc3[nH]2)c1}	&	DDMA-TXO2	&	\smiles{CC1(C)c2ccccc2N(c2ccc3c(c2)C(C)(C)c2cc(N4c5ccccc5C(C)(C)c5ccccc54)ccc2S3(=O)=O)c2ccccc21}	\\
DICzTrz	&	\smiles{N\#Cc1cc(-c2nc(-c3ccccc3)nc(-c3ccccc3)n2)ccc1-n1c2ccccc2c2cc3c4ccccc4n4c5ccccc5c(c21)c34}	&	UGH3	&	\smiles{c1ccc([Si](c2ccccc2)(c2ccccc2)c2cccc([Si](c3ccccc3)(c3ccccc3)c3ccccc3)c2)cc1}	\\
PXZ2PTO	&	\smiles{O=S1(=O)c2ccccc2N(c2ccccc2)c2cc(N3c4ccccc4Oc4ccccc43)ccc21}	&	&	\\
\end{longtable}
}


{
\begin{longtable}{l*{7}{c@{\extracolsep{\fill}}}}
\caption{Complete validation dataset for emission wavelength predictions. Predictions from
sTDA and sTD-DFT methods in gas and toluene phases are compared against
experimental/computational literature values. References are provided in the last column.}
\label{tab:supp-emission-complete}
\\\toprule
\textbf{Molecule} & \textbf{sTDA} & \textbf{sTD-DFT} & \textbf{sTDA} & \textbf{sTD-DFT} & \textbf{Ref.} & \textbf{Type} & \textbf{Citation} \\
 & \textbf{(Gas)} & \textbf{(Gas)} & \textbf{(Tol)} & \textbf{(Tol)} & \textbf{(nm)} &  & \\
\hline
\endfirsthead
\multicolumn{8}{c}{{\tablename\ \thetable{} -- continued from previous page}} \\
\hline
\textbf{Molecule} & \textbf{sTDA} & \textbf{sTD-DFT} & \textbf{sTDA} & \textbf{sTD-DFT} & \textbf{Ref.} & \textbf{Type} & \textbf{Citation} \\
 & \textbf{(Gas)} & \textbf{(Gas)} & \textbf{(Tol)} & \textbf{(Tol)} & \textbf{(nm)} &  & \\
\hline
\endhead
\hline
\multicolumn{8}{r}{{Continued on next page}} \\
\endfoot
\hline
\endlastfoot
-ph-pCzAnBzt & 501.4 & 512.6 & 467.6 & 477.7 & 450.0 & experimental &
\citenum{Guo_2021_10_1016_j_scib_2021_06_018} \\
1SBFN & 420.2 & 421.1 & 399.3 & 400.0 & 385.0 & experimental &
\citenum{Chen_2020_10_1039_d0tc03631e} \\
2,7-MeCzOXD & 524.7 & 524.2 & 507.0 & 506.3 & 520.0 & experimental & \citenum{Tang_2020}
\\
25CzBPym & 631.6 & 632.7 & 645.1 & 646.8 & 487.3 & experimental &
\citenum{Pandidurai_2021} \\
25tCzBPym & 643.1 & 644.4 & 682.5 & 684.2 & 506.0 & experimental &
\citenum{Pandidurai_2021} \\
2BPy-mDTC & 691.4 & 691.6 & 686.1 & 686.0 & 491.0 & experimental &
\citenum{Rajamalli_2017} \\
2C & 560.0 & 556.6 & 595.2 & 591.0 & 534.0 & experimental & \citenum{Barros_2023} \\
2CM & 384.9 & 391.7 & 366.5 & 372.7 & 374.0 & experimental & \citenum{Zassowski_2017} \\
2Cz-DMAC-BTB & 474.2 & 471.3 & 535.1 & 531.4 & 529.0 & experimental &
\citenum{Zhang_2022_10_1016_j_cclet_2021_08_064} \\
2Cz-DMAC-TTR & 469.5 & 472.9 & 517.6 & 522.9 & 601.0 & experimental &
\citenum{Zhang_2022_10_1016_j_cclet_2021_08_064} \\
2Cz-DMAC-TXO & 528.6 & 528.0 & 561.8 & 559.3 & 566.0 & experimental &
\citenum{Zhang_2022_10_1016_j_cclet_2021_08_064} \\
2Cz2tCzBn & 483.6 & 486.9 & 484.3 & 488.3 & 488.0 & experimental &
\citenum{Xie_2020_10_1039_d0tc00718h} \\
2CzPN & 474.8 & 475.8 & 490.5 & 492.1 & 480.4 & experimental & \citenum{Li_2018} \\
2DN & 775.6 & 773.4 & 736.4 & 751.2 & 659.4 & experimental & \citenum{Lin_2022} \\
2F & 776.6 & 788.7 & 700.4 & 711.6 & 493.3 & computational &
\citenum{Wang_2017_10_1039_c7cp02615c} \\
2PN & 823.0 & 825.1 & 1059.4 & 1062.0 & 673.3 & experimental & \citenum{Lin_2022} \\
2SBFN & 412.2 & 409.1 & 394.4 & 391.7 & 382.0 & experimental &
\citenum{Chen_2020_10_1039_d0tc03631e} \\
2TDBA-SBA & 447.6 & 454.2 & 448.2 & 454.8 & 448.0 & experimental & \citenum{Hua_2022} \\
2tCz2CzBn & 477.6 & 481.1 & 482.4 & 486.6 & 472.0 & experimental &
\citenum{Xie_2020_10_1039_d0tc00718h} \\
3,3-NHC & 369.4 & 370.2 & 426.0 & 428.7 & 412.7 & experimental & \citenum{Pramod_2019} \\
3,6-MeCzOXD & 513.1 & 512.5 & 500.5 & 499.9 & 520.0 & experimental & \citenum{Tang_2020}
\\
3,6-TC & 531.4 & 532.0 & 551.0 & 552.0 & 488.0 & experimental & \citenum{Tian_2021} \\
3-BHH & 358.9 & 362.0 & 339.8 & 342.9 & 478.5 & experimental &
\citenum{Pramod_2019_10_1016_j_saa_2019_117288} \\
3-DTPA-BBI & 631.2 & 632.6 & 675.2 & 676.6 & 604.0 & experimental & \citenum{Zhou_2022} \\
35CzBPym & 610.7 & 614.0 & 615.1 & 618.5 & 478.0 & experimental &
\citenum{Pandidurai_2021} \\
35tCzBPym & 622.7 & 625.9 & 647.0 & 650.6 & 496.0 & experimental &
\citenum{Pandidurai_2021} \\
3BPy-mDTC & 581.2 & 584.0 & 580.3 & 583.2 & 478.0 & experimental &
\citenum{Rajamalli_2017} \\
3CzBN & 438.3 & 437.9 & 437.5 & 437.4 & 431.0 & experimental &
\citenum{Zhang_2016_10_1039_c5sc04755b} \\
3CzFTFP & 402.8 & 407.6 & 430.0 & 435.6 & 484.0 & experimental & \citenum{Li_2018} \\
3DPTIA & 496.4 & 503.0 & 474.3 & 480.7 & 452.3 & experimental & \citenum{Yoo_2016} \\
3TCPM & 496.5 & 496.7 & 524.8 & 521.8 & 482.5 & experimental & \citenum{Ma_2021} \\
4CzIPN & 528.4 & 531.6 & 537.7 & 540.9 & 502.7 & experimental & \citenum{Sun_2020} \\
4NTAZ-PPI & 486.9 & 484.6 & 468.1 & 465.7 & 456.7 & experimental & \citenum{Cao_2021} \\
5CzICz & 418.8 & 423.4 & 409.3 & 413.9 & 410.0 & experimental &
\citenum{Zhao_2018_10_1039_c8tc03537g} \\
ACRXTN & 456.8 & 459.1 & 448.7 & 451.4 & 454.7 & experimental & \citenum{Wei_2022} \\
APDC-DTPA & 793.9 & 798.5 & 831.4 & 825.1 & 767.5 & experimental & \citenum{Xiao_2022} \\
AcCN & 562.0 & 557.5 & 571.8 & 567.2 & 512.0 & experimental & \citenum{Choi_2019} \\
BACF & 442.0 & 448.7 & 422.0 & 428.7 & 427.0 & experimental &
\citenum{Li_2022_10_1016_j_dyepig_2021_110030} \\
BACH & 433.6 & 439.5 & 409.9 & 415.7 & 427.7 & experimental &
\citenum{Li_2022_10_1016_j_dyepig_2021_110030} \\
BACN & 459.7 & 465.4 & 448.1 & 454.2 & 443.2 & experimental & \citenum{Qiu_2021} \\
BAPCN & 475.5 & 480.5 & 478.1 & 483.4 & 436.5 & experimental & \citenum{Qiu_2021} \\
BCN & 474.9 & 478.7 & 467.0 & 471.1 & 444.0 & experimental &
\citenum{Li_2022_10_1016_j_dyepig_2022_110456} \\
BDPCC & 437.3 & 440.3 & 425.7 & 428.4 & 483.5 & experimental & \citenum{Kumar_2022} \\
BDQ-tBuCz & 838.1 & 839.6 & 745.9 & 747.6 & 557.0 & experimental & \citenum{Yu_2017} \\
BDQDMAC & 962.9 & 962.8 & 868.4 & 868.3 & 601.7 & experimental & \citenum{Yu_2017} \\
BDQPCz & 820.6 & 814.9 & 721.4 & 715.1 & 502.0 & experimental & \citenum{Yu_2017} \\
BP-mDTC & 522.0 & 524.3 & 527.8 & 529.1 & 461.5 & experimental & \citenum{Rajamalli_2017}
\\
BSBFB & 398.7 & 398.1 & 376.9 & 376.6 & 388.0 & experimental &
\citenum{Chen_2020_10_1039_d0tc03631e} \\
BTZPP & 664.8 & 660.4 & 712.3 & 707.8 & 594.0 & experimental & \citenum{Zhou_2019} \\
BuCzCF3oB & 550.1 & 550.8 & 544.8 & 545.6 & 519.0 & experimental &
\citenum{Kumar_2020_10_1039_c9tc06204a} \\
C-H & 540.9 & 538.1 & 589.8 & 585.5 & 526.0 & experimental & \citenum{Barros_2023} \\
C3F7 & 498.1 & 498.2 & 472.2 & 471.2 & 558.0 & experimental &
\citenum{Kumar_2020_10_1039_c9tc06204a} \\
CBM & 505.0 & 504.9 & 501.4 & 501.4 & 458.0 & experimental & \citenum{Sharma_2020} \\
CDBP-BP-DMAC & 511.7 & 508.8 & 501.7 & 498.6 & 500.7 & experimental &
\citenum{Chen_2022_10_1039_d1qm01625c} \\
CDBP-BP-PXZ & 501.8 & 502.0 & 503.8 & 504.1 & 525.8 & experimental &
\citenum{Chen_2022_10_1039_d1qm01625c} \\
CN-QP & 859.4 & 854.1 & 945.1 & 936.2 & 602.0 & experimental & \citenum{Huang_2019} \\
CPzPC & 581.1 & 584.6 & 524.4 & 527.7 & 468.0 & experimental & \citenum{Wei_2014} \\
CPzPyC & 607.3 & 609.0 & 589.9 & 592.2 & 457.5 & experimental & \citenum{Wei_2014} \\
CPzPzC & 754.3 & 757.6 & 690.7 & 693.9 & 477.0 & experimental & \citenum{Wei_2014} \\
CTM & 473.7 & 477.1 & 485.6 & 483.0 & 502.7 & experimental & \citenum{Tan_2020} \\
CZ-DPS-DMAC & 415.3 & 418.7 & 437.2 & 436.7 & 458.7 & experimental & \citenum{Zheng_2022}
\\
Cz-SO & 368.3 & 372.9 & 367.3 & 371.7 & 424.0 & experimental &
\citenum{Liu_2021_10_1016_j_tet_2020_131869} \\
Cz2ICz & 446.3 & 450.8 & 436.7 & 441.1 & 420.0 & experimental &
\citenum{Zhao_2018_10_1039_c8tc03537g} \\
CzB-FMPIM & 492.0 & 496.3 & 479.7 & 484.5 & 420.0 & experimental &
\citenum{Jia_2019_10_1016_j_orgel_2019_03_044} \\
CzB-FMPPI & 502.9 & 504.4 & 491.9 & 492.9 & 432.0 & experimental &
\citenum{Jia_2019_10_1016_j_orgel_2019_03_044} \\
CzPh-PPO & 458.3 & 459.8 & 450.2 & 451.5 & 390.0 & experimental & \citenum{Xing_2017} \\
CzPrSBI & 473.5 & 479.4 & 444.5 & 450.9 & 477.3 & experimental & \citenum{Danyliv_2023} \\
CzTPA-m-Trz & 552.3 & 555.5 & 549.2 & 553.2 & 849.7 & experimental &
\citenum{Zhang_2020_10_1039_d0tc01614d} \\
CzoB & 497.9 & 498.0 & 472.3 & 471.4 & 517.0 & experimental &
\citenum{Kumar_2020_10_1039_c9tc06204a} \\
Cz–AQ & 727.9 & 732.7 & 710.7 & 715.7 & 586.5 & experimental & \citenum{Bin_Huang_2017} \\
DBQ-3PXZ & 851.7 & 854.2 & 868.0 & 870.1 & 618.0 & experimental & \citenum{Yu_2018} \\
DBT-BZ-DMAC & 636.7 & 634.8 & 665.7 & 664.3 & 535.5 & experimental & \citenum{Barman_2022}
\\
DCbz-Ph-Bz & 456.3 & 460.3 & 437.7 & 441.7 & 404.0 & experimental &
\citenum{Zhang_2020_10_1016_j_dyepig_2020_108477} \\
DCzBNPh-1 & 472.4 & 475.7 & 442.1 & 445.7 & 472.5 & experimental & \citenum{Yan_2022} \\
DCzIPN & 454.2 & 457.1 & 457.6 & 460.4 & 453.0 & experimental & \citenum{Kreiza_2021} \\
DDPA-Ph-Bz & 469.1 & 468.3 & 465.0 & 464.7 & 435.0 & experimental &
\citenum{Zhang_2020_10_1016_j_dyepig_2020_108477} \\
DDPhCz-DCPP & 755.3 & 757.7 & 739.8 & 742.3 & 571.3 & experimental & \citenum{Wang_2019}
\\
DFPFPhe & 503.0 & 506.8 & 475.0 & 478.4 & 537.0 & experimental & \citenum{Nagar_2022} \\
DHID-DBS & 461.5 & 467.0 & 488.2 & 492.2 & 478.0 & experimental & \citenum{Tongsuk_2021}
\\
DHID-DPS & 433.3 & 437.4 & 477.1 & 481.1 & 514.0 & experimental & \citenum{Tongsuk_2021}
\\
DMAC-11-DPPZ & 1031.3 & 1030.1 & 891.9 & 891.3 & 577.5 & experimental &
\citenum{Zhou_2020} \\
DMAC-BPP & 641.0 & 636.9 & 599.3 & 594.8 & 505.0 & experimental & \citenum{Duan_2022} \\
DMAC-Cz-TTR & 500.1 & 503.8 & 581.4 & 584.4 & 449.0 & experimental &
\citenum{Zhang_2021_10_1016_j_orgel_2021_106312} \\
DMAC-DPS & 729.4 & 729.1 & 628.2 & 628.1 & 470.0 & experimental & \citenum{Luo_2016} \\
DMAC-TRZ & 619.5 & 616.8 & 602.9 & 602.4 & 500.6 & theoretical study &
\citenum{Dhali_2021} \\
DMACMNPTO & 348.6 & 350.2 & 346.3 & 347.9 & 382.0 & experimental & \citenum{Guo_2020} \\
DMeCzIPN & 479.2 & 481.5 & 486.4 & 488.8 & 474.0 & experimental & \citenum{Kreiza_2021} \\
DNFPh & 540.6 & 544.7 & 468.9 & 470.1 & 521.5 & experimental & \citenum{Nagar_2022} \\
DPACPhTPI & 495.3 & 496.3 & 483.9 & 485.0 & 436.0 & experimental & \citenum{Huang_2017} \\
DPACpBr & 367.6 & 367.9 & 365.0 & 365.9 & 448.7 & experimental & \citenum{Lee_2021} \\
DPCN & 511.7 & 511.6 & 508.0 & 508.3 & 424.0 & experimental &
\citenum{Xiao_2022_10_1016_j_dyepig_2022_110451} \\
DPFPhe & 494.3 & 494.7 & 461.9 & 464.2 & 532.5 & experimental & \citenum{Nagar_2022} \\
DPPP & 532.0 & 535.3 & 486.6 & 489.5 & 390.7 & experimental &
\citenum{Chen_2020_10_1039_d0tc03631e} \\
DPPSPF & 405.3 & 406.4 & 393.8 & 395.0 & 385.0 & experimental & \citenum{Di_2020} \\
DPS-PXZ & 485.1 & 485.2 & 464.8 & 465.0 & 537.0 & experimental & \citenum{Leng_2020} \\
DPXZ-DPPM & 1090.2 & 1093.3 & 1048.7 & 1051.3 & 636.0 & experimental &
\citenum{Liang_2020} \\
DPhBCz & 395.1 & 399.7 & 381.8 & 386.2 & 365.0 & experimental & \citenum{Ou_2016} \\
DQP & 1207.5 & 1183.6 & 1057.9 & 1036.8 & 562.0 & experimental & \citenum{Barman_2022} \\
DTPCZTZ & 468.9 & 473.7 & 455.8 & 460.7 & 474.6 & experimental & \citenum{Xu_2023} \\
DV-3CzCN & 470.6 & 472.0 & 442.0 & 442.6 & 473.6 & experimental &
\citenum{Sun_2022_10_1016_j_dyepig_2021_109967} \\
DVCz-2CzCN & 434.5 & 438.3 & 453.4 & 454.9 & 461.5 & experimental &
\citenum{Sun_2022_10_1016_j_orgel_2022_106442} \\
Diphenylsulfane & 541.6 & 535.9 & 511.4 & 507.5 & 454.0 & experimental &
\citenum{Cao_2020} \\
F2AcBO & 491.3 & 489.9 & 515.6 & 516.6 & 476.7 & experimental &
\citenum{Chen_2022_10_1016_j_cej_2021_133078} \\
F2PA & 435.9 & 435.7 & 426.3 & 426.2 & 403.3 & experimental &
\citenum{Chen_2020_10_1039_d0tc03631e} \\
FTAT-FBO & 399.3 & 398.4 & 375.0 & 374.2 & 459.0 & experimental &
\citenum{Liu_2022_10_1016_j_cej_2022_137372} \\
FlCz & 553.3 & 548.0 & 525.0 & 520.0 & 412.0 & experimental & \citenum{Zhen_2022} \\
IABTCN & 659.1 & 658.8 & 694.6 & 696.9 & 650.0 & experimental & \citenum{Kim_2019} \\
ImPy-1 & 463.0 & 469.9 & 437.0 & 443.5 & 479.0 & experimental &
\citenum{Girase_2022_10_1039_d2nj01938h} \\
ImPy-2 & 465.7 & 472.2 & 441.1 & 447.4 & 471.0 & experimental &
\citenum{Girase_2022_10_1039_d2nj01938h} \\
ImPy-3 & 557.2 & 556.0 & 533.0 & 531.4 & 480.0 & experimental &
\citenum{Girase_2022_10_1039_d2nj01938h} \\
InDM & 416.2 & 422.6 & 399.3 & 405.0 & 490.0 & experimental & \citenum{Umasankar_2021} \\
MBF & 605.6 & 604.9 & 662.7 & 665.5 & 614.0 & experimental & \citenum{Zhan_2021} \\
MPA & 553.7 & 554.2 & 560.5 & 561.2 & 574.7 & experimental & \citenum{Zhan_2021} \\
MPPA & 805.6 & 809.3 & 941.6 & 947.7 & 679.1 & experimental & \citenum{Sharma_2022} \\
Me-DMAC & 527.8 & 524.1 & 543.5 & 539.4 & 517.0 & experimental &
\citenum{Li_2020_10_1039_c9tc05340a} \\
Me2ACBO & 498.1 & 499.9 & 509.1 & 508.8 & 474.0 & experimental &
\citenum{Chen_2022_10_1016_j_cej_2021_133078} \\
NA-TNA & 648.4 & 645.2 & 683.6 & 680.0 & 619.2 & experimental & \citenum{Zheng_2015} \\
NAI–BiFA & 558.5 & 562.8 & 573.2 & 578.7 & 596.1 & experimental & \citenum{Chen_2019} \\
NAI–PhBiFA & 569.8 & 566.6 & 647.7 & 645.0 & 556.8 & experimental & \citenum{Chen_2019} \\
NAP-1,5-DPA & 501.4 & 502.3 & 477.3 & 478.1 & 507.2 & experimental & \citenum{Dumur_2016}
\\
NAQ3 & 688.2 & 692.4 & 685.2 & 689.2 & 530.0 & experimental & \citenum{Fan_2019} \\
OCzBN & 449.2 & 453.1 & 434.7 & 436.7 & 393.0 & experimental & \citenum{Ban_2019} \\
P4CzCN–BCz & 387.1 & 383.9 & 334.2 & 331.4 & 426.5 & experimental &
\citenum{Zhou_2022_10_1039_d1tc05010a} \\
PBTPA & 736.2 & 732.1 & 753.7 & 750.2 & 642.5 & experimental & \citenum{Yu_2021} \\
PCTXO & 561.6 & 557.6 & 568.3 & 564.3 & 616.4 & experimental & \citenum{Khammultri_2021}
\\
PDA-DP & 477.9 & 477.5 & 472.4 & 472.2 & 480.0 & experimental & \citenum{Gong_2017} \\
PIAnTPh & 498.9 & 506.6 & 462.9 & 470.2 & 440.0 & experimental &
\citenum{Liu_2021_10_1016_j_cej_2021_131351} \\
PO-ABN & 509.7 & 519.7 & 472.9 & 482.2 & 434.0 & experimental & \citenum{Kang_2020} \\
POBBPE & 362.1 & 364.1 & 341.9 & 344.0 & 320.0 & experimental & \citenum{Jeon_2018} \\
PS-BZ-DMAC & 474.2 & 478.9 & 477.1 & 482.3 & 574.0 & experimental &
\citenum{Ma_2021_10_1016_j_dyepig_2021_109649} \\
PSPSF & 406.3 & 407.2 & 394.8 & 395.6 & 386.0 & experimental & \citenum{Di_2020} \\
PT-BZ-DMAC & 573.1 & 568.5 & 583.9 & 580.7 & 531.8 & experimental &
\citenum{Ma_2021_10_1016_j_dyepig_2021_109649} \\
PTZ-2PTO & 459.1 & 460.2 & 434.2 & 434.2 & 452.2 & experimental & \citenum{Xiang_2019} \\
PTZ-DBTO2 & 473.4 & 479.6 & 471.4 & 477.9 & 520.0 & experimental & \citenum{Duan_2022} \\
PX-TRZ & 541.9 & 543.5 & 554.0 & 554.6 & 531.0 & experimental & \citenum{Qu_2020} \\
PXZ-10-DPPZ & 1099.7 & 1102.8 & 958.6 & 960.6 & 614.0 & experimental & \citenum{Zhou_2020}
\\
PXZ-11-DPPZ & 987.5 & 991.7 & 863.8 & 867.0 & 628.5 & experimental & \citenum{Zhou_2020}
\\
PXZ2PTO & 430.0 & 430.4 & 420.1 & 420.5 & 496.7 & experimental &
\citenum{Xiang_2018_10_1039_c8tc03648a} \\
PmPmP & 653.0 & 653.8 & 643.5 & 644.3 & 563.0 & experimental &
\citenum{Shi_2022_10_1039_d1tc03803f} \\
PnB & 518.0 & 514.4 & 488.4 & 484.9 & 465.0 & experimental & \citenum{Umasankar_2021} \\
Pra-2DMAC & 370.2 & 371.9 & 359.0 & 360.5 & 513.0 & experimental &
\citenum{Zhang_2021_10_1039_d0tc04315j} \\
PtB & 527.9 & 525.4 & 487.3 & 483.6 & 446.3 & experimental & \citenum{Umasankar_2021} \\
PxCN & 558.1 & 558.2 & 576.1 & 576.2 & 527.5 & experimental & \citenum{Choi_2019} \\
PyDCN & 423.8 & 422.9 & 403.4 & 402.5 & 472.0 & experimental & \citenum{Dong_2022} \\
PyDCN–DMAC & 508.7 & 504.4 & 496.9 & 493.2 & 494.2 & experimental & \citenum{Dong_2022} \\
PyIAnTPh & 509.5 & 519.5 & 470.9 & 480.2 & 572.5 & experimental &
\citenum{Liu_2021_10_1016_j_cej_2021_131351} \\
SBA & 351.0 & 352.5 & 342.3 & 343.8 & 350.0 & experimental & \citenum{Zeng_2021} \\
SBDBQ-PXZ & 843.2 & 845.7 & 751.8 & 753.8 & 594.0 & experimental & \citenum{Yu_2018} \\
SF-DPSO & 431.1 & 431.1 & 384.5 & 384.4 & 399.5 & experimental & \citenum{Ban_2014} \\
SFAC-BP-SFAC & 592.9 & 592.6 & 597.8 & 597.2 & 506.0 & experimental & \citenum{Huang_2022}
\\
SFCz & 556.1 & 553.2 & 529.4 & 526.6 & 413.0 & experimental & \citenum{Zhen_2022} \\
STAC-BP-STAC & 582.5 & 582.1 & 575.3 & 575.0 & 506.5 & experimental & \citenum{Huang_2022}
\\
SXAC-BP-SXAC & 585.8 & 585.4 & 584.9 & 584.4 & 500.0 & experimental & \citenum{Huang_2022}
\\
SiCz1Py3 & 405.9 & 410.7 & 387.3 & 392.0 & 479.0 & experimental & \citenum{Choi_2017} \\
SiCz2Py2 & 399.8 & 404.0 & 389.3 & 393.7 & 478.0 & experimental & \citenum{Choi_2017} \\
SiCz3Py1 & 409.9 & 414.1 & 386.3 & 390.5 & 476.0 & experimental & \citenum{Choi_2017} \\
SpiroAC-TRZ & 610.1 & 610.1 & 591.5 & 591.0 & 484.5 & experimental & \citenum{Barman_2022}
\\
TATP-BP & 558.1 & 557.8 & 568.1 & 566.4 & 527.7 & experimental &
\citenum{Chen_2018_10_1039_c8tc04721a} \\
TAZ-PPI & 499.8 & 500.6 & 487.4 & 488.5 & 456.7 & experimental & \citenum{Cao_2021} \\
TBN-TPA & 436.7 & 443.9 & 423.1 & 430.0 & 453.0 & experimental & \citenum{Meng_2022} \\
TBPe & 542.6 & 553.7 & 500.3 & 510.7 & 471.0 & experimental & \citenum{Yao_2019} \\
TBRb & 677.3 & 696.5 & 610.9 & 628.2 & 538.0 & experimental & \citenum{Yao_2019} \\
TCz-3PA-TCz & 580.2 & 578.6 & 534.0 & 532.6 & 452.0 & experimental & \citenum{Wei_2019} \\
TDBA-SBA & 422.5 & 428.8 & 422.4 & 428.7 & 464.0 & experimental & \citenum{Hua_2022} \\
TFM-QP & 857.2 & 857.2 & 929.3 & 923.2 & 566.5 & experimental & \citenum{Huang_2019} \\
TMAB & 434.0 & 432.4 & 407.6 & 413.6 & 452.0 & experimental &
\citenum{Berrones_Reyes_2018} \\
TPA-DCPP & 765.0 & 762.7 & 781.4 & 779.2 & 669.3 & experimental & \citenum{Barman_2022} \\
TPA-L-BN & 451.4 & 452.9 & 391.4 & 392.8 & 394.0 & experimental & \citenum{Wu_2019} \\
TPA-PPO & 493.0 & 490.9 & 500.0 & 497.8 & 472.0 & experimental & \citenum{Xing_2017} \\
TPA-PyF2 & 390.2 & 390.8 & 379.4 & 381.1 & 376.0 & experimental & \citenum{Gao_2022} \\
TPADSO2 & 450.5 & 452.0 & 480.0 & 482.0 & 533.3 & experimental &
\citenum{Sun_2017_10_1016_j_dyepig_2016_12_037} \\
TPB-AC & 469.2 & 469.4 & 475.9 & 476.0 & 370.0 & experimental &
\citenum{Lin_2022_10_1039_d2mh00627h} \\
TPBPPI-PBI & 472.9 & 473.4 & 461.4 & 462.6 & 429.0 & experimental & \citenum{Zhu_2020} \\
TPBPPI-PY & 472.2 & 472.0 & 459.4 & 459.9 & 420.0 & experimental & \citenum{Zhu_2020} \\
TPXZ-as-TAZ & 888.7 & 887.8 & 889.5 & 891.5 & 529.0 & computational & \citenum{Zhang_2018}
\\
TRZ-3SO2 & 577.4 & 572.9 & 540.6 & 539.9 & 706.0 & experimental &
\citenum{Zhang_2022_10_1039_d2tc02694e} \\
TSBFB & 404.5 & 403.6 & 385.1 & 384.4 & 395.0 & experimental &
\citenum{Chen_2020_10_1039_d0tc03631e} \\
TXO-TPA & 570.6 & 566.0 & 598.6 & 595.7 & 603.0 & experimental & \citenum{Khammultri_2021}
\\
VBN0 & 432.7 & 433.5 & 419.5 & 421.1 & 482.3 & experimental & \citenum{Sun_2022} \\
XT-T & 360.0 & 360.0 & 349.2 & 350.5 & 478.0 & experimental & \citenum{Fu_2022} \\
{[2,1-b]IF} & 1031.9 & 586.3 & 770.2 & 571.9 & 347.0 & experimental &
\citenum{Romain_2016}
\\
dTolmp & 435.6 & 434.6 & 397.2 & 396.9 & 548.0 & experimental &
\citenum{Wang_2021_10_1039_d1tc00678a} \\
iTPBI-CN & 441.4 & 440.7 & 432.7 & 432.1 & 443.5 & experimental & \citenum{Wang_2015} \\
m-CF3PIBI & 492.2 & 492.6 & 485.5 & 486.1 & 435.3 & experimental & \citenum{Girase_2022}
\\
m-CZ-DPS-DMAC & 425.1 & 428.5 & 423.9 & 427.9 & 490.8 & experimental & \citenum{Dong_2023}
\\
m-DTPACO & 484.6 & 485.3 & 501.6 & 502.8 & 375.0 & experimental & \citenum{Hu_2017} \\
m-PCzTPD & 522.4 & 522.9 & 522.1 & 519.4 & 506.0 & experimental & \citenum{Kautny_2017} \\
m-PO-ABN & 509.5 & 519.6 & 472.0 & 481.3 & 450.0 & experimental & \citenum{Kang_2020} \\
mCP-BP-DMAC & 503.7 & 504.0 & 504.6 & 505.0 & 505.5 & experimental & \citenum{Wu_2023} \\
mCP-BmPy & 395.9 & 399.3 & 384.9 & 389.0 & 410.0 & experimental &
\citenum{Li_2020_10_1016_j_cej_2020_125276} \\
mCP-Ph & 390.6 & 394.3 & 371.4 & 375.4 & 350.0 & experimental &
\citenum{Li_2020_10_1016_j_cej_2020_125276} \\
mCP-Py & 385.2 & 388.5 & 372.2 & 376.2 & 350.0 & experimental &
\citenum{Li_2020_10_1016_j_cej_2020_125276} \\
mCz-TAn-CN & 519.1 & 529.1 & 483.8 & 492.6 & 438.0 & experimental &
\citenum{Wu_2021_10_1039_d1qm00880c} \\
mP2MPC & 450.6 & 452.1 & 430.5 & 432.5 & 399.0 & experimental & \citenum{Peng_2022} \\
mTPA-PPI & 457.5 & 457.9 & 443.8 & 445.0 & 404.0 & experimental &
\citenum{Chen_2020_10_1039_d0tc03631e} \\
o-CzOXD & 500.3 & 500.5 & 477.5 & 477.5 & 492.7 & experimental & \citenum{Tang_2020} \\
o-PCzTPD & 558.4 & 555.3 & 518.5 & 519.1 & 485.0 & experimental & \citenum{Kautny_2017} \\
oCBP & 388.3 & 392.6 & 377.0 & 381.2 & 389.0 & experimental & \citenum{Ha_2022} \\
p-DTPACO & 508.7 & 510.4 & 531.9 & 533.9 & 540.0 & experimental & \citenum{Hu_2017} \\
p-PCzTPD & 555.3 & 553.3 & 563.4 & 561.2 & 476.0 & experimental & \citenum{Kautny_2017} \\
p-PO-ABN & 514.8 & 521.5 & 488.0 & 494.6 & 456.0 & experimental & \citenum{Kang_2020} \\
p-PO15NCzDPA & 508.1 & 518.1 & 467.2 & 477.0 & 427.0 & experimental & \citenum{Ye_2020} \\
pSFIAc2 & 470.9 & 479.8 & 461.4 & 469.7 & 451.5 & experimental & \citenum{Meng_2022} \\
pzpy & 510.1 & 519.9 & 466.8 & 476.8 & 458.0 & experimental &
\citenum{Zhang_2019_10_1039_c9tc00465c} \\
t-DABNA & 431.1 & 438.5 & 418.4 & 425.5 & 464.0 & experimental &
\citenum{Lee_2019_10_1039_c9tc02746g} \\
t3Cz-SO & 434.5 & 440.1 & 442.9 & 448.2 & 446.0 & experimental &
\citenum{Liu_2021_10_1016_j_tet_2020_131869} \\
tBuCzDBA & 757.8 & 751.2 & 684.8 & 678.9 & 558.5 & experimental &
\citenum{Liu_2023_10_1016_j_orgel_2022_106720} \\
tCTM & 483.8 & 487.0 & 483.3 & 483.5 & 459.0 & experimental & \citenum{Tan_2020} \\
thio-ether & 1032.4 & 586.4 & 770.0 & 572.2 & 366.0 & experimental & \citenum{Romain_2016}
\\
t–BuCz2BP & 488.8 & 490.6 & 490.8 & 490.3 & 440.0 & experimental & \citenum{Huang_2016} \\
$\alpha$-2DMAC-DBP & 1177.9 & 1176.8 & 1039.5 & 1038.2 & 637.5 & experimental &
\citenum{Maeng_2021} \\
$\alpha$-DMAC-DBP & 1092.4 & 1092.0 & 952.2 & 951.8 & 622.0 & experimental &
\citenum{Maeng_2021} \\
$\beta$-DI-DBP & 961.4 & 952.3 & 832.6 & 825.2 & 582.5 & experimental &
\citenum{Maeng_2021}
\\
\end{longtable}
}

 
{
\begin{longtable}{l*{7}{c@{\extracolsep{\fill}}}}
\caption{Complete validation dataset for singlet-triplet gap predictions. Predictions from
sTDA and sTD-DFT methods in gas and toluene phases are compared against
experimental/computational literature values. References are provided in the last column.}
\label{tab:supp-st-complete}
\\\toprule
\textbf{Molecule} & \textbf{sTDA} & \textbf{sTD-DFT} & \textbf{sTDA} & \textbf{sTD-DFT} & \textbf{Ref.} & \textbf{Type} & \textbf{Citation} \\
 & \textbf{(Gas)} & \textbf{(Gas)} & \textbf{(Tol)} & \textbf{(Tol)} & \textbf{(eV)} &  & \\
\hline
\endfirsthead
\multicolumn{8}{c}{{\tablename\ \thetable{} -- continued from previous page}} \\
\hline
\textbf{Molecule} & \textbf{sTDA} & \textbf{sTD-DFT} & \textbf{sTDA} & \textbf{sTD-DFT} & \textbf{Ref.} & \textbf{Type} & \textbf{Citation} \\
 & \textbf{(Gas)} & \textbf{(Gas)} & \textbf{(Tol)} & \textbf{(Tol)} & \textbf{(eV)} &  & \\
\hline
\endhead
\hline
\multicolumn{8}{r}{{Continued on next page}} \\
\endfoot
\hline
\endlastfoot
11-AcBPdCN & 0.3009 & 0.2809 & 0.3667 & 0.3437 & 0.0650 & experimental &
\citenum{Kothavale_2022} \\
12AcCz-PM & 0.2661 & 0.2631 & 0.3416 & 0.3407 & 0.3900 & experimental &
\citenum{Zhang_2019_10_1039_c9tc04284a} \\
12BTCzTPN & 0.1955 & 0.1766 & 0.1720 & 0.1610 & 0.0829 & experimental &
\citenum{Shin_2022_10_1016_j_cej_2022_137181} \\
1CM & 0.5074 & 0.4471 & 0.4934 & 0.4316 & 0.2900 & experimental & \citenum{Zassowski_2017}
\\
1CzSO & 0.3195 & 0.2961 & 0.2972 & 0.2779 & 0.2150 & experimental &
\citenum{Li_2017_10_1039_c7tc00119c} \\
1PXZP & 0.1653 & 0.1764 & 0.1692 & 0.1646 & 0.2600 & computational &
\citenum{Wang_2020_10_1016_j_dyepig_2020_108437} \\
23AcCz-PM & 0.2003 & 0.2120 & 0.2544 & 0.2647 & 0.0600 & experimental &
\citenum{Zhang_2019_10_1039_c9tc04284a} \\
25CzBPym & 0.0833 & 0.0797 & 0.0929 & 0.0879 & 0.0604 & experimental &
\citenum{Pandidurai_2021} \\
25DAcBPy & 0.0576 & 0.0575 & 0.3301 & 0.3383 & 0.0235 & computational &
\citenum{Pandidurai_2019} \\
25tCzBPym & 0.0848 & 0.0811 & 0.0917 & 0.0872 & 0.0700 & experimental &
\citenum{Pandidurai_2021} \\
26DAcBPy & 0.1948 & 0.1907 & 0.1799 & 0.1764 & 0.0070 & computational &
\citenum{Pandidurai_2019} \\
26DPXZBPy & 0.2012 & 0.1987 & 0.1685 & 0.1661 & 0.0520 & computational &
\citenum{Pandidurai_2019} \\
2BPy-mDTC & 0.0798 & 0.0793 & 0.0847 & 0.0849 & 0.0650 & experimental &
\citenum{Rajamalli_2017} \\
2BrCPT & 0.2696 & 0.2431 & 0.3015 & 0.2718 & 0.0400 & experimental &
\citenum{Huang_2021_10_1016_j_mtener_2021_100705} \\
2CM & 0.4735 & 0.4172 & 0.4762 & 0.4198 & 0.3100 & experimental & \citenum{Zassowski_2017}
\\
2Cz-DMAC-BTB & 0.2665 & 0.2823 & 0.1276 & 0.1438 & 0.0200 & experimental &
\citenum{Zhang_2022_10_1016_j_cclet_2021_08_064} \\
2Cz-DMAC-TTR & 0.2933 & 0.2743 & 0.3442 & 0.3199 & 0.0470 & experimental &
\citenum{Zhang_2022_10_1016_j_cclet_2021_08_064} \\
2Cz-DMAC-TXO & 0.3679 & 0.3706 & 0.3220 & 0.3319 & 0.0360 & experimental &
\citenum{Zhang_2022_10_1016_j_cclet_2021_08_064} \\
2Cz2tCzBn & 0.1908 & 0.1730 & 0.2168 & 0.1959 & 0.1300 & experimental &
\citenum{Xie_2020_10_1039_d0tc00718h} \\
2CzPN & 0.2427 & 0.2373 & 0.2677 & 0.2593 & 0.2525 & experimental & \citenum{Kumar_2022}
\\
2CzTPEPCz & 0.3096 & 0.3171 & 0.3311 & 0.3352 & 0.1300 & experimental &
\citenum{Zhao_2018} \\
2DMAC-BP-F & 0.2632 & 0.2475 & 0.3177 & 0.2988 & 0.1367 & experimental &
\citenum{Hong_2022_10_1039_d1tc04918f} \\
2DN & 0.0920 & 0.0965 & 0.4000 & 0.3669 & 0.0700 & experimental & \citenum{Lin_2022} \\
2F & 0.2152 & 0.1908 & 0.2374 & 0.2097 & 0.2300 & computational &
\citenum{Wang_2017_10_1039_c7cp02615c} \\
2H-Qz & 0.2138 & 0.2118 & 0.2326 & 0.2307 & 0.1100 & experimental &
\citenum{Li_2021_10_1039_d1tc02633j} \\
2Mi2SB & 0.3464 & 0.3356 & 0.3865 & 0.3721 & 0.0300 & computational &
\citenum{Matsui_2022} \\
2PCzBN & 0.3041 & 0.3010 & 0.3170 & 0.3108 & 0.2800 & experimental & \citenum{Hong_2022}
\\
2PCzBN-tPh & 0.2207 & 0.2191 & 0.2414 & 0.2386 & 0.2600 & experimental &
\citenum{Hong_2022} \\
2PN & 0.1324 & 0.1285 & 0.1074 & 0.1046 & 0.0100 & experimental & \citenum{Lin_2022} \\
2PS-2FPh & 0.2729 & 0.2650 & 0.3064 & 0.2983 & 0.2050 & computational &
\citenum{Skhirtladze_2023} \\
2PXZ-OXD & 0.2031 & 0.2024 & 0.1787 & 0.1783 & 0.2300 & experimental & \citenum{Lee_2013}
\\
2PXZP & 0.2176 & 0.2301 & 0.2422 & 0.2410 & 0.1340 & computational &
\citenum{Wang_2020_10_1016_j_dyepig_2020_108437} \\
2SPBAC-BP & 0.4008 & 0.4004 & 0.4024 & 0.3933 & 0.2200 & experimental &
\citenum{Song_2022} \\
2tCz2CzBn & 0.1995 & 0.1807 & 0.2266 & 0.2046 & 0.1300 & experimental &
\citenum{Xie_2020_10_1039_d0tc00718h} \\
3-ACBPmCN & 0.0470 & 0.0582 & 0.0519 & 0.0651 & 0.1800 & computational &
\citenum{Kothavale_2022} \\
3-AcBPdCN & 0.2112 & 0.2004 & 0.2667 & 0.2540 & 0.1800 & experimental &
\citenum{Kothavale_2022} \\
3-AcBPtCN & 0.0637 & 0.0643 & 0.2413 & 0.2319 & 0.1300 & experimental &
\citenum{Kothavale_2022} \\
3-DPM & 0.1689 & 0.1581 & 0.1558 & 0.1492 & 0.0987 & experimental & \citenum{Sohn_2017} \\
35CzBPym & 0.1491 & 0.1382 & 0.1503 & 0.1393 & 0.1100 & experimental &
\citenum{Pandidurai_2021} \\
35tCzBPym & 0.1484 & 0.1383 & 0.1537 & 0.1431 & 0.1100 & experimental &
\citenum{Pandidurai_2021} \\
36DTPAFM & 0.2073 & 0.2235 & 0.2241 & 0.2391 & 0.1200 & experimental &
\citenum{Li_2022_10_1016_j_orgel_2022_106645} \\
3BPy-mDTC & 0.1555 & 0.1449 & 0.1562 & 0.1456 & 0.1200 & experimental &
\citenum{Rajamalli_2017} \\
3CzBN & 0.3294 & 0.3321 & 0.3069 & 0.3081 & 0.3100 & experimental &
\citenum{Zhang_2016_10_1039_c5sc04755b} \\
3CzPhpPM & 0.2409 & 0.2185 & 0.2659 & 0.2406 & 0.3300 & experimental &
\citenum{Lv_2020_10_1039_d0tc00039f} \\
3DMAC-BP-Br & 0.0203 & 0.0210 & 0.2215 & 0.2151 & 0.0400 & experimental &
\citenum{Huang_2021} \\
3DMAC-BP-CN & 0.1852 & 0.1962 & 0.2209 & 0.2142 & 0.0200 & experimental &
\citenum{Huang_2021} \\
3DPA3CN & 0.2230 & 0.2266 & 0.2186 & 0.2212 & 0.1030 & experimental &
\citenum{Taneda_2015} \\
3Py-DMAC & 0.0653 & 0.0776 & 0.3069 & 0.2963 & 0.0400 & experimental &
\citenum{Liu_2020_10_1016_j_cej_2020_126107} \\
3TCPM & 0.2621 & 0.2614 & 0.1934 & 0.2072 & 0.2600 & experimental & \citenum{Ma_2021} \\
3ai & 0.8060 & 0.7332 & 0.8521 & 0.7687 & 0.1800 & experimental & \citenum{Puntscher_2015}
\\
4-Ac & 0.5593 & 0.5465 & 0.6006 & 0.5815 & 0.0800 & experimental &
\citenum{Keruckiene_2020_10_1016_j_orgel_2019_105576} \\
4-DPM & 0.4478 & 0.4605 & 0.4672 & 0.4791 & 0.1427 & experimental & \citenum{Sohn_2017} \\
4-SCZ & 0.4456 & 0.4094 & 0.4687 & 0.4303 & 0.2400 & experimental &
\citenum{Zhu_2019_10_1016_j_jlumin_2018_10_017} \\
4BPy-mDTC & 0.1541 & 0.1431 & 0.1477 & 0.1375 & 0.1300 & experimental &
\citenum{Rajamalli_2017} \\
4CzIPN & 0.1507 & 0.1367 & 0.1624 & 0.1487 & 0.1103 & experimental &
\citenum{Zhang_2016_10_1039_c6tc00214e} \\
4CzTPN & 0.2100 & 0.1874 & 0.2339 & 0.2079 & 0.0800 & experimental &
\citenum{Zhao_2022_10_1039_d2tc00238h} \\
4F-m-$\nu$-DABNA & 0.2638 & 0.2488 & 0.2657 & 0.2515 & 0.0700 & experimental &
\citenum{Rayappa_Naveen_2022} \\
4F-$\nu$-DABNA & 0.2709 & 0.2560 & 0.2770 & 0.2625 & 0.0500 & experimental &
\citenum{Rayappa_Naveen_2022} \\
5,6PXZ-PIDO & 0.2119 & 0.2186 & 0.1857 & 0.1851 & 0.4500 & computational &
\citenum{Xiang_2018_10_1039_c8tc01656a} \\
5-Cz-Sac & 0.3410 & 0.3079 & 0.3315 & 0.3064 & 0.2800 & experimental & \citenum{Xiao_2021}
\\
6,7-DCNQx-DICz & 0.0088 & 0.0143 & 0.0207 & 0.0279 & 0.0276 & computational &
\citenum{Kothavale_2021} \\
6AcBIQ & 0.0859 & 0.0897 & 0.1004 & 0.1012 & 0.1000 & experimental & \citenum{Yun_2017} \\
9CzFDBFSPO & 0.3962 & 0.3615 & 0.3951 & 0.3550 & 0.1200 & experimental &
\citenum{Zhang_2020}
\\
ACID–BPPZ & 0.2041 & 0.2014 & 0.1581 & 0.1669 & 0.0150 & computational &
\citenum{Chen_2019_10_1039_c8tc06081a} \\
ACRFLCN & 0.0002 & 0.0047 & 0.0001 & 0.0002 & 0.1000 & computational &
\citenum{Zhu_2020_10_1039_d0tc00743a} \\
ACRSA & 0.0025 & 0.0217 & -0.0001 & 0.0000 & 0.0350 & experimental & \citenum{Nasu_2013}
\\
ACRXTN & 0.4781 & 0.4642 & 0.5345 & 0.5175 & 0.1730 & experimental & \citenum{Wei_2022} \\
AI-4Cz & 0.2055 & 0.1909 & 0.2287 & 0.2140 & 0.0200 & experimental &
\citenum{Liu_2021_10_1016_j_orgel_2020_106017} \\
ANQDC & 0.2365 & 0.2634 & 0.2815 & 0.3132 & 0.0200 & computational & \citenum{Gong_2021}
\\
ANQDC-PSTA & 0.0965 & 0.0948 & 0.1344 & 0.1296 & 0.0800 & experimental &
\citenum{Gong_2021}
\\
APDC-DTPA & 0.2802 & 0.2713 & 0.2891 & 0.3004 & 0.1400 & experimental &
\citenum{Chen_2022}
\\
APPT & 0.1857 & 0.1995 & 0.2290 & 0.2449 & 0.0700 & experimental &
\citenum{Yang_2021_10_1016_j_dyepig_2021_109474} \\
APPT-PXZ & 0.1829 & 0.1755 & 0.2201 & 0.2100 & 0.0100 & computational &
\citenum{Yang_2021_10_1016_j_dyepig_2021_109474} \\
Ac-BPCN & 0.3133 & 0.2940 & 0.3797 & 0.3575 & 0.1300 & experimental &
\citenum{Kothavale_2020} \\
AcCN & 0.1610 & 0.1785 & 0.1956 & 0.2129 & 0.0700 & experimental & \citenum{Choi_2019} \\
AcCYP & 0.0725 & 0.0850 & 0.0824 & 0.0981 & 0.2000 & experimental & \citenum{Choi_2020} \\
B-dpa-Cz & 0.4397 & 0.3954 & 0.4512 & 0.4060 & 0.1500 & experimental & \citenum{Park_2022}
\\
B-dpa-SpiroAC & 0.4062 & 0.3635 & 0.4135 & 0.3709 & 0.1400 & experimental &
\citenum{Park_2022} \\
BCZ-DPS-AD & 0.3282 & 0.2926 & 0.3383 & 0.3004 & 0.0967 & experimental &
\citenum{Huo_2023}
\\
BDAPM & 0.3393 & 0.3383 & 0.3593 & 0.3599 & 0.0600 & experimental &
\citenum{Sun_2017_10_1039_c6tc04653c} \\
BDPCC & 0.3333 & 0.3139 & 0.3425 & 0.3241 & 0.1100 & experimental & \citenum{Kumar_2022}
\\
BDQ-tBuCz & 0.1592 & 0.1565 & 0.1872 & 0.1835 & 0.1800 & experimental & \citenum{Yu_2017}
\\
BDQC-2 & 0.1647 & 0.1628 & 0.1856 & 0.1829 & 0.1200 & experimental &
\citenum{Kim_2017_10_1039_c6tc04786f} \\
BDQDMAC & 0.0645 & 0.0647 & 0.0735 & 0.0738 & 0.0750 & experimental & \citenum{Yu_2017} \\
BDQPCz & 0.1517 & 0.1623 & 0.1728 & 0.1880 & 0.4950 & experimental & \citenum{Yu_2017} \\
BFCZPZ1 & 0.2679 & 0.2584 & 0.2793 & 0.2685 & 0.3267 & experimental & \citenum{Cai_2021}
\\
BFCZPZ2 & 0.2685 & 0.2615 & 0.2917 & 0.2840 & 0.3100 & experimental & \citenum{Cai_2021}
\\
BMIM & 0.2801 & 0.2712 & 0.2898 & 0.3011 & 0.1400 & experimental & \citenum{Chen_2022} \\
BP-mDTC & 0.1928 & 0.1825 & 0.1832 & 0.1774 & 0.0700 & experimental &
\citenum{Rajamalli_2017} \\
BPCN-2Cz & 0.3558 & 0.3243 & 0.3590 & 0.3265 & 0.2200 & experimental &
\citenum{Zhang_2020_10_1016_j_orgel_2020_105660} \\
BPCN-Cz2Ph & 0.3042 & 0.2773 & 0.3250 & 0.2965 & 0.3100 & experimental &
\citenum{Zhang_2020_10_1016_j_orgel_2020_105660} \\
BPCP-2CPC & 0.1738 & 0.1724 & 0.2227 & 0.2182 & 0.1660 & experimental &
\citenum{Kim_2019_10_1039_c9cc05391c} \\
BPZPP & 0.1329 & 0.1344 & 0.1529 & 0.1531 & 0.2000 & experimental & \citenum{Chen_2023} \\
BTCZPZ1 & 0.2457 & 0.2342 & 0.2592 & 0.2483 & 0.2400 & experimental & \citenum{Cai_2021}
\\
BTD-PXZ & 0.1863 & 0.1829 & 0.1718 & 0.1684 & 0.1300 & experimental & \citenum{Ma_2023} \\
BTITrz & 0.2219 & 0.2224 & 0.2289 & 0.2290 & 0.1025 & computational &
\citenum{Konidena_2019_10_1039_c9tc04679h} \\
BTPBFCz & 0.2662 & 0.2440 & 0.2687 & 0.2490 & 0.0900 & experimental &
\citenum{Lee_2021_10_1016_j_cej_2020_127293} \\
BTPBFCz-D3 & 0.2264 & 0.2157 & 0.2175 & 0.2055 & 0.1200 & experimental &
\citenum{Lee_2021_10_1016_j_cej_2020_127293} \\
BTPO & 0.4356 & 0.4211 & 0.4061 & 0.3864 & 0.4300 & experimental &
\citenum{Ma_2021_10_1039_d1sc04118e} \\
BTZPP & 0.3066 & 0.3191 & 0.2829 & 0.2942 & 0.3000 & experimental & \citenum{Zhou_2019} \\
BTZ–DMAC & 0.1274 & 0.1340 & 0.1126 & 0.1194 & 0.0200 & computational & \citenum{Ni_2017}
\\
BZC-PXZ & 0.2407 & 0.2284 & 0.2292 & 0.2175 & 0.0200 & experimental & \citenum{Wang_2020}
\\
BuPCzBCO & 0.3702 & 0.3307 & 0.3938 & 0.3517 & 0.1850 & computational &
\citenum{Shin_2023}
\\
BzITz & 0.2513 & 0.2477 & 0.2575 & 0.2534 & 0.1850 & computational &
\citenum{Raikwar_2022}
\\
CCDC & 0.3066 & 0.2759 & 0.1276 & 0.1415 & 0.2220 & computational &
\citenum{Zhao_2017_10_1039_c7qm00195a} \\
CDBP-BP-DMAC & 0.1829 & 0.1967 & 0.2389 & 0.2546 & 0.0393 & experimental &
\citenum{Chen_2022_10_1039_d1qm01625c} \\
CDBP-BP-PXZ & 0.2877 & 0.2868 & 0.2877 & 0.2862 & 0.0605 & experimental &
\citenum{Chen_2022_10_1039_d1qm01625c} \\
CN-QP & 0.1557 & 0.1648 & 0.1216 & 0.1341 & 0.0300 & experimental & \citenum{Huang_2019}
\\
CNPP-TPA & 0.2225 & 0.2259 & 0.2549 & 0.2570 & 0.1450 & experimental &
\citenum{Wang_2022_10_1016_j_cej_2022_135080} \\
CNTPA-CZ & 0.5282 & 0.5166 & 0.5729 & 0.5571 & 0.3500 & experimental & \citenum{Ran_2019}
\\
CNTPA-PXZ & 0.3266 & 0.3257 & 0.3370 & 0.3364 & 0.1300 & experimental & \citenum{Ran_2019}
\\
CZ-DPS-DMAC & 0.3076 & 0.2840 & 0.2851 & 0.2885 & 0.0700 & experimental &
\citenum{Zheng_2022} \\
CZ-TRZ & 0.2454 & 0.2302 & 0.2816 & 0.2641 & 0.1500 & experimental & \citenum{Kumar_2020}
\\
CZ9CZPZ & 0.2859 & 0.2751 & 0.3162 & 0.3040 & 0.3150 & experimental & \citenum{Cai_2021}
\\
Cz-SO & 0.3954 & 0.3540 & 0.3935 & 0.3540 & 0.4100 & experimental &
\citenum{Liu_2021_10_1016_j_tet_2020_131869} \\
CzDCbTrz & 0.1759 & 0.1760 & 0.2349 & 0.2348 & 0.1200 & experimental & \citenum{Kim_2017}
\\
CzPm & 0.2873 & 0.2766 & 0.3039 & 0.2960 & 0.1800 & experimental & \citenum{Leitonas_2022}
\\
CzmPPC & 0.1031 & 0.1273 & 0.0742 & 0.0936 & 0.2800 & experimental &
\citenum{Chen_2022_10_1039_d2tc01435a} \\
CzoB & 0.0786 & 0.0780 & 0.0975 & 0.1027 & 0.0900 & computational & \citenum{Kumar_2019}
\\
DAcIPN & 0.3427 & 0.3327 & 0.3332 & 0.3215 & 0.0350 & experimental &
\citenum{Skuodis_2018}
\\
DBCP & 0.0634 & 0.0645 & 0.0814 & 0.1017 & 0.0813 & experimental &
\citenum{Chen_2022_10_1016_j_cej_2022_134731} \\
DBQ-3PXZ & 0.1057 & 0.1015 & 0.0955 & 0.0920 & 0.0300 & experimental & \citenum{Yu_2018}
\\
DBT-BZ-DMAC & 0.3142 & 0.3200 & 0.3062 & 0.3102 & 0.0800 & experimental &
\citenum{Barman_2022} \\
DC-ACR & 0.0033 & 0.0030 & 0.0366 & 0.0414 & 0.0100 & experimental &
\citenum{Chen_2017_10_1039_c7tc01164d} \\
DC-TC & 0.0052 & 0.0048 & 0.0056 & 0.0054 & 0.1400 & experimental &
\citenum{Chen_2017_10_1039_c7tc01164d} \\
DCPDAPM & 0.0701 & 0.0701 & 0.5171 & 0.4998 & 0.1300 & computational &
\citenum{Zhao_2018_10_1039_c7tc04934j} \\
DCz-DPPZ & 0.1560 & 0.1490 & 0.2024 & 0.1931 & 0.1900 & experimental &
\citenum{Zhang_2022_10_1016_j_dyepig_2022_110634} \\
DCzBNPh-1 & 0.2563 & 0.2378 & 0.3209 & 0.2983 & 0.5000 & experimental & \citenum{Yan_2022}
\\
DCzIPN & 0.3588 & 0.3410 & 0.3593 & 0.3429 & 0.1433 & experimental & \citenum{Kreiza_2021}
\\
DDMA-TXO2 & 0.5414 & 0.5304 & 0.4180 & 0.3968 & 0.2200 & experimental &
\citenum{Ward_2020}
\\
DDPhCz-DPPZ & 0.1229 & 0.1229 & 0.1514 & 0.1470 & 0.1600 & experimental &
\citenum{Zhang_2022_10_1016_j_dyepig_2022_110634} \\
DFDBQPXZ & 0.1692 & 0.1691 & 0.1596 & 0.1589 & 0.0000 & experimental &
\citenum{Yu_2018_10_1039_c7cc09925h} \\
DHID-DBS & 0.4448 & 0.4129 & 0.4697 & 0.4490 & 0.3300 & experimental &
\citenum{Tongsuk_2021}
\\
DHID-DPS & 0.5477 & 0.5208 & 0.4687 & 0.4470 & 0.1800 & experimental &
\citenum{Tongsuk_2021}
\\
DHPZ-2BI & 0.2115 & 0.2147 & 0.2053 & 0.2075 & 0.1900 & experimental & \citenum{Lee_2015}
\\
DMAC-11-DPPZ & 0.0560 & 0.0574 & 0.0714 & 0.0723 & 0.0803 & experimental &
\citenum{Zhou_2020} \\
DMAC-AQ & 0.4361 & 0.4090 & 0.4923 & 0.4663 & 0.0633 & experimental &
\citenum{Huang_2022_10_1016_j_orgel_2022_106657} \\
DMAC-BP & 0.0644 & 0.0658 & 0.5157 & 0.4999 & 0.0700 & experimental &
\citenum{Shakeel_2018}
\\
DMAC-Cz & 0.4201 & 0.3912 & 0.4611 & 0.4309 & 0.0007 & experimental &
\citenum{Zhang_2021_10_1016_j_orgel_2021_106312} \\
DMAC-Cz-TTR & 0.2709 & 0.2527 & 0.2312 & 0.2200 & 0.0660 & experimental &
\citenum{Zhang_2021_10_1016_j_orgel_2021_106312} \\
DMAC-DPS & 0.2222 & 0.2229 & 0.2340 & 0.2343 & 0.0100 & experimental & \citenum{Luo_2016}
\\
DMAC-FO & 0.1247 & 0.1195 & 0.3839 & 0.3734 & 0.1400 & experimental &
\citenum{Yu_2019_10_1016_j_orgel_2019_06_008} \\
DMAC-PCN & 0.1341 & 0.1504 & 0.1133 & 0.1292 & 0.1300 & experimental &
\citenum{Li_2020_10_1016_j_cej_2020_125276} \\
DMAC-TRZ & 0.0485 & 0.0571 & 0.0494 & 0.0510 & 0.0500 & experimental &
\citenum{Braveenth_2023} \\
DMACMNPTO & 0.5377 & 0.5222 & 0.5471 & 0.5302 & 0.2050 & experimental & \citenum{Guo_2020}
\\
DMBFTX & 0.5437 & 0.4736 & 0.5723 & 0.5002 & 0.2800 & experimental &
\citenum{Wang_2018_10_1016_j_cclet_2017_07_025} \\
DMIC-TRZ & 0.5193 & 0.6832 & 0.7550 & 1.3766 & 0.0560 & experimental &
\citenum{Chen_2022_10_1016_j_orgel_2022_106536} \\
DMeCzIPN & 0.2790 & 0.2669 & 0.2781 & 0.2657 & 0.0910 & experimental &
\citenum{Kreiza_2021}
\\
DPACPhTPI & 0.2551 & 0.2503 & 0.2498 & 0.2437 & 0.7200 & experimental &
\citenum{Huang_2017}
\\
DPACpBr & 0.3348 & 0.3322 & 0.4044 & 0.3961 & 0.0020 & experimental & \citenum{Lee_2021}
\\
DPE & 0.2845 & 0.2840 & 0.2987 & 0.2980 & 0.1800 & computational &
\citenum{He_2018_10_1039_c7tc02763j} \\
DPE-DDMAc & 0.1751 & 0.1961 & 0.2781 & 0.3000 & 0.1800 & computational &
\citenum{He_2018_10_1039_c7tc02763j} \\
DPE-DPXZ & 0.2852 & 0.2849 & 0.2996 & 0.2988 & 0.2000 & experimental &
\citenum{He_2018_10_1039_c7tc02763j} \\
DPIBPZ-DPXZ & 0.1101 & 0.1057 & 0.1073 & 0.1034 & 0.0500 & experimental &
\citenum{Wang_2022_10_1016_j_cej_2021_134423} \\
DPS-PXZ & 0.2879 & 0.2873 & 0.3129 & 0.3120 & 0.0500 & experimental & \citenum{Leng_2020}
\\
DPXZ-DPPM & 0.1047 & 0.1014 & 0.1043 & 0.1013 & 0.0470 & experimental &
\citenum{Liang_2020}
\\
DQP & 0.0189 & 0.0396 & 0.0164 & 0.0402 & 0.0555 & experimental & \citenum{Barman_2022} \\
DTPTCzDP-CN & 0.2642 & 0.2402 & 0.2493 & 0.2321 & 0.0300 & experimental &
\citenum{Shi_2022_10_1016_j_dyepig_2022_110430} \\
DV-3CzCN & 0.2477 & 0.2398 & 0.3404 & 0.3364 & 0.2300 & experimental &
\citenum{Sun_2022_10_1016_j_dyepig_2021_109967} \\
DVCz-2CzCN & 0.2868 & 0.2623 & 0.2421 & 0.2332 & 0.1600 & experimental &
\citenum{Sun_2022_10_1016_j_orgel_2022_106442} \\
DiCz-Sac & 0.2404 & 0.2322 & 0.2562 & 0.2469 & 0.2000 & experimental & \citenum{Xiao_2021}
\\
ECDPTT & 0.3447 & 0.3433 & 0.3385 & 0.3331 & 0.0800 & experimental & \citenum{Dong_2018}
\\
F2AcBO & 0.2847 & 0.2918 & 0.2588 & 0.2545 & 0.0750 & experimental &
\citenum{Chen_2022_10_1016_j_cej_2021_133078} \\
FAC & 0.5171 & 0.4920 & 0.5458 & 0.5144 & 0.0900 & computational &
\citenum{Zhu_2022_10_1039_d2tc00344a} \\
FDQPXZ & 0.1276 & 0.1217 & 0.1281 & 0.1225 & 0.0450 & experimental & \citenum{Liu_2022} \\
FSF4A & 0.2862 & 0.2912 & 0.2898 & 0.2923 & 0.0220 & experimental & \citenum{Song_2021} \\
Fene & 0.1297 & 0.1378 & 0.1967 & 0.2089 & 0.0250 & computational &
\citenum{Zhang_2021_10_1016_j_cclet_2020_07_041} \\
Fens & 0.2404 & 0.2561 & 0.3481 & 0.3589 & 0.0200 & computational &
\citenum{Zhang_2021_10_1016_j_cclet_2020_07_041} \\
FlCz & 0.1952 & 0.2170 & 0.2129 & 0.2360 & 0.5100 & experimental & \citenum{Zhen_2022} \\
IABTCN & 0.3551 & 0.3561 & 0.3843 & 0.3786 & 0.3998 & experimental & \citenum{Kim_2019} \\
ICBNTrz4 & 0.2041 & 0.1993 & 0.2097 & 0.2048 & 0.1000 & experimental &
\citenum{Lee_2022_10_1016_j_cej_2021_130988} \\
IDAC-MCO & 0.4472 & 0.3996 & 0.4200 & 0.3755 & 0.0900 & experimental &
\citenum{Chen_2018_10_1016_j_orgel_2018_03_024} \\
IDFL-2DPA & 0.2454 & 0.2221 & 0.2493 & 0.2259 & 0.4500 & experimental &
\citenum{Patil_2020_10_1016_j_cej_2020_125125} \\
IQ-Se & 0.3316 & 0.3254 & 0.4327 & 0.4125 & 1.1800 & computational &
\citenum{Zhang_2022_10_1039_d1cc06380d} \\
IpCm-PhBzAc & 0.5082 & 0.4902 & 0.5207 & 0.5015 & 0.3500 & experimental &
\citenum{Chen_2018_10_1039_c8tc02739k} \\
MBF & 0.2198 & 0.2225 & 0.2260 & 0.2181 & 0.0100 & experimental & \citenum{Zhan_2021} \\
MCZ-P2-DTM & 0.2697 & 0.2587 & 0.2716 & 0.2621 & 0.2300 & experimental & \citenum{Ma_2019}
\\
MPA & 0.2598 & 0.2577 & 0.2572 & 0.2543 & 0.1400 & experimental & \citenum{Zhan_2021} \\
MPPA & 0.1960 & 0.1890 & 0.1952 & 0.1867 & 0.1743 & experimental & \citenum{Sharma_2022}
\\
Me-DMAC & 0.1117 & 0.1283 & 0.1387 & 0.1562 & 0.0100 & experimental &
\citenum{Li_2020_10_1039_c9tc05340a} \\
Me-DOC & 0.4294 & 0.3786 & 0.4302 & 0.3763 & 0.2400 & computational & \citenum{Liu_2020}
\\
Me-MOC & 0.3803 & 0.3386 & 0.3800 & 0.3370 & 0.2500 & computational & \citenum{Liu_2020}
\\
Me2ACBO & 0.3142 & 0.3054 & 0.2883 & 0.2898 & 0.0500 & experimental &
\citenum{Chen_2022_10_1016_j_cej_2021_133078} \\
MeOCzPCN & 0.4683 & 0.4467 & 0.4954 & 0.4673 & 0.0100 & experimental &
\citenum{Tsiko_2021}
\\
NA-TNA & 0.4044 & 0.4138 & 0.4040 & 0.4137 & 0.2400 & experimental & \citenum{Zheng_2015}
\\
NAI–PhBiFA & 0.3697 & 0.3820 & 0.2973 & 0.3055 & 0.1200 & experimental &
\citenum{Chen_2019}
\\
NIDPA-1 & 0.3876 & 0.3588 & 0.4076 & 0.3773 & 0.6000 & experimental & \citenum{Patil_2020}
\\
NPE-AcDPS & 0.4903 & 0.4850 & 0.3559 & 0.3563 & 0.0300 & experimental &
\citenum{Huang_2022_10_1016_j_dyepig_2021_109860} \\
NyDPAc & 0.1657 & 0.1819 & 0.1808 & 0.1987 & 0.1520 & experimental & \citenum{Lee_2020} \\
NyDPO & 0.1494 & 0.1611 & 0.1945 & 0.2037 & 0.0333 & computational &
\citenum{Lee_2020_10_1039_d0tc01644f} \\
OCzBN & 0.3704 & 0.3466 & 0.4579 & 0.4450 & 0.1400 & experimental & \citenum{Ban_2019} \\
OPDPO & 0.1450 & 0.1626 & 0.0857 & 0.0813 & 0.1300 & experimental & \citenum{Zhao_2018} \\
P4CzCN–BCz & 0.6006 & 0.6278 & 0.6999 & 0.7319 & 0.1700 & experimental &
\citenum{Zhou_2022_10_1039_d1tc05010a} \\
PBICT & 0.2132 & 0.1980 & 0.3424 & 0.3194 & 0.0280 & experimental &
\citenum{Kothavale_2021}
\\
PFBP-2b & 0.1515 & 0.1519 & 0.1663 & 0.1583 & 0.0600 & experimental &
\citenum{Hladka_2018}
\\
PIAnTPh & 0.7770 & 0.7391 & 0.8291 & 0.7873 & 1.1400 & experimental &
\citenum{Liu_2021_10_1016_j_cej_2021_131351} \\
PO-TPA & 0.1934 & 0.1883 & 0.3355 & 0.3190 & 0.2500 & experimental & \citenum{Liang_2023}
\\
PPO21 & 0.3775 & 0.3353 & 0.3825 & 0.3406 & 0.0253 & computational &
\citenum{WeiXiaofang_2022} \\
PPZPPI & 0.2203 & 0.2236 & 0.2151 & 0.2175 & 0.1100 & experimental &
\citenum{Huang_2018_10_1039_c7tc05576e} \\
PPZTPI & 0.2232 & 0.2276 & 0.2160 & 0.2192 & 0.1100 & experimental &
\citenum{Huang_2018_10_1039_c7tc05576e} \\
PS-BZ-DMAC & 0.5560 & 0.5307 & 0.5994 & 0.5716 & 0.0045 & experimental &
\citenum{Ma_2021_10_1016_j_dyepig_2021_109649} \\
PSPBP & 0.2941 & 0.2927 & 0.2819 & 0.2814 & 0.1100 & experimental & \citenum{Wu_2021} \\
PSPP & 0.3194 & 0.3148 & 0.3162 & 0.3119 & 0.1100 & experimental & \citenum{Wu_2021} \\
PT-BZ-DMAC & 0.0678 & 0.0852 & 0.0520 & 0.0640 & 0.0084 & experimental &
\citenum{Ma_2021_10_1016_j_dyepig_2021_109649} \\
PTPC & 0.2440 & 0.2387 & 0.2293 & 0.2245 & 0.6355 & experimental & \citenum{Ying_2022} \\
PTZ-2PTO & 0.2932 & 0.2868 & 0.2926 & 0.2929 & 0.2700 & experimental &
\citenum{Xiang_2019}
\\
PTZ-DCPP & 0.2776 & 0.2665 & 0.3429 & 0.3333 & 0.1867 & experimental & \citenum{Wang_2018}
\\
PTZ-NAI & 0.0003 & 0.0003 & 0.1172 & 0.1133 & 0.0867 & experimental & \citenum{Wang_2021}
\\
PX-TRZ & 0.2424 & 0.2355 & 0.2529 & 0.2507 & 0.0130 & experimental & \citenum{Qu_2020} \\
PXZ-10-DPPZ & 0.0924 & 0.0891 & 0.0913 & 0.0886 & 0.0365 & experimental &
\citenum{Zhou_2020}
\\
PXZ-11-DPPZ & 0.1207 & 0.1154 & 0.1309 & 0.1255 & 0.0507 & experimental &
\citenum{Zhou_2020}
\\
PXZ-BOO & 0.2410 & 0.2297 & 0.2536 & 0.2410 & 0.0600 & experimental &
\citenum{Chen_2020_10_1039_d0tc03747h} \\
PXZ-NAI & 0.2654 & 0.2530 & 0.3109 & 0.2946 & 0.0600 & experimental & \citenum{Wang_2021}
\\
PXZ-PPO & 0.5104 & 0.4859 & 0.5382 & 0.5229 & 0.6000 & experimental &
\citenum{Chen_2020_10_1039_d0tc03747h} \\
PXZ-QL & 0.3702 & 0.3444 & 0.3846 & 0.3576 & 0.1000 & computational & \citenum{Shen_2021}
\\
PXZ2PTO & 0.3385 & 0.3358 & 0.3417 & 0.3392 & 0.0233 & experimental &
\citenum{Xiang_2018_10_1039_c8tc03648a} \\
PhQLPXZ & 0.1369 & 0.1437 & 0.1971 & 0.2118 & 0.0900 & experimental & \citenum{Liu_2022}
\\
PrDPhAc & 0.4722 & 0.4514 & 0.5308 & 0.5046 & 0.0750 & computational &
\citenum{Liu_2016_10_1039_c6cc02856j} \\
Pra-2DMAC & 0.5074 & 0.4925 & 0.5244 & 0.5098 & 0.4400 & experimental &
\citenum{Zhang_2021_10_1039_d0tc04315j} \\
PxCN & 0.2607 & 0.2601 & 0.2562 & 0.2557 & 0.0367 & experimental & \citenum{Choi_2019} \\
PxCYP & 0.2325 & 0.2438 & 0.2431 & 0.2422 & 0.2000 & experimental & \citenum{Choi_2020} \\
PxPYM & 0.1691 & 0.1801 & 0.2358 & 0.2454 & 0.0100 & computational &
\citenum{Lee_2020_10_1039_d0ra07865d} \\
Py-TPA & 0.1253 & 0.1357 & 0.1337 & 0.1335 & 0.1400 & computational & \citenum{Tai_2023}
\\
PyDCN–DMAC & 0.2456 & 0.2665 & 0.3176 & 0.3359 & 0.0100 & experimental &
\citenum{Dong_2022}
\\
PyIAnTPh & 0.7218 & 0.6752 & 0.7861 & 0.7350 & 1.1400 & experimental &
\citenum{Liu_2021_10_1016_j_cej_2021_131351} \\
SAF-2NP & 0.0816 & 0.0812 & 0.0892 & 0.0874 & 0.0400 & computational &
\citenum{Zhang_2022_10_1039_d2mh00639a} \\
SBA-2DPS & 0.4557 & 0.4531 & 0.3631 & 0.3818 & 0.0900 & experimental & \citenum{Zeng_2019}
\\
SBDBQ-PXZ & 0.1089 & 0.1046 & 0.1198 & 0.1154 & 0.0700 & experimental & \citenum{Yu_2018}
\\
SBF-PXZ & 0.3045 & 0.3048 & 0.3249 & 0.3241 & 0.2600 & experimental &
\citenum{Jesuraj_2020}
\\
SBTF & 0.5022 & 0.4954 & 0.5297 & 0.5222 & 0.2000 & experimental & \citenum{Cha_2015} \\
SFAC-BP-SFAC & 0.0774 & 0.0786 & 0.0898 & 0.0919 & 0.0360 & experimental &
\citenum{Huang_2022} \\
SFCz & 0.2013 & 0.2131 & 0.2129 & 0.2257 & 0.4800 & experimental & \citenum{Zhen_2022} \\
SFI23pPM & 0.2328 & 0.2146 & 0.2600 & 0.2406 & 0.3300 & experimental &
\citenum{Lv_2020_10_1039_d0tc00039f} \\
SIR-A & 0.0001 & 0.0002 & 0.0000 & 0.0001 & 0.0160 & experimental & \citenum{Pei_2022} \\
SIR-B & 0.0070 & 0.0083 & 0.0015 & 0.0017 & 0.0140 & experimental & \citenum{Pei_2022} \\
SIR-C & 0.0082 & 0.0089 & -0.0005 & -0.0002 & 0.0150 & experimental & \citenum{Pei_2022}
\\
SPAC-FO & 0.1398 & 0.1339 & 0.1384 & 0.1326 & 0.1300 & experimental &
\citenum{Yu_2019_10_1016_j_orgel_2019_06_008} \\
SPFS-PXZ & 0.3291 & 0.3293 & 0.3238 & 0.3229 & 0.0153 & experimental &
\citenum{Fan_2019_10_1039_c9tc02032b} \\
STAC-BP-STAC & 0.0858 & 0.0874 & 0.1261 & 0.1271 & 0.0520 & experimental &
\citenum{Huang_2022} \\
SXAC-BP-SXAC & 0.0835 & 0.0851 & 0.0929 & 0.0948 & 0.0610 & experimental &
\citenum{Huang_2022} \\
SiTCNCz & 0.2497 & 0.2300 & 0.2626 & 0.2416 & 0.1300 & experimental &
\citenum{Yun_2021_10_1039_d1tc03225a} \\
Spiro-CN & 0.2889 & 0.2968 & 0.2696 & 0.2768 & 5.5000 & experimental &
\citenum{Nakagawa_2012} \\
SpiroAC-TRZ & 0.0677 & 0.0679 & 0.0597 & 0.0616 & 0.0720 & experimental &
\citenum{Barman_2022} \\
TATP-BP & 0.1452 & 0.1464 & 0.1292 & 0.1354 & 0.1290 & experimental &
\citenum{Chen_2018_10_1039_c8tc04721a} \\
TBPe & 0.7123 & 0.6662 & 0.7738 & 0.7233 & 0.1200 & computational &
\citenum{Duan_2022_10_1039_d1sc05753g} \\
TCPZ & 0.5340 & 0.5069 & 0.5735 & 0.5469 & 0.1400 & experimental & \citenum{Zhang_2013} \\
TCZ1 & 0.4017 & 0.3602 & 0.4205 & 0.3767 & 0.0160 & experimental & \citenum{Skuodis_2018}
\\
TCz1 & 0.4013 & 0.3599 & 0.4206 & 0.3771 & 0.1000 & experimental &
\citenum{Bezvikonnyi_2019}
\\
TCzTrz & 0.1798 & 0.1783 & 0.2079 & 0.2050 & 0.1400 & experimental & \citenum{Kim_2017} \\
TMDBP & 0.5203 & 0.5192 & 0.5528 & 0.5493 & 0.1000 & experimental & \citenum{Chai_2020} \\
TNPZ & 0.0818 & 0.0930 & 0.0960 & 0.1087 & 0.1700 & computational &
\citenum{Liu_2023_10_1039_d2sc06335b} \\
TPA-2 & 0.0416 & 0.0478 & 0.1131 & 0.1203 & 0.1450 & experimental &
\citenum{Zhao_2021_10_1016_j_mtchem_2021_100625} \\
TPA-DCPP & 0.1746 & 0.1796 & 0.1812 & 0.1855 & 0.1480 & experimental &
\citenum{Barman_2022}
\\
TPADSO2 & 0.4055 & 0.3964 & 0.3925 & 0.3818 & 0.3800 & experimental &
\citenum{Sun_2017_10_1016_j_dyepig_2016_12_037} \\
TPAPm & 0.4500 & 0.4557 & 0.4763 & 0.4820 & 0.3200 & experimental & \citenum{Tsiko_2022}
\\
TPD4PA & 0.3244 & 0.2950 & 0.3337 & 0.3056 & 0.0500 & experimental & \citenum{Naveen_2023}
\\
TPSi-F & 0.6295 & 0.6174 & 0.6494 & 0.6375 & 0.1000 & experimental & \citenum{M_hes_2014}
\\
TPXZ-as-TAZ & 0.0953 & 0.0968 & 0.1079 & 0.1048 & 0.2200 & experimental &
\citenum{Zhang_2018} \\
TPXZ–as-TAZ & 0.0983 & 0.0984 & 0.0271 & 0.0357 & 0.0300 & experimental &
\citenum{Xiang_2016} \\
TRZ-pIC & 0.4296 & 0.3745 & 0.4447 & 0.3877 & 0.2900 & experimental & \citenum{Maeng_2020}
\\
TX-CzPh & 0.5513 & 0.4893 & 0.5517 & 0.4916 & 0.6700 & computational &
\citenum{WangShuai_2021} \\
TXO-P-Si & -0.0053 & -0.0056 & -0.0027 & -0.0027 & 0.5222 & experimental &
\citenum{WeiXiaofang_2022} \\
Tri-PXZ-PCN & 0.1726 & 0.1682 & 0.1583 & 0.1532 & 0.0500 & experimental &
\citenum{Chen_2018_10_1039_c8tc01698d} \\
TrisCz-Trz & 0.1356 & 0.1368 & 0.2203 & 0.2251 & 0.0700 & experimental &
\citenum{Cha_2017}
\\
UGH3 & -0.0012 & 0.0083 & 0.0170 & 0.0280 & 0.2200 & experimental & \citenum{Ward_2020} \\
VBC-PTZ & 0.4336 & 0.4376 & 0.4823 & 0.4872 & 0.5300 & experimental & \citenum{Su_2022} \\
VBN0 & 0.3552 & 0.3499 & 0.4127 & 0.4015 & 0.1100 & experimental & \citenum{Sun_2022} \\
V–TPA & 0.2797 & 0.2756 & 0.2992 & 0.2941 & 0.6210 & experimental & \citenum{Ban_2022} \\
Yad & 0.3868 & 0.3562 & 0.4022 & 0.3704 & 0.0200 & computational &
\citenum{Zhang_2021_10_1016_j_cclet_2020_07_041} \\
czpzpy & 0.4202 & 0.3853 & 0.4351 & 0.3990 & 0.1300 & experimental & \citenum{Chen_2015}
\\
dbG2TAZ & 0.4466 & 0.4517 & 0.4754 & 0.4820 & 0.5300 & experimental &
\citenum{Albrecht_2017}
\\
m-CZ-DPS-DMAC & 0.3204 & 0.2975 & 0.3722 & 0.3450 & 0.2100 & experimental &
\citenum{Dong_2023} \\
m-DTPACO & 0.3172 & 0.3132 & 0.3495 & 0.3436 & 0.2100 & experimental & \citenum{Hu_2017}
\\
m-$\nu$-DABNA & 0.2892 & 0.2722 & 0.2900 & 0.2738 & 0.0700 & experimental &
\citenum{Rayappa_Naveen_2022} \\
m2Cz2Trz & 0.1515 & 0.1484 & 0.1759 & 0.1742 & 0.0900 & experimental & \citenum{Lee_2018}
\\
mCP-BP-DMAC & 0.2872 & 0.2856 & 0.2875 & 0.2859 & 0.0160 & experimental &
\citenum{Wu_2023} \\
mDPBPZ-DPXZ & 0.1066 & 0.1030 & 0.1053 & 0.1025 & 0.0200 & computational &
\citenum{Wang_2022_10_1016_j_cej_2021_134423} \\
mP2MPC & 0.2837 & 0.2741 & 0.2978 & 0.2844 & 0.6500 & experimental & \citenum{Peng_2022}
\\
mPTBC & 0.2343 & 0.2243 & 0.2175 & 0.2081 & 0.0060 & experimental &
\citenum{Wang_2017_10_1016_j_dyepig_2017_04_024} \\
mTPA–PPI & 0.2582 & 0.2558 & 0.2701 & 0.2628 & 0.7600 & experimental & \citenum{Liu_2015}
\\
oPTBC & 0.1336 & 0.1284 & 0.1281 & 0.1232 & 0.0070 & experimental &
\citenum{Wang_2017_10_1016_j_dyepig_2017_04_024} \\
p-DTPACO & 0.4180 & 0.4096 & 0.4450 & 0.4365 & 0.0500 & experimental & \citenum{Hu_2017}
\\
pDPAPQ-PXZ & 0.1777 & 0.1708 & 0.1956 & 0.1877 & 0.0600 & experimental &
\citenum{Zhou_2022_10_1016_j_cej_2022_135775} \\
pDPBPZ-PXZ & 0.0906 & 0.0982 & 0.0988 & 0.1070 & 0.0600 & experimental &
\citenum{Zhou_2022_10_1016_j_cej_2022_135775} \\
pDTCz-2DPyS & 0.2269 & 0.2264 & 0.2872 & 0.2793 & 0.2100 & experimental &
\citenum{Rajamalli_2019} \\
pSFIAc2 & 0.3667 & 0.3178 & 0.3779 & 0.3305 & 0.3100 & experimental & \citenum{Meng_2022}
\\
t-BuPCz & 0.3330 & 0.2985 & 0.3504 & 0.3140 & 0.0800 & experimental &
\citenum{Zhou_2019_10_1039_c9cc06804j} \\
t-DABNA & 0.4476 & 0.3993 & 0.4608 & 0.4111 & 0.1600 & experimental & \citenum{Kim_2021}
\\
t3Cz-SO & 0.3555 & 0.3196 & 0.3436 & 0.3105 & 0.2600 & experimental &
\citenum{Liu_2021_10_1016_j_tet_2020_131869} \\
tCbz-mPYRs & 0.2990 & 0.2843 & 0.3273 & 0.3107 & 0.4150 & computational &
\citenum{Serevi_ius_2020_10_1016_j_orgel_2020_105723} \\
t–BuCz2BP & 0.3370 & 0.3279 & 0.3345 & 0.3371 & 0.2150 & experimental &
\citenum{Huang_2016} \\
$\alpha$-2CbPN & 0.2931 & 0.2856 & 0.3139 & 0.3036 & 0.2600 & experimental &
\citenum{Kim_2017} \\
$\alpha$-2DMAC-DBP & 0.0552 & 0.0562 & 0.0580 & 0.0595 & 0.0250 & experimental &
\citenum{Maeng_2021} \\
$\alpha$-DMAC-DBP & 0.0521 & 0.0525 & 0.0511 & 0.0517 & 0.0925 & experimental &
\citenum{Maeng_2021} \\
$\beta$-DI-DBP & 0.1070 & 0.1192 & 0.1275 & 0.1408 & 0.1600 & experimental &
\citenum{Maeng_2021} \\
$\delta$-2CbPN & 0.2610 & 0.2545 & 0.2760 & 0.2673 & 0.1300 & experimental &
\citenum{Kim_2017} \\
\end{longtable}
}

 
%%%%%%%%%%%%%%%%%%%%%%%%%%%%%%%%%%%%%%%%%%%%%%%%%%%%%%%%%%%%%%%%%%%%%
%% Table S3: Outlier Analysis
%%%%%%%%%%%%%%%%%%%%%%%%%%%%%%%%%%%%%%%%%%%%%%%%%%%%%%%%%%%%%%%%%%%%%
{\small
\begin{longtable}{l l S[table-format=3.4] S[table-format=4.4] S[table-format=4.4] l l}
\caption{Outlier analysis for emission wavelength and singlet-triplet gap predictions. 
Outliers were identified using the Median Absolute Deviation (MAD) method with a threshold 
of $3 \times \text{MAD}$ from the median residual. For emission wavelengths, 
MAD = \qty{43.3}{\nano\meter} (threshold: \qty{130.0}{\nano\meter}); for \deltaest, 
MAD = \qty{0.11}{\electronvolt} (threshold: \qty{0.32}{\electronvolt}). 
All predictions are from sTDA-xTB calculations in toluene solvent.}
\label{tab:outliers}
\\
\toprule
\textbf{Molecule} & \textbf{Property} & {\textbf{Exp./Ref.}} & {\textbf{Predicted}} & {\textbf{Residual}} & \textbf{Architecture} & \textbf{Citation} \\
\midrule
\endfirsthead
\multicolumn{7}{c}{{\tablename\ \thetable{} -- continued from previous page}} \\
\toprule
\textbf{Molecule} & \textbf{Property} & {\textbf{Exp./Ref.}} & {\textbf{Predicted}} & {\textbf{Residual}} & \textbf{Architecture} & \textbf{Citation} \\
\midrule
\endhead
\midrule
\multicolumn{7}{r}{{Continued on next page}} \\
\endfoot
\bottomrule
\endfoot

% Emission Wavelength Outliers (12 molecules)
DQP                                 & $\lambda_{\text{PL}}$ [nm] & 562.0 & 1057.9 &   495.9 & Polycyclic acceptor    & \citenum{Barman_2022} \\
{[2,1-b]IF}                           & $\lambda_{\text{PL}}$ [nm] & 347.0 &  770.2 &   423.2 & Other                  & \citenum{Romain_2016} \\
DPXZ-DPPM                           & $\lambda_{\text{PL}}$ [nm] & 636.0 & 1048.7 &   412.7 & Acridine/Phenoxazine   & \citenum{Liang_2020} \\
thio-ether                          & $\lambda_{\text{PL}}$ [nm] & 366.0 &  770.0 &   404.0 & Other                  & \citenum{Romain_2016} \\
$\alpha$-2DMAC-DBP                  & $\lambda_{\text{PL}}$ [nm] & 637.5 & 1039.5 &   402.0 & Acridine/Phenoxazine   & \citenum{Maeng_2021} \\
2PN                                 & $\lambda_{\text{PL}}$ [nm] & 673.3 & 1059.4 &   386.1 & Polycyclic acceptor    & \citenum{Lin_2022} \\
TFM-QP                              & $\lambda_{\text{PL}}$ [nm] & 566.5 &  929.3 &   362.8 & Polycyclic acceptor    & \citenum{Huang_2019} \\
TPXZ-as-TAZ                         & $\lambda_{\text{PL}}$ [nm] & 529.0 &  889.5 &   360.5 & Acridine/Phenoxazine   & \citenum{Zhang_2018} \\
PXZ-10-DPPZ                         & $\lambda_{\text{PL}}$ [nm] & 614.0 &  958.6 &   344.6 & Acridine/Phenoxazine   & \citenum{Zhou_2020} \\
CN-QP                               & $\lambda_{\text{PL}}$ [nm] & 602.0 &  945.1 &   343.1 & Polycyclic acceptor    & \citenum{Huang_2019} \\
$\alpha$-DMAC-DBP                   & $\lambda_{\text{PL}}$ [nm] & 622.0 &  952.2 &   330.2 & Acridine/Phenoxazine   & \citenum{Maeng_2021} \\
DMAC-11-DPPZ                        & $\lambda_{\text{PL}}$ [nm] & 577.5 &  891.9 &   314.4 & Acridine/Phenoxazine   & \citenum{Zhou_2020} \\

% Singlet-Triplet Gap Outliers (8 molecules, Spiro-CN excluded - discussed separately)
IQ-Se                               & \deltaest [eV] & 1.1800 & 0.4327 & -0.7473 & Heteroatom-containing  & \citenum{Zhang_2022_10_1039_d1cc06380d} \\
DMIC-TRZ                            & \deltaest [eV] & 0.0560 & 0.7550 &  0.6990 & Triazine acceptor      & \citenum{Chen_2022_10_1016_j_orgel_2022_106536} \\
3ai                                 & \deltaest [eV] & 0.1800 & 0.8521 &  0.6721 & Other                  & \citenum{Puntscher_2015} \\
TBPe                                & \deltaest [eV] & 0.1200 & 0.7738 &  0.6538 & Other                  & \citenum{Duan_2022_10_1039_d1sc05753g} \\
PS-BZ-DMAC                          & \deltaest [eV] & 0.0045 & 0.5994 &  0.5949 & Acridine/Phenoxazine   & \citenum{Ma_2021_10_1016_j_dyepig_2021_109649} \\
TPSi-F                              & \deltaest [eV] & 0.1000 & 0.6494 &  0.5494 & Heteroatom-containing  & \citenum{M_hes_2014} \\
P4CzCN–BCz                          & \deltaest [eV] & 0.1700 & 0.6999 &  0.5299 & Carbazole-based        & \citenum{Zhou_2022_10_1039_d1tc05010a} \\
TXO-P-Si                            & \deltaest [eV] & 0.5222 & -0.0027 & -0.5250 & Heteroatom-containing  & \citenum{Wei_2022} \\

\end{longtable}
}
% Notes on outliers:
% - Emission wavelength outliers (12): Large prediction errors > 130 nm, primarily affecting 
%   polycyclic and extended conjugated systems where the vertical approximation breaks down.
% - Singlet-triplet gap outliers (8): Residuals > 0.32 eV, representing molecules with either 
%   very small experimental gaps (< 0.1 eV) or unusual electronic structures.
% - Spiro-CN (ΔE_ST = 5.5 eV) is excluded from this table as it represents a non-TADF molecule
%   and is discussed separately in the main text.


%%%%%%%%%%%%%%%%%%%%%%%%%%%%%%%%%%%%%%%%%%%%%%%%%%%%%%%%%%%%%%%%%%%%%
%% Table S4: TD-DFT Benchmarking
%%%%%%%%%%%%%%%%%%%%%%%%%%%%%%%%%%%%%%%%%%%%%%%%%%%%%%%%%%%%%%%%%%%%%

\begin{table}[H]
\centering
\caption{Benchmarking of xTB-based methods against TD-DFT calculations. 
Vertical excitation energies and singlet-triplet gaps were calculated at the 
$\omega$B97X-D/def2-TZVP level of theory using ORCA 6.1.0, with CPCM solvation in toluene. 
The benchmark set comprises 8 representative TADF emitters spanning different 
donor-acceptor architectures. All structures were optimized at the GFN2-xTB level 
in toluene solvent. Both sTDA-xTB and sTD-DFT-xTB systematically underestimate 
\deltaest\, by approximately 0.77--0.79 eV compared to TD-DFT, but preserve the relative 
ranking of molecules.}
\label{tab:tddft_benchmark}
\footnotesize
\begin{tabular}{l S[table-format=1.3] S[table-format=1.3] S[table-format=1.3] S[table-format=1.3] S[table-format=1.3] S[table-format=+1.3] S[table-format=+1.3]}
\toprule
\textbf{Molecule} & 
{\textbf{TD-DFT}} & 
{\textbf{TD-DFT}} & 
{\textbf{TD-DFT}} & 
{\textbf{sTDA-xTB}} & 
{\textbf{sTD-DFT-xTB}} & 
{\textbf{Deviation}} & 
{\textbf{Deviation}} \\
& 
{$S_1$ [eV]} & 
{$T_1$ [eV]} & 
{$\Delta E_{\text{ST}}$ [eV]} & 
{$\Delta E_{\text{ST}}$ [eV]} & 
{$\Delta E_{\text{ST}}$ [eV]} & 
{sTDA [eV]} & 
{sTD-DFT [eV]} \\
\midrule
2TCz-DPS   & 4.251 & 3.159 & 1.092 & 0.321 & 0.296 & +0.771 & +0.796 \\
4CzIPN     & 3.685 & 2.801 & 0.884 & 0.234 & 0.208 & +0.650 & +0.676 \\
CzS2       & 4.455 & 3.196 & 1.259 & 0.381 & 0.352 & +0.878 & +0.907 \\
DMAC-DPS   & 4.168 & 3.329 & 0.839 & 0.234 & 0.234 & +0.605 & +0.605 \\
DMAC-TRZ   & 4.095 & 3.071 & 1.024 & 0.049 & 0.051 & +0.975 & +0.973 \\
PSPCz      & 4.481 & 3.197 & 1.284 & 0.404 & 0.363 & +0.880 & +0.921 \\
Px2BP      & 3.827 & 2.991 & 0.836 & 0.226 & 0.214 & +0.610 & +0.622 \\
TDBA-DI    & 4.052 & 2.889 & 1.163 & 0.380 & 0.344 & +0.783 & +0.819 \\
\midrule
\textbf{Mean}     & 4.127 & 3.079 & 1.048 & 0.279 & 0.258 & +0.769 & +0.790 \\
\textbf{MAE}      & \multicolumn{3}{c}{---} & \multicolumn{2}{c}{0.769} & \multicolumn{2}{c}{0.790} \\
\textbf{RMSE}     & \multicolumn{3}{c}{---} & \multicolumn{2}{c}{0.780} & \multicolumn{2}{c}{0.801} \\
\bottomrule
\end{tabular}
\end{table}

\textbf{Notes:}
\begin{itemize}
\item \textbf{Computational details:} TD-DFT calculations performed with ORCA 6.1.0 using 
$\omega$B97X-D functional with def2-TZVP basis set and CPCM solvation model (toluene, $\varepsilon=2.374$). 
Vertical excitations calculated on GFN2-xTB/GBSA(toluene) optimized geometries.

\item \textbf{Systematic underestimation:} Both sTDA-xTB and sTD-DFT-xTB consistently 
underestimate \deltaest\ compared to TD-DFT (mean deviation $\sim\qty{0.77}{\electronvolt}$--$\qty{0.79}{\electronvolt}$). This systematic 
bias does not affect relative ranking for high-throughput screening applications.

\item \textbf{Correlation:} Despite the systematic offset, xTB methods correctly identify 
molecules with smaller vs. larger \deltaest\ values (e.g., Px2BP and 4CzIPN have the smallest 
gaps in both TD-DFT and xTB calculations).

\item \textbf{Computational cost:} TD-DFT calculations require $\sim2-6$ hours per molecule 
on 8 CPU cores, while xTB methods complete in $\sim2-5$ minutes, enabling the screening of 
747 molecules presented in this work.

\item \textbf{Recommendation:} Use xTB methods for initial high-throughput screening to 
identify promising candidates, followed by TD-DFT validation for quantitative accuracy.
\end{itemize}



% \subsection{Data Availability}
% 
% All computational data, analysis scripts, and molecular structures are available at
% \href{https://doi.org/10.5281/zenodo.17436069}{10.5281/zenodo.17436069}.



%%%%%%%%%%%%%%%%%%%%%%%%%%%%%%%%%%%%%%%%%%%%%%%%%%%%%%%%%%%%%%%%%%%%%
%% If you are using classical BibTeX rather than biblatex,
%% remove the \printbibliography and uncomment the \bibliograpy one
%%%%%%%%%%%%%%%%%%%%%%%%%%%%%%%%%%%%%%%%%%%%%%%%%%%%%%%%%%%%%%%%%%%%%
%\printbibliography
\bibliography{TADF_Article_References}
\end{document}
